% In this file you should put the actual content of the blueprint.
% It will be used both by the web and the print version.
% It should *not* include the \begin{document}
%
% If you want to split the blueprint content into several files then
% the current file can be a simple sequence of \input. Otherwise It
% can start with a \section or \chapter for instance.

\part{Response-adaptive targeting for multi-treatment experiments}

\textbf{Overview}

This blueprint guides the Lean formalization of the paper
\emph{On Response-Adaptive Targeting Strategies for Multi-Treatment Experiments}
by Redouane Yagouti, Rémy Degenne and Emilie Kaufmann. The paper introduces the
$\alpha$-Rebalancing Targeting Strategies ($\aRTS$), a family of response-adaptive
randomization procedures for clinical trials with $K \geq 2$ arms, and proves
that every member of the family enjoys the same asymptotic guarantees: strong
consistency, asymptotic normality of the estimators and allocation proportions,
and asymptotic efficiency. A forced-exploration variant ($\aRTSFE$) extends these
guarantees, in particular to sparse target allocations.

The mathematical development is organized so that the two main theorems are
proved once, from an abstract set of \emph{generic sufficient conditions}
(Chapter~\ref{chap:generic}), and each concrete design family is then handled by
verifying those conditions. The blueprint is deliberately decomposed into many
small lemmas to make the formalization tractable.

\textbf{Reading guide}
\begin{itemize}
  \item Chapters~\ref{chap:model}--\ref{chap:auxiliary} set up the model,
        estimators, conditions, the $\aRTS$ family and the auxiliary processes.
  \item Chapter~\ref{chap:technical} collects self-contained technical lemmas
        (good starting points for formalization).
  \item Chapter~\ref{chap:generic} states the generic conditions and verifies
        them for $\aRTS$.
  \item Chapters~\ref{chap:consistency}--\ref{chap:normality} prove the two main
        theorems from the generic conditions.
  \item Chapter~\ref{chap:forced} treats forced exploration and sparse targets.
  \item Chapter~\ref{chap:testing} gives an application to hypothesis testing.
\end{itemize}

Part~\ref{part:prereq} is separate: it details the probabilistic results this
formalization needs that are \emph{not yet in Mathlib}, decomposed into small
lemmas with proof sketches, so that they can be contributed to Mathlib
independently.

\textbf{Status} The formalization has not started.

\textbf{Blueprint author} Rémy Degenne. Anyone is welcome to contribute!

\chapter{Model, notation and assumptions}
\label{chap:model}

This chapter sets up the probabilistic model of a sequential multi-treatment
experiment, the estimators used, and the family of allocation rules
($\aRTS$) whose asymptotic behaviour is analysed in the rest of the blueprint.
It follows Section~2 and Section~3 of the paper.

\textbf{Lean remark.}
Throughout, everything lives on a fixed probability space $(\Omega, \filt, \Prob)$.
We work with $K \geq 2$ treatments (``arms''), indexed by $[K] := \{1, \dots, K\}$.
Random objects are indexed by a time $m \in \N$ (the patient number).
A convenient Lean encoding is to fix once and for all:
\begin{itemize}
  \item a natural number $K$ with $2 \le K$,
  \item a probability space $(\Omega, \filt, \Prob)$,
  \item the assignment process $X : \N \to \Omega \to (\text{Fin } K \to \{0,1\})$,
  \item the response process $\xi : \N \to \text{Fin } K \to \Omega \to \R$.
\end{itemize}

\section{The sequential experiment}
\label{sec:model_experiment}

\begin{definition}[Assignment vector]\label{def:assignment}
\lean{Learning.actionIndicator}\leanok
For each patient $m \ge 1$, the \emph{assignment vector}
$X_m = (X_{m,1}, \dots, X_{m,K})$ takes values in $\{0,1\}^K$ and satisfies
$\sum_{k=1}^K X_{m,k} = 1$. The event $\{X_{m,k} = 1\}$ means that patient $m$ is
assigned to arm $k$.
\end{definition}

\begin{definition}[Responses and parameters]\label{def:responses}
For each arm $k \in [K]$, $\{\xi_{m,k} : m \ge 1\}$ is a sequence of i.i.d.\
real random variables, the potential responses of patients on treatment $k$.
The parameter of interest for arm $k$ is the mean
$\theta_k := \E[\xi_{1,k}]$, and $\Params = (\theta_1, \dots, \theta_K) \in \R^K$.
We write $V_k := \Var[\xi_{1,k}]$ for the variance of the response of arm $k$.
\end{definition}

\begin{definition}[Allocation counts]\label{def:counts}
\lean{AlphaRAR.count}\leanok
\uses{def:assignment}
The number of patients assigned to arm $k$ up to time $n$ is
\[
  N_{n,k} := \sum_{j=1}^n X_{j,k}, \qquad N_n := (N_{n,1}, \dots, N_{n,K}).
\]
The vector $N_n / n$ is the \emph{observed allocation proportion} at time $n$.
\end{definition}

\begin{lemma}[Counts sum to time]\label{lem:counts_sum}
\lean{AlphaRAR.counts_sum}\leanok
\uses{def:counts}
For all $n \ge 1$, $\sum_{k=1}^K N_{n,k} = n$, and hence
$\sum_{k=1}^K \frac{N_{n,k}}{n} = 1$.
\end{lemma}

\begin{proof}\leanok
\uses{def:assignment}
Immediate from $\sum_{k=1}^K X_{j,k} = 1$ and summing over $j \le n$.
\end{proof}

\begin{definition}[Filtration]\label{def:filtration}
\lean{Learning.IsAlgEnvSeq.filtration}\leanok
\uses{def:assignment, def:responses}
$\filt_m := \sigma\big((X_i, \xi_i)_{1 \le i \le m}\big)$ is the filtration
generated by the assignments and observed outcomes of the first $m$ patients,
with $\filt_0$ the trivial $\sigma$-algebra.
\end{definition}

\begin{definition}[RAR procedure]\label{def:rar}
\uses{def:filtration}
A \emph{response-adaptive randomization (RAR) procedure} is a sequence
$(p_m)_{m \ge 1}$ of $\filt_m$-measurable random vectors taking values in the
probability simplex
\[
  \simplex_K := \Big\{ p \in [0,1]^K : \sum_{k=1}^K p_k = 1 \Big\},
\]
such that the assignment $X_{m+1}$ is drawn conditionally on $\filt_m$ according
to $p_m$:
\[
  \Prob(X_{m+1} = k \mid \filt_m) = p_{m,k}, \qquad k \in [K].
\]
\end{definition}

\textbf{Lean remark.} The conditional-law requirement is best captured through
the conditional expectation $\E[X_{m+1,k} \mid \filt_m] = p_{m,k}$
(see Definition~\ref{def:cond_prob}), which is the only consequence used in the proofs.

\begin{definition}[Target function and target allocation]\label{def:target}
A \emph{target function} is a map $\Target : \hilb \to \simplex_K$ defined on an
open domain $\hilb \subseteq \R^K$, written
$\Target(z) = (\Target_1(z), \dots, \Target_K(z))$ with
$\sum_{k=1}^K \Target_k(z) = 1$. The \emph{target allocation} is
$v := \Target(\Params)$. A RAR procedure is an \emph{adaptive targeting
mechanism} if $N_n/n \to v$.
\end{definition}

\section{Sequential estimation}
\label{sec:model_estimation}

\begin{definition}[Sequential estimator]\label{def:estimator}
\lean{AlphaRAR.estimator}\leanok
\uses{def:counts, def:responses}
Given initial values $\theta_{0,k} \in \R$, the sequential estimator of
$\theta_k$ at stage $m$ is
\[
  \hat{\theta}_{m,k} := \frac{\sum_{j=1}^m X_{j,k}\, \xi_{j,k} + \theta_{0,k}}{N_{m,k} + 1}.
\]
We write $\estParams_m := (\hat{\theta}_{m,1}, \dots, \hat{\theta}_{m,K})$.
\end{definition}

\begin{definition}[Plug-in target]\label{def:plugin_target}
\lean{AlphaRAR.aRTSTarget}\leanok
\uses{def:estimator, def:target}
The plug-in estimator of the target allocation is
$\estTarget_m := \Target(\estParams_m)$, with components $\hat{\rho}_{m,k}$.
It is well defined once $\estParams_m \in \hilb$, which is guaranteed by
Condition~\ref{def:condB}.
\end{definition}

\begin{definition}[Bahadur representation]\label{def:bahadur}
\uses{def:estimator}
An estimator $\hat{\theta}_{m,k}$ of $\theta_k$ admits a \emph{Bahadur
representation} if
\[
  \hat{\theta}_{m,k} = \frac{1}{N_{m,k}}\sum_{j=1}^m X_{j,k}\, \xi_{j,k}
    + o\!\big(N_{m,k}^{-1/2}\big), \qquad m \to \infty.
\]
\end{definition}

\begin{lemma}[The sequential estimator is Bahadur]\label{lem:estimator_bahadur}
\lean{AlphaRAR.estimator_sub_eq}\leanok
\uses{def:estimator, def:bahadur, def:Q}
On the event $\{N_{m,k} \to \infty\}$, the estimator of
Definition~\ref{def:estimator} admits the Bahadur representation of
Definition~\ref{def:bahadur}. The formalized statement \texttt{estimator\_sub\_eq} is the sharper
\emph{exact} error identity
\[
  \hat{\theta}_{m,k} - \theta_k = \frac{Q_{m,k} + (\theta_{0,k} - \theta_k)}{N_{m,k} + 1},
\]
which holds pathwise for every $m$ (with $N_{m,k}+1 \ne 0$, automatic for $\{0,1\}$-indicators) and
from which the $o(N^{-1/2})$ Bahadur form is immediate.
\end{lemma}

\begin{proof}\leanok
The numerator of $\hat{\theta}_{m,k} - \theta_k = \frac{\sum_j X_{j,k}\xi_{j,k} + \theta_{0,k}}{N_{m,k}
+ 1} - \theta_k$ is $\sum_j X_{j,k}\xi_{j,k} + \theta_{0,k} - \theta_k(N_{m,k}+1) = (\sum_j
X_{j,k}\xi_{j,k} - \theta_k N_{m,k}) + (\theta_{0,k} - \theta_k) = Q_{m,k} + (\theta_{0,k} -
\theta_k)$, giving the exact identity. On $\{N_{m,k}\to\infty\}$ the offset
$\frac{\theta_{0,k}-\theta_k}{N_{m,k}+1} = O(N_{m,k}^{-1}) = o(N_{m,k}^{-1/2})$, which is the Bahadur
form.
\end{proof}

\textbf{Lean remark.} The Bahadur property (Definition~\ref{def:bahadur}) is the
only property of the estimator used in the asymptotic analysis. It is worth
stating the main theorems for an abstract estimator satisfying
Definition~\ref{def:bahadur}, then deriving the concrete case from
Lemma~\ref{lem:estimator_bahadur}.

\section{The conditions}
\label{sec:model_conditions}

\begin{definition}[Condition A: second moments]\label{def:condA}
\uses{def:responses}
Condition~\textbf{A} holds if for every arm $k \in [K]$ the response satisfies
$\E|\xi_{1,k}|^2 < \infty$ (equivalently $V_k < \infty$).
\end{definition}

\begin{definition}[Attainable parameter values]\label{def:attainable}
\lean{AlphaRAR.attainableSet}\leanok
\uses{def:estimator}
For each arm $k$, let $V_k^{\mathrm{att}}$ be the set of all values attainable by
the estimator $\hat{\theta}_{n,k}$, over all $n$ and all realizations of the
responses $\{\xi_{j,k}\}$, and set $I_k := \{\theta_k\} \cup V_k^{\mathrm{att}}$.
(The formalized \texttt{AlphaRAR.attainableSet} is the \emph{closure} of the range of
$\hat{\theta}_{\cdot,k}$ over all times and outcomes; being closed it already contains every a.s.\
limit $z_k$ — in particular $\theta_k$ on the infinitely-sampled event — so the explicit union with
$\{\theta_k\}$ is unnecessary.)
\end{definition}

\begin{definition}[Condition B: regular non-sparse target]\label{def:condB}
\uses{def:target, def:attainable}
Condition~\textbf{B} holds if the domain $\hilb$ is open and contains
$I_1 \times \dots \times I_K$, and moreover
\begin{itemize}
  \item $\Target$ is twice differentiable on $\hilb$;
  \item for every $z \in I_1 \times \dots \times I_K$, $\Target(z) \in (0,1)^K$.
\end{itemize}
In particular the target allocation $v = \Target(\Params)$ is non-sparse:
$v_k > 0$ for all $k$.
\end{definition}

\begin{definition}[Sparse target allocation]\label{def:sparse}
\uses{def:target}
A target function $\Target$ is \emph{sparse at $\Params$} if some component of
$v = \Target(\Params)$ is zero.
\end{definition}

\section{Stochastic orders}
\label{sec:model_orders}

The rates in the paper are stated with the probabilistic Landau notation
$o_p$ and $O_p$. We record their definitions and the two facts about them that
are used repeatedly.

\begin{definition}[$o_p$ and $O_p$]\label{def:op_Op}
\lean{AlphaRAR.IsLittleOpOne, AlphaRAR.IsBigOpOne}\leanok
Let $(a_n)_{n\ge1}$ be positive reals and $(X_n)_{n\ge1}$ random variables.
\begin{itemize}
  \item $X_n = o_p(a_n)$ if for all $\varepsilon > 0$,
        $\Prob\!\big(|X_n|/a_n > \varepsilon\big) \to 0$.
  \item $X_n = O_p(a_n)$ if for all $\varepsilon > 0$ there is $M > 0$ with
        $\limsup_n \Prob\!\big(|X_n|/a_n > M\big) \le \varepsilon$
        (boundedness in probability of $X_n / a_n$).
\end{itemize}
\end{definition}

\textbf{Lean remark.} These are exactly the notions
\texttt{Filter.Tendsto} to $0$ in probability and
tightness/boundedness-in-probability. Mathlib already has
\texttt{MeasureTheory} machinery for convergence in probability; a small
API around $o_p/O_p$ (sum, product, scalar multiple, comparison) should be
built first as it is used everywhere.

\begin{lemma}[$o_p$ and $O_p$ under scaling]\label{lem:op_smul}
\uses{def:op_Op}
If $X_n = o_p(a_n)$ and $b_n / a_n \to c \in \R$, then $X_n = o_p(b_n)$
whenever $c \neq 0$; similarly for $O_p$. Moreover $o_p(a_n) + o_p(a_n) = o_p(a_n)$,
$O_p(a_n) + O_p(a_n) = O_p(a_n)$, and $o_p(a_n) + O_p(a_n) = O_p(a_n)$.
\end{lemma}

\begin{proof}
Direct from the definitions and a union bound.
\end{proof}

\chapter{The \texorpdfstring{$\aRTS$}{aRTS} family of designs}
\label{chap:design}

This chapter defines the $\alpha$-Rebalancing Targeting Strategies ($\aRTS$),
the central object of the paper, and records the example designs that belong to
the family. It corresponds to Section~3.1 of the paper.

\section{Definition of the family}
\label{sec:design_def}

\begin{definition}[$\aRTS$ family]\label{def:aRTS}
\uses{def:rar, def:plugin_target, def:counts}
Let $K \ge 2$, $\alpha \in [0,1)$ and $m_0 \in \N$ (a burn-in parameter).
A RAR procedure is an \emph{$\alpha$-Rebalancing Targeting Strategy} ($\aRTS$)
if its allocation probabilities $p_{m+1} \in \simplex_K$ satisfy, for all
$m \ge K m_0$, the \emph{throttling condition}
\begin{equation}\label{eq:throttle}
  \forall k \in [K], \quad
  \frac{N_{m,k}}{m} > \hat{\rho}_{m,k}
  \ \Longrightarrow\
  p_{m+1,k} \le \alpha\, \hat{\rho}_{m,k}.
\end{equation}
In words: whenever arm $k$ is currently over-sampled relative to its estimated
target, the probability of selecting it next is throttled down to at most
$\alpha \hat{\rho}_{m,k}$.
\end{definition}

\begin{remark}\label{rem:burnin}
The burn-in parameter $m_0$ is the number of initial samples of each arm that
may be taken before adaptive randomization starts. It plays no role in the
proofs; one may take $m_0 = 0$ when the regularization $\theta_0$ of the
estimator is chosen appropriately.
\end{remark}

\textbf{Lean remark.} The throttling condition \eqref{eq:throttle} is a
\emph{property} of a RAR procedure, not a construction. The cleanest Lean design
is to define a predicate \texttt{IsAlphaRTS} on RAR procedures and prove all
asymptotic theorems from that predicate, then verify it for each concrete
example below. This is realized by \texttt{AlphaRAR.IsARTS} (the throttle predicate) and, for
constructing concrete designs, the shared plumbing
\texttt{AlphaRAR.aRTSAlgorithmOfProb}: it turns any measurable, history-dependent
\emph{probability-vector} function $p$ into an \texttt{Algorithm} whose policy is the finite-support
kernel $\sum_a p_a \delta_a$ (\texttt{AlphaRAR.probVecKernel}); the companion
\texttt{AlphaRAR.aRTSAlgorithmOfProb\_isARTS} reduces verifying \texttt{IsARTS} for a design to the
pure arithmetic of its throttle inequality (\texttt{AlphaRAR.probVecKernel\_apply\_singleton}
identifies the policy mass on arm $k$ with $p_k$). Each concrete example below is then a
probability-vector function whose simplex and throttle properties are elementary.

\section{Examples}
\label{sec:design_examples}

Each of the following designs is shown to belong to the $\aRTS$ family, i.e.\ to
satisfy \eqref{eq:throttle}. These lemmas are independent of one another and are
good small formalization targets.

\begin{definition}[Distance-based design]\label{def:distance_design}
\lean{AlphaRAR.distanceProb, AlphaRAR.distanceAlgorithm}\leanok
\uses{def:plugin_target, def:counts}
Given $\alpha \in [0,1)$, the \emph{distance-based} design sets, for
$m \ge K m_0$ and $k \in [K]$,
\[
  p_{m+1,k} = \alpha\, \hat{\rho}_{m,k}
  + (1-\alpha)\frac{\delta_{m,k}}{\sum_{i=1}^K \delta_{m,i}},
  \qquad
  \delta_{m,k} := \max\!\Big(0,\ \hat{\rho}_{m,k} - \frac{N_{m,k}}{m}\Big),
\]
with the convention that if $\delta_{m,j} = 0$ for all $j$, then
$p_{m+1,j} = \hat{\rho}_{m,j}$ for all $j$.
\end{definition}

\begin{lemma}[Distance-based is a valid probability vector]\label{lem:distance_simplex}
\lean{AlphaRAR.distanceProb_nonneg, AlphaRAR.distanceProb_sum}\leanok
\uses{def:distance_design}
For all $m \ge K m_0$, the vector $p_{m+1}$ of Definition~\ref{def:distance_design}
lies in $\simplex_K$.
\end{lemma}

\begin{proof}\leanok
Each coordinate is nonnegative, and
$\sum_k p_{m+1,k} = \alpha \sum_k \hat{\rho}_{m,k}
 + (1-\alpha)\frac{\sum_k \delta_{m,k}}{\sum_i \delta_{m,i}} = \alpha + (1-\alpha) = 1$
when some $\delta_{m,i}>0$; in the degenerate case $p_{m+1}=\estTarget_m \in \simplex_K$.
\end{proof}

\begin{lemma}[Distance-based is $\aRTS$]\label{lem:distance_isARTS}
\lean{AlphaRAR.distance_isARTS}\leanok
\uses{def:distance_design, def:aRTS}
The distance-based design satisfies the throttling condition \eqref{eq:throttle},
hence belongs to the $\aRTS$ family.
\end{lemma}

\begin{proof}\leanok
\uses{lem:distance_simplex}
Membership in $\simplex_K$ is Lemma~\ref{lem:distance_simplex}. For the throttle:
if $N_{m,k}/m > \hat{\rho}_{m,k}$ then $\delta_{m,k} = 0$, so
$p_{m+1,k} = \alpha\hat{\rho}_{m,k} \le \alpha \hat{\rho}_{m,k}$.

\emph{Formalization.} The design is \texttt{AlphaRAR.distanceProb} (a probability-vector function of
the history) and the algorithm \texttt{AlphaRAR.distanceAlgorithm}; \texttt{AlphaRAR.distanceProb\_%
throttle} is the throttle inequality. The only subtlety is that the mixing normaliser $\sum_i
\delta_{m,i}$ is positive whenever some arm is over-sampled: the signed deficits
$N_{m,i}/m - \hat\rho_{m,i}$ sum to $0$ (\texttt{sum\_histProp}, \texttt{sum\_histTarget}), so an
over-sampled arm forces an under-sampled one, giving $\delta_{m,j}>0$ for some $j$.
\end{proof}

\begin{definition}[ERADE 2025 design]\label{def:erade2025}
\lean{AlphaRAR.eradeProb, AlphaRAR.eradeAlgorithm}\leanok
\uses{def:plugin_target, def:counts}
Given $\alpha \in [0,1)$, the \emph{ERADE 2025} design of
\cite{alkhnefr2025efficient} sets, for $m \ge K m_0$ and $k \in [K]$,
\[
  p_{m+1,k} =
  \begin{cases}
    \alpha \hat{\rho}_{m,k}, & N_{m,k}/m > \hat{\rho}_{m,k},\\
    \hat{\rho}_{m,k}, & N_{m,k}/m = \hat{\rho}_{m,k},\\
    \hat{\rho}_{m,k} + (1-\alpha)\dfrac{\sum_{j\in S}\hat{\rho}_{m,j}}{|T|},
      & N_{m,k}/m < \hat{\rho}_{m,k},
  \end{cases}
\]
where $S = \{ j : N_{m,j}/m > \hat{\rho}_{m,j}\}$ is the set of over-sampled arms
and $T = \{ j : N_{m,j}/m < \hat{\rho}_{m,j}\}$ the set of under-sampled arms.
\end{definition}

\begin{lemma}[ERADE 2025 is a valid probability vector]\label{lem:erade2025_simplex}
\lean{AlphaRAR.eradeProb_nonneg, AlphaRAR.eradeProb_sum}\leanok
\uses{def:erade2025}
For all $m \ge K m_0$, the vector $p_{m+1}$ of Definition~\ref{def:erade2025}
lies in $\simplex_K$.
\end{lemma}

\begin{proof}\leanok
Nonnegativity is clear. For the sum, the total mass removed from over-sampled
arms is $\sum_{j\in S}(1-\alpha)\hat{\rho}_{m,j}$, which is exactly the mass
added on the $|T|$ under-sampled arms, and arms with equality keep their target;
using $\sum_k \hat{\rho}_{m,k} = 1$ gives $\sum_k p_{m+1,k} = 1$.
\end{proof}

\begin{lemma}[ERADE 2025 is $\aRTS$]\label{lem:erade2025_isARTS}
\lean{AlphaRAR.erade_isARTS}\leanok
\uses{def:erade2025, def:aRTS}
The ERADE 2025 design satisfies \eqref{eq:throttle}, hence belongs to the
$\aRTS$ family.
\end{lemma}

\begin{proof}\leanok
\uses{lem:erade2025_simplex}
Membership in $\simplex_K$ is Lemma~\ref{lem:erade2025_simplex}. For the throttle:
if $N_{m,k}/m > \hat{\rho}_{m,k}$ then $p_{m+1,k} = \alpha\hat{\rho}_{m,k}$.

\emph{Formalization.} \texttt{AlphaRAR.eradeProb}/\texttt{eradeAlgorithm} and
\texttt{AlphaRAR.eradeProb\_throttle}. The simplex identity uses that the mass freed from the
over-sampled arms, $\sum_{j\in S}(1-\alpha)\hat\rho_{m,j}$, is redistributed exactly over the $|T|$
under-sampled arms; the degenerate case $|T|=0$ forces $S=\varnothing$ by mass conservation
(\texttt{sum\_histProp}, \texttt{sum\_histTarget}).
\end{proof}

\begin{definition}[Interpolated D-Tracking]\label{def:d_tracking}
\lean{AlphaRAR.dTrackingProb, AlphaRAR.dTrackingAlgorithm}\leanok
\uses{def:plugin_target, def:counts}
For $\alpha \in [0,1)$, the \emph{Interpolated D-Tracking} rule sets, for
$m \ge K m_0$ and $k \in [K]$,
\[
  p_{m+1,k} = \alpha\, \hat{\rho}_{m,k} + (1-\alpha)\,\indic_{\{k = k^\star\}},
  \qquad
  k^\star \in \argmax_{a \in [K]} \big\{ m\hat{\rho}_{m,a} - N_{m,a} \big\}.
\]
The quantity $m\hat{\rho}_{m,a} - N_{m,a}$ is the deficit of arm $a$ relative to
its target, so $k^\star$ is the most under-sampled arm.
\end{definition}

\begin{lemma}[Interpolated D-Tracking is a valid probability vector]\label{lem:d_tracking_simplex}
\lean{AlphaRAR.dTrackingProb_nonneg, AlphaRAR.dTrackingProb_sum}\leanok
\uses{def:d_tracking}
For all $m \ge K m_0$, $p_{m+1} \in \simplex_K$.
\end{lemma}

\begin{proof}\leanok
$\sum_k p_{m+1,k} = \alpha \sum_k \hat{\rho}_{m,k} + (1-\alpha) = 1$, and each
coordinate is nonnegative.
\end{proof}

\begin{lemma}[Interpolated D-Tracking is $\aRTS$]\label{lem:d_tracking_isARTS}
\lean{AlphaRAR.dTracking_isARTS}\leanok
\uses{def:d_tracking, def:aRTS}
For every $\alpha \in [0,1)$, the Interpolated D-Tracking rule satisfies
\eqref{eq:throttle}, hence belongs to the $\aRTS$ family.
\end{lemma}

\begin{proof}\leanok
\uses{lem:counts_sum, lem:d_tracking_simplex}
Membership in $\simplex_K$ is Lemma~\ref{lem:d_tracking_simplex}.
If $N_{m,k}/m > \hat{\rho}_{m,k}$ then $m\hat{\rho}_{m,k} - N_{m,k} < 0$. Since
$\sum_a (m\hat{\rho}_{m,a} - N_{m,a}) = 0$ by Lemma~\ref{lem:counts_sum}, the
maximizing deficit is $\ge 0$, so $k \ne k^\star$ and
$p_{m+1,k} = \alpha\hat{\rho}_{m,k}$.

\emph{Formalization.} The favoured arm $k^\star=\arg\max_a(m\hat\rho_{m,a}-N_{m,a})$ is
\texttt{argmax} of \texttt{AlphaRAR.dtDeficit} (LeanMachineLearning's measurable finite-tuple
\texttt{argmax}, so the policy kernel is measurable without needing an order on the arms); the
throttle is \texttt{AlphaRAR.dTrackingProb\_throttle}.
\end{proof}

\begin{remark}\label{rem:d_tracking_zero}
When $\alpha = 0$, the Interpolated D-Tracking rule becomes
$p_{m+1,k} = \indic_{\{k = k^\star\}}$, which is the D-Tracking rule of
\cite{garivier2016optimal} (without its forced-exploration component, treated in
Chapter~\ref{chap:forced}).
\end{remark}

\chapter{Auxiliary processes and martingales}
\label{chap:auxiliary}

The analysis of the $\aRTS$ family rests on a handful of auxiliary stochastic
processes built from the assignments and responses. This chapter collects their
definitions and elementary (mostly martingale) properties. It corresponds to the
setup of the proof sketch (Section~3.2) and of Appendix~A.

\section{Conditional probabilities and the assignment martingale}
\label{sec:aux_assignment}

\begin{definition}[Conditional selection probability]\label{def:cond_prob}
\uses{def:assignment, def:filtration}
For $k \in [K]$ and $m \ge 1$,
\[
  p_{m,k} := \E[X_{m,k} \mid \filt_{m-1}].
\]
For a RAR procedure these coincide with the allocation probabilities: with the
indexing convention of Definition~\ref{def:rar}, $p_{m,k}$ is the $k$-th
component of the vector used to draw $X_m$.
\end{definition}

\begin{definition}[Assignment martingale]\label{def:M}
\lean{AlphaRAR.assignMG}\leanok
\uses{def:cond_prob}
For $k \in [K]$ set $\Delta M_{m,k} := X_{m,k} - p_{m,k}$ and
\[
  M_{n,k} := \sum_{m=1}^n (X_{m,k} - p_{m,k}), \qquad M_n := (M_{n,1},\dots,M_{n,K}).
\]
\end{definition}

\begin{lemma}[$M$ is a martingale]\label{lem:M_martingale}
\lean{AlphaRAR.martingale_assignMart}\leanok
\uses{def:M, def:filtration}
For each $k \in [K]$, $(M_{n,k})_{n\ge0}$ is an $(\filt_n)$-martingale with
$M_{0,k}=0$ and bounded increments $|\Delta M_{m,k}| \le 1$.
\end{lemma}

\begin{proof}\leanok
\uses{def:cond_prob}
$\E[\Delta M_{m,k}\mid\filt_{m-1}] = \E[X_{m,k}\mid\filt_{m-1}] - p_{m,k} = 0$
by Definition~\ref{def:cond_prob}, and $X_{m,k},p_{m,k}\in[0,1]$.
\end{proof}

\begin{lemma}[The assignment martingale is $O_p(\sqrt n)$]\label{lem:M_Op}
\lean{AlphaRAR.isBigOpOne_assignMart_div_sqrt}\leanok
\uses{def:M, lem:M_martingale}
For each $k \in [K]$, $M_{n,k} = O_p(\sqrt n)$.
\end{lemma}

\begin{proof}\leanok
\uses{lem:mart_bdd_incr_Op}
$M_{\cdot,k}$ is a martingale with $M_{0,k}=0$ and increments bounded by $1$
(Lemma~\ref{lem:M_martingale}), so Lemma~\ref{lem:mart_bdd_incr_Op} applies with
$c = 1$.
\end{proof}

\begin{lemma}[Count decomposition]\label{lem:count_decomp}
\lean{AlphaRAR.count_eq}\leanok
\uses{def:M, def:counts}
For all $n\ge1$ and $k\in[K]$,
$\displaystyle N_{n,k} = \sum_{m=1}^n p_{m,k} + M_{n,k}$.
\end{lemma}

\begin{proof}\leanok
Sum the identity $X_{m,k} = p_{m,k} + \Delta M_{m,k}$ over $m \le n$.
\end{proof}

\section{The response martingale}
\label{sec:aux_response}

\begin{definition}[Centered response martingale]\label{def:Q}
\lean{AlphaRAR.respMG, AlphaRAR.respMart}\leanok
\uses{def:assignment, def:responses}
For $k \in [K]$ set $Y_{m,k} := \xi_{m,k} - \theta_k$ and
\[
  Q_{n,k} := \sum_{m=1}^n X_{m,k}\,(\xi_{m,k} - \theta_k)
           = \sum_{m=1}^n X_{m,k}\,Y_{m,k},
  \qquad Q_n := (Q_{n,1}, \dots, Q_{n,K}).
\]
\end{definition}

\begin{lemma}[$Q$ is a martingale]\label{lem:Q_martingale}
\lean{AlphaRAR.martingale_respMart}\leanok
\uses{def:Q, def:filtration}
For each $k \in [K]$, if the responses are integrable then $(Q_{n,k})_{n\ge0}$ is an
$(\filt_n)$-martingale.
\end{lemma}

\begin{proof}\leanok
Write $\Delta Q_{n,k} = X_{n,k}(\xi_{n,k} - \theta_k)$, where $X_{n,k}$ is the
indicator that the assignment $A_n$ equals arm $k$. The subtle point is the
\emph{filtration convention}: $X_{n,k}$ is measurable with respect to the
action-augmented $\sigma$-algebra $\mathcal{G}_n := \filt_{n-1}\vee\sigma(A_n)$
(\texttt{filtrationAction}) but \emph{not} with respect to $\filt_{n-1}$ --- indeed,
if it were, the assignment martingale $M_{n,k}$ of Definition~\ref{def:M} would be
identically zero. Conditioning on $\mathcal{G}_n$, the response $\xi_{n,k}$ is drawn
from the arm's reward kernel independently of the past given the arm, so its
conditional mean is $\theta_k$; pulling out the $\mathcal{G}_n$-measurable factor
$X_{n,k}$ gives $\E[\Delta Q_{n,k}\mid\mathcal{G}_n] = X_{n,k}(\theta_k-\theta_k)=0$.
The tower property ($\filt_{n-1}\subseteq\mathcal{G}_n$) yields
$\E[\Delta Q_{n,k}\mid\filt_{n-1}] = 0$.

This is formalized in the \texttt{LeanMachineLearning} algorithm--environment
framework: the arm is the action and the response is the feedback of a stationary
environment with per-arm reward kernel $\nu$, and $\theta_k = \int x\,\nu_k(dx)$ is
the kernel mean. The conditional-mean computation is \texttt{condExp\_feedback}, and
the two filtrations $\filt$ and $\mathcal{G}$ are \texttt{filtration} and
\texttt{filtrationAction}.
\end{proof}

\begin{lemma}[Quadratic variation of $Q$]\label{lem:Q_quad_var}
\uses{def:Q, def:counts}
Under Condition~\textbf{A}, the predictable quadratic variation of $Q_{\cdot,k}$
is $\langle Q_{\cdot,k}\rangle_n = V_k N_{n,k}$, and for $k \ne j$ the cross
variation $\langle Q_{\cdot,k}, Q_{\cdot,j}\rangle_n = 0$.
\end{lemma}

\begin{proof}
\uses{lem:Q_martingale, lem:qv_incr, lem:qv_orthogonal}
Sum the conditional second moments of Lemma~\ref{lem:Q_martingale} using
$\sum_{m\le n} X_{m,k} = N_{n,k}$. The cross terms vanish because
$X_{m,k}X_{m,j} = 0$ for $k \ne j$ (each patient is assigned to one arm).
\end{proof}

\begin{lemma}[Estimator error via $Q$]\label{lem:theta_error_Q}
\lean{AlphaRAR.theta_error_Q}\leanok
\uses{def:estimator, def:Q, def:bahadur}
On $\{N_{m,k}\ge1\}$, the leading term of the estimator error is
$\dfrac{1}{N_{m,k}}\sum_{j=1}^m X_{j,k}\xi_{j,k} - \theta_k
 = \dfrac{Q_{m,k}}{N_{m,k}}$.
Consequently, under the Bahadur representation
(Definition~\ref{def:bahadur}),
$\hat{\theta}_{m,k} - \theta_k = \dfrac{Q_{m,k}}{N_{m,k}} + o(N_{m,k}^{-1/2})$.
\end{lemma}

\begin{proof}
$\frac{1}{N_{m,k}}\sum_j X_{j,k}\xi_{j,k} - \theta_k
= \frac{1}{N_{m,k}}\sum_j X_{j,k}(\xi_{j,k}-\theta_k)
= \frac{Q_{m,k}}{N_{m,k}}$, using $\sum_j X_{j,k} = N_{m,k}$.
\leanok
\end{proof}

\section{The auxiliary process \texorpdfstring{$U$}{U} and the last hitting time}
\label{sec:aux_U}

\begin{definition}[Auxiliary process $U$]\label{def:U}
\lean{AlphaRAR.auxU}\leanok
\uses{def:M, def:plugin_target}
For a fixed $\alpha \in [0,1)$, $k \in [K]$ and $n \ge 1$,
\[
  U_{n,k} := \sum_{m=1}^{n-1} \alpha\, \hat{\rho}_{m,k} + M_{n,k} - n\, \hat{\rho}_{n,k}.
\]
\end{definition}

\begin{lemma}[Increment of $U$]\label{lem:U_increment}
\lean{AlphaRAR.auxU_succ_sub}\leanok
\uses{def:U}
For all $n \ge 1$ and $k \in [K]$,
\[
  \Delta U_{n,k} = U_{n,k} - U_{n-1,k}
  = \alpha\, \hat{\rho}_{n-1,k} - p_{n,k} + \Delta\big(N_{n,k} - n\hat{\rho}_{n,k}\big),
\]
where we used $\Delta M_{n,k} = X_{n,k} - p_{n,k}$ and $\Delta N_{n,k} = X_{n,k}$.
\end{lemma}

\begin{proof}\leanok
Direct computation from Definition~\ref{def:U}, expanding
$M_{n,k} - M_{n-1,k} = X_{n,k} - p_{n,k}$ and telescoping.
\end{proof}

\begin{definition}[Last under-sampling time]\label{def:hitting}
\lean{AlphaRAR.hitting}\leanok
\uses{def:counts, def:plugin_target}
For $k \in [K]$ and $n \ge 1$, the \emph{last under-sampling time} before $n$ is
\[
  \ell_{n,k} := \max\Big\{\, m \le n : \frac{N_{m,k}}{m} \le \hat{\rho}_{m,k} \,\Big\}.
\]
It is the last time up to $n$ at which arm $k$ is under-sampled (its proportion
is at most its estimated target).
\end{definition}

\begin{lemma}[Basic properties of the hitting time]\label{lem:hitting_basic}
\lean{AlphaRAR.hitting_basic}\leanok
\uses{def:hitting}
For each $k$, $(\ell_{n,k})_{n\ge1}$ is non-decreasing and $\ell_{n,k} \le n$.
\end{lemma}

\begin{proof}\leanok
Immediate from the definition: enlarging $n$ enlarges the set over which the max
is taken, and the max is bounded by $n$.
\end{proof}

\textbf{Lean remark.} Well-definedness of $\ell_{n,k}$ requires the set
$\{m\le n : N_{m,k}/m \le \hat\rho_{m,k}\}$ to be nonempty; a convenient
convention (used implicitly in the paper) is that it contains a base time such
as the end of the burn-in, or one sets $\ell_{n,k}=0$ when the set is empty.
The proofs only use Lemma~\ref{lem:hitting_basic} together with the sign
property of the next lemma.

\begin{lemma}[Sign at the hitting time]\label{lem:hitting_sign}
\lean{AlphaRAR.hitting_sign}\leanok
\uses{def:hitting, def:aRTS}
For $m > \ell_{n,k}$ (and $m \le n$), $\frac{N_{m,k}}{m} > \hat{\rho}_{m,k}$;
in particular, for an $\aRTS$ design with $m \ge Km_0$, $p_{m+1,k} \le \alpha\hat{\rho}_{m,k}$.
\end{lemma}

\begin{proof}\leanok
By maximality of $\ell_{n,k}$, no $m$ with $\ell_{n,k} < m \le n$ satisfies
$N_{m,k}/m \le \hat{\rho}_{m,k}$. The second statement is the throttling
condition \eqref{eq:throttle}.
\end{proof}

\chapter{Technical lemmas}
\label{chap:technical}

This chapter gathers the probabilistic and analytic lemmas that are used as
black boxes in the asymptotic analysis. They are self-contained and are natural
first formalization targets. It corresponds to Appendix~C of the paper.

\section{Probabilistic lemmas}
\label{sec:tech_prob}

\begin{lemma}[Bounded expectation gives $O_p$]\label{lem:expectation_to_O}
\uses{def:op_Op}
Let $(u_n)_{n\ge1}$ be positive reals and $(X_n)_{n\ge1}$ real random variables.
If there is $C > 0$ with $\E|X_n| \le C u_n$ for all $n$, then $X_n = O_p(u_n)$.
\end{lemma}

\begin{proof}
By Markov's inequality, $\Prob(|X_n| > M u_n) \le \E|X_n|/(M u_n) \le C/M$.
Given $\varepsilon > 0$, take $M \ge C/\varepsilon$ to get
$\Prob(|X_n| > M u_n) \le \varepsilon$ for all $n$, which is the definition of
$X_n = O_p(u_n)$.
\end{proof}

\begin{lemma}[Doob $L^2$ maximal bound for martingale differences]\label{lem:mart_maximal}
Let $\{M_n, \filt_n\}$ be a martingale with $\E|\Delta M_n|^2 \le C_0$ for all
$n$. Then for any integer $L$,
\[
  \E\Big[\max_{1 \le m \le L} |M_n - M_{n-m}|\Big] \le 4\sqrt{C_0 L}.
\]
\end{lemma}

\begin{proof}
Fix $n$ and set $S_k := M_{n-L+k} - M_{n-L}$ for $0 \le k \le L$, a martingale for
the shifted filtration $\mathcal{G}_k := \filt_{n-L+k}$. First,
$\max_{1\le m\le L}|M_n - M_{n-m}| \le 2\max_{1\le k\le L}|S_k|$ by the triangle
inequality. Since $\{S_k^2\}$ is a submartingale, Doob's maximal inequality gives
$\Prob(\max_k |S_k| \ge x) \le \E[S_L^2]/x^2$. Orthogonality of martingale
increments yields $\E[S_L^2] = \sum_{k=1}^L \E|\Delta M_{n-L+k}|^2 \le L C_0$.
Combining, $\Prob(\max_m |M_n - M_{n-m}| \ge x) \le 4 L C_0 / x^2$, and integrating
the tail,
$\E[\max_m |M_n-M_{n-m}|] \le \int_0^{2\sqrt{LC_0}} 1\,dx
 + 4\int_{2\sqrt{LC_0}}^\infty \frac{LC_0}{x^2}\,dx = 4\sqrt{LC_0}$.
\end{proof}

\begin{lemma}[Dyadic maximal bound]\label{lem:mart_maximal_dyadic}
Under the hypotheses of Lemma~\ref{lem:mart_maximal}, for any $L$ with
$4 < L < n$,
\[
  \E\Big[\max_{L \le m \le n} \frac{|M_n - M_{n-m}|}{m}\Big] \le 12\sqrt{\frac{C_0}{L}}.
\]
\end{lemma}

\begin{proof}
\uses{lem:mart_maximal}
Split the range $[L,n]$ into dyadic blocks $2^{j-1} \le m < 2^j$. On such a block
$\frac1m \le \frac{2}{2^j}$, so
$\E[\max_{L\le m\le n}\frac{|M_n-M_{n-m}|}{m}]
 \le \sum_{j \ge \log_2 L} \frac{2}{2^j}\, \E[\max_{m<2^j}|M_n-M_{n-m}|]
 \le \sum_{j\ge\log_2 L} \frac{2}{2^j}\, 4\sqrt{C_0 2^j}$
by Lemma~\ref{lem:mart_maximal}. The geometric sum is bounded by
$12\sqrt{C_0/L}$.
\end{proof}

\begin{lemma}[Increment control from maximal bounds]\label{lem:increment_control}
For any sequence $(Q_n)_{n\ge0}$ in a normed space, any $\ell$ with
$0 \le \ell \le n$ and any $L < n$,
\[
  \|Q_n - Q_\ell\| \le (n - \ell)\max_{L \le m \le n}\frac{\|Q_n - Q_{n-m}\|}{m}
   + \max_{m < L}\|Q_n - Q_{n-m}\|.
\]
\end{lemma}

\begin{proof}
Let $d := n - \ell$. If $d < L$, then $\|Q_n - Q_\ell\| = \|Q_n - Q_{n-d}\| \le
\max_{m<L}\|Q_n-Q_{n-m}\|$. If $d \ge L$, then
$\|Q_n - Q_\ell\| = d\cdot \frac{\|Q_n-Q_{n-d}\|}{d}
 \le (n-\ell)\max_{L\le m\le n}\frac{\|Q_n-Q_{n-m}\|}{m}$.
Taking the larger of the two bounds proves the claim.
\end{proof}

\begin{lemma}[Kolmogorov-type inequality for martingales]\label{lem:chebyshev_mart}
Let $\{X_n\}$ be a square-integrable martingale with $\E[X_1] = 0$. Then for
$\lambda > 0$,
\[
  \Prob\Big(\max_{1\le i\le n} X_i \ge \lambda\Big)
  \le \frac{\Var(X_n)}{\Var(X_n) + \lambda^2}.
\]
\end{lemma}

\begin{proof}
For $\alpha > 0$, $\{\max_k X_k \ge \lambda\} \subseteq
\{\max_k (X_k+\alpha)^2 \ge (\lambda+\alpha)^2\}$. As $\{(X_k+\alpha)^2\}$ is a
submartingale, Doob's inequality gives
$\Prob(\max_k(X_k+\alpha)^2 \ge (\lambda+\alpha)^2)
 \le \frac{\E[(X_n+\alpha)^2]}{(\lambda+\alpha)^2}
 = \frac{\Var(X_n)+\alpha^2}{(\lambda+\alpha)^2}$.
Optimizing at $\alpha = \Var(X_n)/\lambda$ gives
$\frac{\Var(X_n)}{\Var(X_n)+\lambda^2}$.
\end{proof}

\begin{corollary}[Martingale is $O_p(\sqrt n)$]\label{cor:mart_Op}
\uses{def:op_Op}
Let $\{X_n\}$ be a square-integrable martingale with $\E[X_1]=0$ and stationary
increment variances, i.e.\ $\Var(X_n) = n\sigma^2$ for some $\sigma^2 < \infty$.
Then $X_n = O_p(\sqrt n)$.
\end{corollary}

\begin{proof}
\uses{lem:chebyshev_mart}
By Lemma~\ref{lem:chebyshev_mart},
$\Prob(X_n \ge M\sqrt n) \le \frac{n\sigma^2}{n\sigma^2 + M^2 n}
 = \frac{\sigma^2}{\sigma^2 + M^2}$, uniformly in $n$; the right-hand side is
$\le \varepsilon$ for $M$ large. Applying the same to $-X_n$ gives boundedness in
probability of $X_n/\sqrt n$.
\end{proof}

\section{Exponential families}
\label{sec:tech_expfam}

\begin{proposition}[Variance equals inverse Fisher information]\label{prop:exp_fam}
For a regular one-parameter exponential family with density
$f(y\mid\eta) = \exp\{\eta y - A(\eta)\}h(y)$, if the parameter of interest
$\theta_k = A'(\eta_k)$ is the mean parameter and $\xi_{1,k} = Y$, then
$\Var(\xi_{1,k}) = I_k(\theta_k)^{-1}$, where $I_k$ is the Fisher information for
one observation.
\end{proposition}

\begin{proof}
For the natural family, $\E_\eta[Y] = A'(\eta)$ and $\Var_\eta(Y) = A''(\eta)$.
With $\theta_k = A'(\eta_k)$ one has $\Var(\xi_{1,k}) = A''(\eta_k)$. By the
reparameterization formula for Fisher information,
$I_k(\theta_k) = I_{\eta_k}(\eta_k)\big(\frac{d\eta_k}{d\theta_k}\big)^2
 = A''(\eta_k)\big(\frac{1}{A''(\eta_k)}\big)^2 = \frac{1}{A''(\eta_k)}$,
so $\Var(\xi_{1,k}) = I_k(\theta_k)^{-1}$.
\end{proof}

\section{An analytic convergence lemma}
\label{sec:tech_analytic}

\begin{lemma}[Positive-part convergence]\label{lem:convergence}
Let $(a_n)_{n\ge1}$ be a non-decreasing sequence with $a_n \le n$, let
$\alpha, v \in [0,1]$, and let $(X_n)_{n\ge1}$ be real with
$\lim_{n\to\infty} \frac{X_n}{n} = -\alpha v$. Then for any real sequence
$\varepsilon_n$ with $\limsup_n \varepsilon_n \le 0$,
\[
  \lim_{n\to\infty}\Big(\varepsilon_n + \frac{X_n - X_{a_n}}{n}\Big)^+ = 0.
\]
\end{lemma}

\begin{proof}
If $(a_n)$ is bounded, it is eventually constant, so $X_{a_n}/n \to 0$ and
$\varepsilon_n + \frac{X_n - X_{a_n}}{n} \to -\alpha v \le 0$, whence the positive
part tends to $0$. If $a_n \to \infty$, then $X_{a_n}/a_n \to -\alpha v$; for
$\varepsilon>0$ and $n$ large, $-\frac{X_{a_n}}{n} \le (\varepsilon + \alpha v)\frac{a_n}{n}
 \le \varepsilon + \alpha v$ (using $a_n \le n$), so
$\varepsilon_n + \frac{X_n - X_{a_n}}{n} \le \varepsilon_n + \frac{X_n}{n} + \varepsilon + \alpha v$.
Since $\frac{X_n}{n} + \alpha v \to 0$ and $\limsup \varepsilon_n \le 0$, the
limsup of the left side is $\le \varepsilon$; as $\varepsilon$ is arbitrary the
positive part tends to $0$.
\end{proof}

\chapter{Generic sufficient conditions}
\label{chap:generic}

The heart of the paper is a set of abstract sufficient conditions, phrased in
terms of the auxiliary process $U$ and an arbitrary sequence of times
$\ell_{n,k}$, under which the main asymptotic theorems hold. The two concrete
families ($\aRTS$ in Chapter~\ref{chap:consistency}--\ref{chap:normality} and
$\aRTSFE$ in Chapter~\ref{chap:forced}) are then handled by exhibiting a suitable
$\ell_{n,k}$. This chapter states the abstract conditions and verifies them for
the $\aRTS$ family. It corresponds to Lemma~3.2 and Lemma~A.1 of the paper.

\section{The abstract conditions}
\label{sec:generic_conditions}

\begin{definition}[Generic condition set]\label{def:generic_cond}
\uses{def:counts, def:plugin_target, def:U}
Let $(\ell_{n,k})_{n\ge1, k\in[K]}$ be integer-valued random variables. We say
that a RAR procedure \emph{satisfies the generic conditions with $\ell$} if:
\begin{enumerate}
  \item[(i)] for all $k$, $(\ell_{n,k})_n$ is non-decreasing and $\ell_{n,k}\le n$;
  \item[(ii)] for all $k,n$,
        \begin{equation}\label{eq:generic_ineq}
          N_{n,k} - n\hat{\rho}_{n,k}
          \le 1 + N_{\ell_{n,k},k} - \ell_{n,k}\hat{\rho}_{\ell_{n,k},k}
             + U_{n,k} - U_{\ell_{n,k},k};
        \end{equation}
  \item[(iii)] for all $k,n$,
        \begin{equation}\label{eq:generic_small}
          N_{\ell_{n,k},k} - \ell_{n,k}\hat{\rho}_{\ell_{n,k},k} = o(\sqrt n)
          \quad\text{a.s.}
        \end{equation}
\end{enumerate}
\end{definition}

\begin{theorem}[Generic conditions imply the main theorems]\label{thm:generic_main}
\lean{AlphaRAR.consistency_of_generic_ae}
\uses{def:generic_cond, def:condA, def:condB, lem:consistency_matching}
Suppose a RAR procedure satisfies the generic conditions with some $\ell$, and
that Conditions~\textbf{A}--\textbf{B} hold. Then the conclusions of the strong
consistency theorem (Theorem~\ref{thm:LLN}) and of the asymptotic normality
theorem (Theorem~\ref{thm:normality}) hold. (The Lean statement
\texttt{consistency\_of\_generic\_ae} formalizes the strong-consistency direction:
the generic conditions and the plug-in-target limit imply $N_{n,k}/n \to u_k$ a.s.\ for every
arm; the asymptotic-normality direction is not yet formalized.)
\end{theorem}

\begin{proof}\uses{thm:LLN, thm:normality}
This is the content of Chapters~\ref{chap:consistency} and \ref{chap:normality}:
those chapters derive Theorem~\ref{thm:LLN} and Theorem~\ref{thm:normality}
using only the generic conditions of Definition~\ref{def:generic_cond} (together
with Conditions~\textbf{A}--\textbf{B}), never the specific form of the design.
\end{proof}

\textbf{Lean remark.} Theorem~\ref{thm:generic_main} is the key modularization:
it should be a standalone Lean theorem taking a hypothesis bundle
\texttt{GenericConditions} and returning the two conclusions. Both the $\aRTS$
and $\aRTSFE$ results then reduce to constructing $\ell$ and discharging
Definition~\ref{def:generic_cond}.

\section{Verification for the \texorpdfstring{$\aRTS$}{aRTS} family}
\label{sec:generic_verify}

We now show that the last under-sampling time $\ell_{n,k}$ of
Definition~\ref{def:hitting} satisfies the generic conditions for any $\aRTS$
design. The proof of the key inequality \eqref{eq:generic_ineq} is isolated into
a telescoping identity.

\begin{lemma}[Telescoping identity for $U$]\label{lem:U_telescope}
\lean{AlphaRAR.auxU_telescope}\leanok
\uses{def:U, def:hitting, def:aRTS}
For any $\aRTS$ design and any $\ell_{n,k} \le n$,
\[
  N_{n,k} - n\hat{\rho}_{n,k}
  = N_{\ell_{n,k},k} - \ell_{n,k}\hat{\rho}_{\ell_{n,k},k}
    + \sum_{m=\ell_{n,k}}^{n-1}\big(p_{m,k} - \alpha\hat{\rho}_{m,k}\big)
    + U_{n,k} - U_{\ell_{n,k},k}.
\]
\end{lemma}

\begin{proof}\leanok
\uses{lem:U_increment}
Summing the increment identity of Lemma~\ref{lem:U_increment} from
$m = \ell_{n,k}$ to $n-1$ gives
$U_{n,k} - U_{\ell_{n,k},k}
 = \sum_{m=\ell_{n,k}}^{n-1} (\alpha\hat{\rho}_{m,k} - p_{m,k})
   + (N_{n,k}-n\hat{\rho}_{n,k}) - (N_{\ell_{n,k},k}-\ell_{n,k}\hat{\rho}_{\ell_{n,k},k})$.
Rearranging yields the claim.
\end{proof}

\begin{lemma}[Key inequality for $\aRTS$]\label{lem:preliminary_ineq}
\lean{AlphaRAR.preliminary_ineq}\leanok
\uses{def:aRTS, def:hitting, def:U}
For any $\aRTS$ design, inequality \eqref{eq:generic_ineq} holds with the last
under-sampling time of Definition~\ref{def:hitting}:
\[
  N_{n,k} - n\hat{\rho}_{n,k}
  \le 1 + N_{\ell_{n,k},k} - \ell_{n,k}\hat{\rho}_{\ell_{n,k},k}
    + U_{n,k} - U_{\ell_{n,k},k}.
\]
\end{lemma}

\begin{proof}\leanok
\uses{lem:U_telescope, lem:hitting_sign}
In Lemma~\ref{lem:U_telescope}, each summand $p_{m,k} - \alpha\hat{\rho}_{m,k}$
with $\ell_{n,k} < m \le n-1$ is $\le 0$: indeed $m > \ell_{n,k}$ so arm $k$
is over-sampled at time $m$ and the throttling condition
(Lemma~\ref{lem:hitting_sign}) gives $p_{m,k} \le \alpha\hat{\rho}_{m,k}$.
Dropping this nonpositive sum and bounding the single leftover boundary term
$p_{\ell_{n,k},k} \le 1$ gives the inequality.
\end{proof}

\begin{lemma}[Smallness at the hitting time for $\aRTS$]\label{lem:preliminary_small}
\lean{AlphaRAR.preliminary_small}\leanok
\uses{def:hitting, def:aRTS}
For any $\aRTS$ design, \eqref{eq:generic_small} holds:
$N_{\ell_{n,k},k} - \ell_{n,k}\hat{\rho}_{\ell_{n,k},k} \le Km_0 = o(\sqrt n)$.
\end{lemma}

\begin{proof}
If $\ell_{n,k} \ge Km_0+1$, then by definition of the hitting time
$N_{\ell_{n,k},k}/\ell_{n,k} \le \hat{\rho}_{\ell_{n,k},k}$, so the difference is
$\le 0$. If $\ell_{n,k} < Km_0+1$, then $N_{\ell_{n,k},k} \le \ell_{n,k} \le Km_0$,
so the difference is $\le Km_0$. Either way it is bounded by the constant $Km_0$,
which is $o(\sqrt n)$.
\leanok
\end{proof}

\begin{proposition}[$\aRTS$ satisfies the generic conditions]\label{prop:aRTS_generic}
\lean{AlphaRAR.generic_ineq_of_hitting, AlphaRAR.generic_small_of_hitting}
\uses{def:aRTS, def:generic_cond, def:hitting}
Every $\aRTS$ design satisfies the generic conditions
(Definition~\ref{def:generic_cond}) with $\ell_{n,k}$ the last under-sampling
time of Definition~\ref{def:hitting}.
\end{proposition}

\begin{proof}
\uses{lem:hitting_basic, lem:preliminary_ineq, lem:preliminary_small}
Condition (i) is Lemma~\ref{lem:hitting_basic}, condition (ii) is
Lemma~\ref{lem:preliminary_ineq}, and condition (iii) is
Lemma~\ref{lem:preliminary_small}. (The Lean lemmas
\texttt{generic\_ineq\_of\_hitting} and \texttt{generic\_small\_of\_hitting} package
Lemmas~\ref{lem:preliminary_ineq}--\ref{lem:preliminary_small} at the last under-sampling time
$\ell_{n,k}$ into the operational forms \eqref{eq:generic_ineq}--\eqref{eq:generic_small},
given the $\aRTS$ throttle $\neg P_m \Rightarrow p_{m}\le\alpha\hat\rho_m$ (patient $m$ throttled
by the target at the same index $m$). A single boundary summand $p_{\ell,k}-\alpha\hat\rho_{\ell,k}$
is uncontrolled by the throttle, so the constant in \eqref{eq:generic_ineq} is $1$. No lower bound
on $\ell_{n,k}$ is needed: the telescoping identity holds at $\ell_{n,k}=0$ as well (the
$\alpha\hat\rho$ sum of $U$ grows at every step), so the $0$-indexed edge case is automatic. The
remaining step is to verify that a concrete $\aRTS$ design's selection probabilities satisfy the
throttle — this is Definition~\ref{def:aRTS_alg} and Corollary~\ref{cor:aRTS_alg_consistency}.)
\end{proof}

\begin{corollary}[Strong consistency of the \texorpdfstring{$\aRTS$}{aRTS} proportions]
\label{cor:aRTS_consistency}
\lean{AlphaRAR.aRTS_consistency}
\uses{prop:aRTS_generic, thm:generic_main, lem:rho_converges, lem:M_lln}
Consider an $\aRTS$ design driven by an algorithm–environment sequence under a stationary
environment, with $Y_n\in L^2$ (Condition~\textbf{A}) and a continuous simplex-valued target map
$\Target$. If the design's selection probabilities satisfy the throttle
$\neg P_{m,k}\Rightarrow p_{m,k}\le\alpha\hat\rho_{m,k}$ a.s., then almost surely there is a common
limit vector $v$ with $N_{n,k}/n\to v_k$ and $\hat\rho_{n,k}\to v_k$ for every arm $k$.
\end{corollary}

\begin{proof}
\uses{prop:aRTS_generic, thm:generic_main}
Instantiate Theorem~\ref{thm:generic_main} with the assignment indicators
$Y_{n,k}=\mathbf 1\{A_n=k\}$, the selection probabilities
$p_{n,k}=P[\mathbf 1\{A_n=k\}\mid\filt_{n-1}]$, the plug-in target
$\hat\rho_{n,k}=\Target(\estParams_n)_k$, and $\ell_{n,k}$ the last under-sampling time. The generic
conditions are discharged by Proposition~\ref{prop:aRTS_generic} (the throttle is the stated design
hypothesis); the vanishing normalized assignment martingale $M_{n,k}/n\to0$ is
Lemma~\ref{lem:M_lln}, and the plug-in-target convergence $\hat\rho_{n,k}\to v_k$ is
Lemma~\ref{lem:rho_converges}. The common limit $v_k$ is that of $\hat\rho_{\cdot,k}$.
\end{proof}

\section{The \texorpdfstring{$\aRTS$}{aRTS} family as a property of an algorithm}
\label{sec:aRTS_alg}

Corollary~\ref{cor:aRTS_consistency} takes the process-level throttle as a hypothesis. We now
express it as a membership condition on the driving algorithm and close the loop.

\begin{definition}[$\aRTS$ algorithm]\label{def:aRTS_alg}
\lean{AlphaRAR.IsARTS}\leanok
\uses{def:aRTS}
An algorithm chooses its next action from a history $h$ of action–reward pairs for patients
$0,\dots,n$ via a Markov kernel; the probability it assigns to arm $k$ is
$\mathrm{policy}_n(h)(\{k\})$. The algorithm
is an \emph{$\alpha$-throttled $\aRTS$ design} with offsets $\Params_0$ and target map $\Target$ if
its policy throttles every over-sampled arm: for every history $h$ and arm $k$, if
$N_{n+1,k}(h) > (n+1)\,\hat\rho_k(h)$ then $\mathrm{policy}_n(h)(\{k\}) \le \alpha\,\hat\rho_k(h)$,
where the count $N(h)$ and target $\hat\rho(h)=\Target(\estParams(h))$ are computed from $h$.
\end{definition}

\begin{corollary}[Consistency of an $\aRTS$ algorithm]\label{cor:aRTS_alg_consistency}
\lean{AlphaRAR.aRTS_consistency_of_isARTS}\leanok
\uses{def:aRTS_alg, cor:aRTS_consistency}
If the driving algorithm is an $\aRTS$ design (Definition~\ref{def:aRTS_alg}), under the hypotheses
of Corollary~\ref{cor:aRTS_consistency}, then almost surely $N_{n,k}/n\to v_k$ and
$\hat\rho_{n,k}\to v_k$ for every arm $k$.
\end{corollary}

\begin{proof}\leanok
\uses{cor:aRTS_consistency}
The policy-level throttle discharges the process-level throttle hypothesis of
Corollary~\ref{cor:aRTS_consistency}. The link is that the selection probability
$p_{n+1,k}=P[\mathbf 1\{A_{n+1}=k\}\mid\filt_n]$ equals $\mathrm{policy}_n(\text{history}_n)(\{k\})$
almost surely — a conditional expectation computed from the algorithm's conditional distribution of
the next action given the history (\texttt{HasCondDistrib.condExp\_comp\_eq}) — and the history-level
count and target agree with the process count $N_{n+1,k}$ and target $\hat\rho_{n+1,k}$.
\end{proof}

\chapter{Strong consistency and rates}
\label{chap:consistency}

This chapter proves the first main theorem: under the generic conditions and
Conditions~\textbf{A}--\textbf{B}, the allocation proportions converge almost
surely to the target and the estimators converge at the law-of-iterated-logarithm
rate. The proof is decomposed into small steps, each a candidate Lean lemma. It
corresponds to Theorem~3.1 and Appendix~A.2 of the paper.

Throughout the chapter we fix a RAR procedure satisfying the generic conditions
(Definition~\ref{def:generic_cond}) with some $\ell_{n,k}$, and we assume
Conditions~\textbf{A}--\textbf{B}. By Proposition~\ref{prop:aRTS_generic} this
covers every $\aRTS$ design.

\section{Convergence of the plug-in target}
\label{sec:cons_rho_limit}

\begin{lemma}[SLLN for the response martingale]\label{lem:respMart_div_count_slln}
\lean{AlphaRAR.respMart_div_pullCount_ae_tendsto_zero}\leanok
\uses{def:Q, def:counts, thm:mart_slln_bracket, lem:Q_quad_var}
For each arm $k$, on the event $\{N_{n,k}\to\infty\}$ the bracket-normalized response
martingale vanishes: $\dfrac{Q_{n,k}}{N_{n,k}}\to 0$ a.s.
\end{lemma}

\begin{proof}\leanok
\uses{thm:mart_slln_bracket, lem:Q_quad_var}
The response martingale $Q_{\cdot,k}$ has predictable quadratic variation
$\langle Q_k\rangle_n = V_k N_{n,k}$ (Lemma~\ref{lem:Q_quad_var}). If $V_k>0$, then on
$\{N_{n,k}\to\infty\}$ the bracket $\langle Q_k\rangle_n = V_k N_{n,k}\to\infty$, so the
bracket-normalized strong law (Theorem~\ref{thm:mart_slln_bracket}) gives
$Q_{n,k}/\langle Q_k\rangle_n\to0$, whence
$Q_{n,k}/N_{n,k} = V_k\,(Q_{n,k}/\langle Q_k\rangle_n)\to0$. If $V_k=0$, then
$\E[Q_{n,k}^2] = V_k\,\E[N_{n,k}] = 0$, so $Q_{n,k}=0$ a.s.\ for every $n$ and the ratio is
identically $0$.
\end{proof}

\begin{lemma}[Estimator limit on $\{N\to\infty\}$]\label{lem:theta_limit_infinite}
\lean{AlphaRAR.estimator_ae_tendsto_of_pullCount_atTop}\leanok
\uses{def:estimator, lem:estimator_bahadur, lem:respMart_div_count_slln}
For each arm $k$, on the event $\{N_{n,k}\to\infty\}$ the estimator converges a.s.\ to the arm
mean: $\hat\theta_{n,k}\to\theta_k$.
\end{lemma}

\begin{proof}\leanok
\uses{lem:estimator_bahadur, lem:respMart_div_count_slln}
By the exact error identity (Lemma~\ref{lem:estimator_bahadur}),
$\hat\theta_{n,k}-\theta_k = (Q_{n,k}+(\theta_{0,k}-\theta_k))/(N_{n,k}+1)$. On
$\{N_{n,k}\to\infty\}$ the constant offset is $O(N_{n,k}^{-1})\to0$, while
$\frac{Q_{n,k}}{N_{n,k}+1} = \frac{Q_{n,k}}{N_{n,k}}\cdot\frac{N_{n,k}}{N_{n,k}+1}\to0\cdot1=0$
by Lemma~\ref{lem:respMart_div_count_slln}. Hence $\hat\theta_{n,k}\to\theta_k$.
\end{proof}

\begin{lemma}[Estimator eventually constant on $\{\sup_n N<\infty\}$]\label{lem:theta_limit_finite}
\lean{AlphaRAR.exists_tendsto_estimator_of_not_pullCount_atTop}\leanok
\uses{def:estimator, def:counts}
For each arm $k$ and each realization, if $N_{n,k}\not\to\infty$ then $\hat\theta_{n,k}$ is
eventually constant, and in particular converges.
\end{lemma}

\begin{proof}\leanok
The count $N_{\cdot,k}$ is a nondecreasing $\N$-valued sequence, so $N_{n,k}\not\to\infty$
forces it to be bounded, hence eventually constant at its supremum $N_{N_0,k}$. For $n\ge N_0$
no further patient is assigned to arm $k$, so both the numerator $\sum_{j<n}X_{j,k}\xi_{j,k}$
and the denominator $N_{n,k}+1$ of $\hat\theta_{n,k}$ are frozen; hence $\hat\theta_{n,k}$ is
eventually constant.
\end{proof}

\begin{lemma}[The sequential estimator converges a.s.]\label{lem:theta_converges}
\lean{AlphaRAR.estimator_ae_tendsto}\leanok
\uses{lem:theta_limit_infinite, lem:theta_limit_finite}
For each arm $k$, the estimator $\hat\theta_{n,k}$ converges a.s.: to $\theta_k$ on
$\{N_{n,k}\to\infty\}$ and to an eventually-attained value on $\{\sup_n N_{n,k}<\infty\}$.
\end{lemma}

\begin{proof}\leanok
\uses{lem:theta_limit_infinite, lem:theta_limit_finite}
The alternative $N_{n,k}\to\infty$ versus $\sup_n N_{n,k}<\infty$ is exhaustive (as $N_{n,k}$
is monotone in $n$); branch~1 is Lemma~\ref{lem:theta_limit_infinite} and branch~2 is
Lemma~\ref{lem:theta_limit_finite}.
\end{proof}

\begin{lemma}[Dichotomy for the estimator limit]\label{lem:theta_limit_dichotomy}
\uses{def:estimator, def:condA, def:attainable, lem:theta_converges}
Under Condition~\textbf{A}, for each $k$ the estimator $\hat{\theta}_{n,k}$
converges a.s.\ to a limit in $I_k$: on $\{N_{n,k}\to\infty\}$ it converges to
$\theta_k$ by the strong law of large numbers; on $\{\sup_n N_{n,k} < \infty\}$
it is eventually constant and its limit lies in $V_k^{\mathrm{att}}$.
\end{lemma}

\begin{proof}
\uses{lem:theta_limit_infinite}
On $\{N_{n,k}\to\infty\}$, the estimator converges a.s.\ to $\theta_k$ by
Lemma~\ref{lem:theta_limit_infinite} (the martingale strong law applied to $Q_{\cdot,k}$). On
$\{\sup_n N_{n,k}<\infty\}$, only finitely many responses on arm $k$ are ever
observed, so $\hat{\theta}_{n,k}$ is eventually constant with a value in
$V_k^{\mathrm{att}}$.
\end{proof}

\begin{lemma}[The estimator vector converges a.s.]\label{lem:theta_converges_vec}
\lean{AlphaRAR.estimator_ae_tendsto_pi}\leanok
\uses{lem:theta_converges}
With finitely many arms, the per-arm a.s.\ limits of Lemma~\ref{lem:theta_converges} combine
into a single a.s.\ limit vector: almost surely there is $z\in\R^K$ with
$\estParams_n = (\hat\theta_{n,1},\dots,\hat\theta_{n,K}) \to z$.
\end{lemma}

\begin{proof}\leanok
\uses{lem:theta_converges}
Each component converges a.s.\ by Lemma~\ref{lem:theta_converges}; intersecting the finitely
many almost-sure events and choosing the (unique) limits componentwise gives a single a.s.\
event on which the whole vector converges.
\end{proof}

\begin{lemma}[Convergence of $\estTarget$ by continuity]\label{lem:rho_converges_cont}
\lean{AlphaRAR.rho_converges}\leanok
\uses{def:plugin_target, lem:theta_converges_vec}
For any continuous target map $\Target$, the plug-in target $\estTarget_n = \Target(\estParams_n)$
converges a.s.: almost surely there is $u\in\R^K$ with $\hat\rho_{n,k}\to u_k$ for every $k$,
where $u = \Target(z)$ and $z$ is the a.s.\ limit of $\estParams_n$.
\end{lemma}

\begin{proof}\leanok
\uses{lem:theta_converges_vec}
By Lemma~\ref{lem:theta_converges_vec}, $\estParams_n \to z$ a.s.\ (in the product topology of
$\R^K$). Continuity of $\Target$ transports this to
$\estTarget_n = \Target(\estParams_n) \to \Target(z) =: u$, and taking components gives
$\hat\rho_{n,k}\to u_k$.
\end{proof}

\begin{lemma}[Non-sparse limit of $\estTarget$]\label{lem:rho_converges}
\uses{def:condB, def:plugin_target, lem:rho_converges_cont}
Under Conditions~\textbf{A}--\textbf{B} there is a random $u \in (0,1)^K$ with
$\estTarget_n \to u$ a.s. Concretely $u = \Target(z)$ where $z\in I_1\times\dots\times I_K$
is the a.s.\ limit of $\estParams_n$.
\end{lemma}

\begin{proof}
\uses{lem:rho_converges_cont, lem:theta_limit_dichotomy}
By Lemma~\ref{lem:rho_converges_cont} (with $\Target$ continuous on $\hilb$ by Condition~\textbf{B})
and Lemma~\ref{lem:theta_limit_dichotomy}, $\estParams_n \to z$ a.s.\ with
$z \in I_1\times\dots\times I_K$. By Condition~\textbf{B}, $\Target$ is
continuous on the open set $\hilb \supseteq I_1\times\dots\times I_K$ and maps
$I_1\times\dots\times I_K$ into $(0,1)^K$; hence
$\estTarget_n = \Target(\estParams_n) \to \Target(z) =: u \in (0,1)^K$.
\end{proof}

\section{Convergence of the auxiliary process}
\label{sec:cons_U}

\begin{lemma}[LLN for the assignment martingale]\label{lem:M_lln}
\lean{AlphaRAR.assignMart_div_atTop_ae_tendsto_zero}\leanok
\uses{def:M}
$M_{n,k} = o(n)$ a.s.\ for each $k$.
\end{lemma}

\begin{proof}\leanok
\uses{lem:M_martingale, cor:mart_slln_bounded}
$(M_{n,k})$ is a martingale with bounded increments $|\Delta M_{m,k}|\le1$, so its
increments have summable normalized variances $\sum_m \frac{\Var(\Delta M_{m,k})}{m^2}<\infty$;
the strong law of large numbers for martingales gives $M_{n,k}/n \to 0$ a.s.
\end{proof}

\begin{lemma}[Limit of $U/n$]\label{lem:U_over_n}
\lean{AlphaRAR.auxU_div_tendsto}\leanok
\uses{def:U}
Under Conditions~\textbf{A}--\textbf{B}, for each $k$,
$\displaystyle \lim_{n\to\infty} \frac{U_{n,k}}{n} = -(1-\alpha)\,u_k$ a.s.,
where $u$ is the limit from Lemma~\ref{lem:rho_converges}. (The formalized statement
\texttt{auxU\_div\_tendsto} is the deterministic core: for real sequences with $\rho_n\to u$ and
$M_n/n\to0$, one has $U_n/n\to-(1-\alpha)u$; the a.s.\ statement is the pathwise application with the
a.s.\ limits from Lemmas~\ref{lem:M_lln} and~\ref{lem:rho_converges}.)
\end{lemma}

\begin{proof}\leanok
\uses{lem:M_lln, lem:rho_converges}
From Definition~\ref{def:U},
$\frac{U_{n,k}}{n} = \frac{\alpha}{n}\sum_{m=0}^{n-1}\hat{\rho}_{m,k}
 - \hat{\rho}_{n,k} + \frac{M_{n,k}}{n}$.
By Lemma~\ref{lem:M_lln} the last term is $o(1)$; by Lemma~\ref{lem:rho_converges}
$\hat{\rho}_{n,k}\to u_k$, and by Cesàro
$\frac1n\sum_{m<n}\hat{\rho}_{m,k}\to u_k$. Hence
$\frac{U_{n,k}}{n} \to \alpha u_k - u_k = -(1-\alpha)u_k$.
\end{proof}

\section{Matching of proportions and target}
\label{sec:cons_match}

\begin{lemma}[Positive part vanishes]\label{lem:pos_part_vanishes}
\lean{AlphaRAR.pos_part_vanishes, AlphaRAR.pos_part_vanishes_ae}\leanok
\uses{def:generic_cond}
Under the generic conditions and Conditions~\textbf{A}--\textbf{B}, for each $k$,
$\big(\frac{N_{n,k}}{n} - \hat{\rho}_{n,k}\big)^+ \to 0$ a.s. (The formalized statement
\texttt{pos\_part\_vanishes} is the pathwise core: it takes as hypotheses the plug-in convergence
$\rho_n\to u$ (\ref{lem:rho_converges}), the vanishing $M_n/n\to0$ (\ref{lem:M_lln}), and the two
generic conditions \eqref{eq:generic_ineq}--\eqref{eq:generic_small} in operational form.)
\end{lemma}

\begin{proof}\leanok
\uses{lem:convergence, lem:U_over_n}
Apply Lemma~\ref{lem:convergence} with $X_n := U_{n,k}$, $a_n := \ell_{n,k}$
(non-decreasing, $\le n$ by the generic conditions), and
$\varepsilon_n := \frac1n + \frac{N_{\ell_{n,k},k}-\ell_{n,k}\hat{\rho}_{\ell_{n,k},k}}{n}$,
which satisfies $\limsup\varepsilon_n \le 0$ by \eqref{eq:generic_small}. Since
$\frac{U_{n,k}}{n}\to -(1-\alpha)u_k$ (Lemma~\ref{lem:U_over_n}), the lemma yields
$\big(\varepsilon_n + \frac{U_{n,k}-U_{\ell_{n,k},k}}{n}\big)^+ \to 0$. Combined
with the generic inequality \eqref{eq:generic_ineq}, which bounds
$\frac{N_{n,k}}{n}-\hat{\rho}_{n,k}$ by that quantity, we get
$\big(\frac{N_{n,k}}{n}-\hat{\rho}_{n,k}\big)^+\to0$.
\end{proof}

\begin{lemma}[Negative part vanishes]\label{lem:neg_part_vanishes}
\lean{AlphaRAR.neg_part_vanishes, AlphaRAR.neg_part_vanishes_ae}\leanok
\uses{def:counts, def:plugin_target}
Under the same hypotheses, for each $k$,
$\big(\hat{\rho}_{n,k} - \frac{N_{n,k}}{n}\big)^+ \to 0$ a.s.
\end{lemma}

\begin{proof}\leanok
\uses{lem:pos_part_vanishes, lem:counts_sum}
Using $\sum_j N_{n,j}/n = 1$ (Lemma~\ref{lem:counts_sum}) and
$\sum_j \hat{\rho}_{n,j} = 1$,
$\big(\hat{\rho}_{n,k} - \frac{N_{n,k}}{n}\big)^+
 = \big(\sum_{j\ne k}(\frac{N_{n,j}}{n} - \hat{\rho}_{n,j})\big)^+
 \le \sum_{j\ne k}\big(\frac{N_{n,j}}{n} - \hat{\rho}_{n,j}\big)^+ \to 0$
by Lemma~\ref{lem:pos_part_vanishes}.
\end{proof}

\begin{lemma}[Proportions match the plug-in target]\label{lem:match}
\lean{AlphaRAR.match_proportion, AlphaRAR.match_proportion_ae}\leanok
\uses{def:counts, def:plugin_target}
Under the same hypotheses, for each $k$,
$\frac{N_{n,k}}{n} - \hat{\rho}_{n,k} \to 0$ and
$\frac{N_{n,k}}{n} \to u_k \in (0,1)$ a.s. (The formalized statement \texttt{match\_proportion}
delivers $\frac{N_{n,k}}{n}\to u_k$ from the two vanishing positive parts and $\hat{\rho}_{n,k}\to
u_k$.)
\end{lemma}

\begin{proof}\leanok
\uses{lem:pos_part_vanishes, lem:neg_part_vanishes, lem:rho_converges}
The absolute value $|\frac{N_{n,k}}{n}-\hat{\rho}_{n,k}|$ is the sum of the two
positive parts, both tending to $0$ by
Lemmas~\ref{lem:pos_part_vanishes}--\ref{lem:neg_part_vanishes}. Since
$\hat{\rho}_{n,k}\to u_k$ (Lemma~\ref{lem:rho_converges}), also
$\frac{N_{n,k}}{n}\to u_k$.
\end{proof}

\section{Identification of the limit and consistency}
\label{sec:cons_final}

\begin{lemma}[All arms sampled infinitely often]\label{lem:all_arms_infinite}
\lean{AlphaRAR.all_arms_infinite, AlphaRAR.all_arms_infinite_ae}\leanok
\uses{def:counts}
Under the same hypotheses, $N_{n,k}\to\infty$ a.s.\ for every $k$.
\end{lemma}

\begin{proof}\leanok
\uses{lem:match}
Since $\frac{N_{n,k}}{n}\to u_k > 0$ (Lemma~\ref{lem:match}), $N_{n,k}\sim u_k n
\to\infty$.
\end{proof}

\begin{lemma}[A.s. consistency of the proportions]\label{lem:consistency_matching}
\lean{AlphaRAR.consistency_ae}\leanok
\uses{lem:pos_part_vanishes, lem:neg_part_vanishes, lem:match, lem:rho_converges_cont}
Given the a.s.\ vanishing positive gaps for every arm (Lemma~\ref{lem:pos_part_vanishes}) and the
plug-in-target limits $\hat\rho_{n,k}\to u_k$ (Lemma~\ref{lem:rho_converges_cont}), together with
the simplex constraints ($\sum_k X_{j,k}=1$ and $\sum_k\hat\rho_{n,k}=1$), the proportions converge
a.s.\ to the target: $N_{n,k}/n \to u_k$ for all arms $k$ simultaneously.
\end{lemma}

\begin{proof}\leanok
\uses{lem:neg_part_vanishes, lem:match}
For each arm the negative gap vanishes by Lemma~\ref{lem:neg_part_vanishes}, and
Lemma~\ref{lem:match} then yields $N_{n,k}/n\to u_k$; intersecting the finitely many a.s.\ events
over the arms gives the simultaneous statement.
\end{proof}

\begin{lemma}[Estimator consistency]\label{lem:theta_consistent}
\uses{def:estimator, def:target, lem:consistency_matching}
Under the same hypotheses, $\estParams_n \to \Params$ a.s., and consequently
$\hat{\rho}_n \to v = \Target(\Params)$ and $\frac{N_{n,k}}{n}\to v_k$ a.s.
\end{lemma}

\begin{proof}
\uses{lem:all_arms_infinite, lem:theta_limit_dichotomy, lem:match}
By Lemma~\ref{lem:all_arms_infinite} every arm is sampled infinitely often, so
the first branch of Lemma~\ref{lem:theta_limit_dichotomy} applies to every $k$:
$\hat{\theta}_{n,k}\to\theta_k$. Continuity of $\Target$ gives
$\hat{\rho}_n\to\Target(\Params)=v$, and by Lemma~\ref{lem:match} the limit $u$ of
$N_n/n$ equals $v$.
\end{proof}

\section{Convergence rates}
\label{sec:cons_rates}

\begin{lemma}[LIL rate for the estimator]\label{lem:theta_LIL}
\lean{AlphaRAR.abs_estimator_sub_le_rate, AlphaRAR.abs_estimator_sub_le_rate_ae}\leanok
\uses{def:estimator}
Under the same hypotheses,
$\estParams_n - \Params = O\!\big(\sqrt{\tfrac{\log n}{n}}\big)$ a.s. (the $\log$, not $\log\log$,
form; see the LIL chapter's remark on the constant). The formalized statement
\texttt{abs\_estimator\_sub\_le\_rate} is the pathwise core: from the exact error identity
(Lemma~\ref{lem:estimator_bahadur}), a two-sided bound $|Q_{n,k}|\le C\sqrt{n\log n}$
(Lemma~\ref{lem:lil_truncation}) and $N_{n,k}/n\to v_k>0$ (Lemma~\ref{lem:match}), it gives
$|\hat\theta_{n,k}-\theta_k|\le C'\,\sqrt{n\log n}/n = C'\sqrt{\log n/n}$ for large $n$.
\end{lemma}

\begin{proof}\leanok
\uses{lem:all_arms_infinite, lem:Q_quad_var, lem:estimator_bahadur, lem:lil_truncation, lem:match}
By Lemma~\ref{lem:estimator_bahadur} the error is exactly
$\hat{\theta}_{n,k}-\theta_k = (Q_{n,k}+(\theta_{0,k}-\theta_k))/(N_{n,k}+1)$. The martingale
$Q_{\cdot,k}$ has quadratic variation $V_k N_{n,k}$ (Lemma~\ref{lem:Q_quad_var}), so the
finite-variance LIL (Lemma~\ref{lem:lil_truncation}, applied two-sided) gives
$|Q_{n,k}| = O(\sqrt{n\log n})$ a.s. Since $N_{n,k}\sim v_k n$ (Lemma~\ref{lem:match}), $N_{n,k}+1
\gtrsim (v_k/2)n$, and dividing bounds $|\hat{\theta}_{n,k}-\theta_k| \le
(2/v_k)(C+|\theta_{0,k}-\theta_k|)\sqrt{n\log n}/n = O(\sqrt{\log n / n})$.
\end{proof}

\begin{lemma}[Rate for the plug-in target]\label{lem:rho_rate}
\uses{def:plugin_target}
Under the same hypotheses,
$\hat{\rho}_n - v = O\!\big(\sqrt{\tfrac{\log\log n}{n}}\big)$ and hence
$n(\hat{\rho}_n - v) = O(\sqrt{n\log\log n})$ a.s.
\end{lemma}

\begin{proof}
\uses{lem:theta_LIL, def:condB, lem:taylor_rho}
A first-order Taylor expansion of $\Target$ around $\Params$ (twice
differentiable by Condition~\textbf{B}) gives
$\hat{\rho}_n - v = G(\estParams_n-\Params) + O(\|\estParams_n-\Params\|^2)$
with $G := \nabla\Target|_\Params$. Both terms are
$O(\sqrt{\log\log n / n})$ by Lemma~\ref{lem:theta_LIL}.
\end{proof}

\begin{theorem}[Strong consistency and rates]\label{thm:LLN}
\uses{def:generic_cond, def:condA, def:condB, def:counts, def:plugin_target}
Under the generic conditions and Conditions~\textbf{A}--\textbf{B}, for every
$k\in[K]$, as $n\to\infty$,
\[
  \frac{N_{n,k}}{n}\xrightarrow{a.s.} v_k, \qquad
  \estParams_n - \Params = O\!\Big(\sqrt{\tfrac{\log\log n}{n}}\Big)\ \text{a.s.},\qquad
  n(\hat{\rho}_n - v) = O(\sqrt{n\log\log n})\ \text{a.s.}
\]
\end{theorem}

\begin{proof}
\uses{lem:theta_consistent, lem:theta_LIL, lem:rho_rate}
Combine Lemma~\ref{lem:theta_consistent} (first statement), Lemma~\ref{lem:theta_LIL}
(second) and Lemma~\ref{lem:rho_rate} (third).
\end{proof}

\chapter{Asymptotic normality}
\label{chap:normality}

This chapter proves the second main theorem: refined $o_p(\sqrt n)$ rates for the
deviation between proportions and plug-in targets, and joint central limit
theorems for the estimators and the allocation proportions. As before, everything
is derived from the generic conditions, so it applies to every $\aRTS$ design via
Proposition~\ref{prop:aRTS_generic}. It corresponds to Theorem~3.2 and
Appendix~A.3 of the paper.

We keep the standing assumptions of Chapter~\ref{chap:consistency}: a RAR
procedure satisfying the generic conditions with some $\ell_{n,k}$, and
Conditions~\textbf{A}--\textbf{B}. We introduce the notation
\[
  G := \nabla\Target\big|_\Params, \qquad
  V := \mathrm{diag}\!\Big(\tfrac{V_1}{v_1}, \dots, \tfrac{V_K}{v_K}\Big), \qquad
  \Omega := \begin{pmatrix} GVG^\top & GVG^\top\\ GVG^\top & GVG^\top\end{pmatrix}.
\]

\section{A refined decomposition of \texorpdfstring{$U$}{U}}
\label{sec:norm_U}

\begin{lemma}[Additive decomposition of the $U$-increment]\label{lem:diff_U_decomp}
\lean{AlphaRAR.diff_U_decomp}\leanok
\uses{def:U}
For any non-decreasing sequence $(\ell_n)$ with $\ell_n\le n$, under
Conditions~\textbf{A}--\textbf{B}, for each $k$,
\[
  U_{n,k} - U_{\ell_n,k}
  = (n-\ell_n)\big[-(1-\alpha)v_k + o(1)\big]
    + (M_{n,k} - M_{\ell_n,k})
    + \ell_n\big(\hat{\rho}_{\ell_n,k} - \hat{\rho}_{n,k}\big).
\]
\end{lemma}

\begin{proof}\leanok
\uses{thm:LLN}
Expanding $U_{n,k}-U_{\ell_n,k}$ from Definition~\ref{def:U} and grouping the
target terms around $v_k$,
$U_{n,k}-U_{\ell_n,k}
 = (n-\ell_n)\big[-(1-\alpha)v_k - (\hat{\rho}_{n,k}-v_k)
   + \frac{\alpha}{n-\ell_n}\sum_{m=\ell_n}^{n-1}(\hat{\rho}_{m,k}-v_k)\big]
   + M_{n,k}-M_{\ell_n,k} + \ell_n(\hat{\rho}_{\ell_n,k}-\hat{\rho}_{n,k})$.
By Theorem~\ref{thm:LLN}, $\hat{\rho}_n\to v$, so both $\hat{\rho}_{n,k}-v_k$ and
the Cesàro average $\frac1{n-\ell_n}\sum_{m=\ell_n}^{n-1}(\hat{\rho}_{m,k}-v_k)$
are $o(1)$ (treating separately the bounded and diverging cases of $\ell_n$).

\emph{Formalization.} \texttt{AlphaRAR.diff\_U\_decomp}. The exact $M$-explicit identity is
\texttt{AlphaRAR.auxU\_sub}; the $o(1)$ remainder $\alpha\sum_{[\ell_n,n)}(\hat\rho_m-v)
-(n-\ell_n)(\hat\rho_n-v)$ is made rigorous as ``for every $\varepsilon>0$, eventually
$|\cdot|\le\varepsilon(n-\ell_n)$'', its first term controlled by the windowed Cesàro bound
\texttt{AlphaRAR.abs\_sum\_Ico\_sub\_le\_of\_tendsto} (valid for an arbitrary window, bounded or
diverging) and its second by $\hat\rho_n\to v$.
\end{proof}

\begin{lemma}[Control of $\ell_n(\hat{\rho}_{\ell_n}-\hat{\rho}_n)$]\label{lem:ell_rho_control}
\lean{AlphaRAR.ell_rho_control}\leanok
\uses{def:plugin_target, def:Q}
Under Conditions~\textbf{A}--\textbf{B} there is a constant $C\ge0$ with, for any
non-decreasing $(\ell_n)$, $\ell_n\le n$,
\[
  \big|\ell_n(\hat{\rho}_{\ell_n} - \hat{\rho}_n)\big|
  \le o(1)(n-\ell_n) + C\|Q_n - Q_{\ell_n}\| + o_p(\sqrt n).
\]
\end{lemma}

\begin{proof}\leanok
\uses{lem:theta_error_Q, lem:rho_rate, cor:mart_Op, lem:taylor_rho}
By Corollary~\ref{cor:mart_Op}, $Q_{n,k}=O_p(\sqrt n)$, hence
$\hat{\theta}_{n,k}-\theta_k = Q_{n,k}/N_{n,k} + o_p(N_{n,k}^{-1/2}) = O_p(n^{-1/2})$
using $N_{n,k}\sim v_k n$ (Lemma~\ref{lem:theta_error_Q}). The Taylor expansion of
$\Target$ gives $\hat{\rho}_n - v = G(\estParams_n-\Params)+O_p(1/n)$ and similarly
at $\ell_n$; subtracting,
$\ell_n(\hat{\rho}_{\ell_n}-\hat{\rho}_n)
 = \ell_n G(\estParams_{\ell_n}-\estParams_n) + O_p(1)$, so
$|\ell_n(\hat{\rho}_{\ell_n}-\hat{\rho}_n)| \le C\ell_n\|\estParams_{\ell_n}-\estParams_n\| + O_p(1)$.
Writing the estimator difference through $Q$ (splitting into the increment
$\sum_{j=\ell_n+1}^n$ and the reweighting of $\sum_{j\le\ell_n}$) and bounding
each piece using $N_{n,k}-N_{\ell_n,k}\le n-\ell_n$ and the SLLN, the three
resulting terms are respectively $C\|Q_n-Q_{\ell_n}\|$, $o(1)(n-\ell_n)$ and
$o_p(\sqrt n)$.

\emph{Formalization.} \texttt{AlphaRAR.ell\_rho\_control} produces the bound in the shape
$\ell_n|\hat\rho_{\ell_n,k}-\hat\rho_{n,k}| \le V_\rho(n-\ell_n)+W_\rho$ with $V_\rho=o_p(1)$,
$W_\rho=o_p(\sqrt n)$, which is what the drift-sign step consumes. The Condition-\textbf{B} input
enters as the single hypothesis \texttt{hlip} (the target is Lipschitz near $\Params$, so
$|\hat\rho_{\ell,k}-\hat\rho_{n,k}|\le L\sum_{k'}|\hat\theta_{\ell,k'}-\hat\theta_{n,k'}|$
eventually) --- mirroring how \texttt{clt\_rho} takes the derivative as a hypothesis. The
estimator-increment control $\ell|\hat\theta_{\ell,k'}-\hat\theta_{n,k'}|\le g_{k'}(n-\ell)+h_{k'}
|Q_{n,k'}-Q_{\ell,k'}|$ is \texttt{AlphaRAR.abs\_estimator\_diff\_le}; the coefficients are
$g_{k'}=o_p(1)$ (from the Doob-maximal bound \texttt{isBigOpOne\_sup'\_abs\_div\_sqrt} and the a.s.\
bound \texttt{isBigOpOne\_of\_ae\_bounded}) and the reweighting coefficient
$h_{k'}=\ell_n/(N_{n,k'}+1)$, which is only \emph{bounded in probability} ($O_p(1)$, converging a.s.\
to $1/v_{k'}$, again via \texttt{isBigOpOne\_of\_ae\_bounded}), not by a uniform constant; it
therefore enters the $o_p$ bookkeeping through $O_p\cdot o_p=o_p$
(\texttt{IsBigOpOne.mul\_littleOp}). The $Q$-increment control
$|Q_{n,k'}-Q_{\ell,k'}|\le Q_v(n-\ell)+Q_w$ is \texttt{qm\_increments\_resp}.
\end{proof}

\section{Sharp increment bounds}
\label{sec:norm_bounds}

\begin{lemma}[$Q$- and $M$-increments are $o_p(\sqrt n)$ plus drift]\label{lem:QM_increments}
\lean{AlphaRAR.qm_increments_resp, AlphaRAR.qm_increments_of_bdd,
AlphaRAR.norm_increment_le_vmaxSeq_wmaxSeq}\leanok
\uses{def:Q, def:M, def:hitting}
Under Conditions~\textbf{A}--\textbf{B}, for the hitting time $\ell_{n,k}$,
\[
  \|Q_n - Q_{\ell_{n,k}}\| \le o_p(1)(n-\ell_{n,k}) + o_p(\sqrt n),
  \qquad
  |M_{n,k} - M_{\ell_{n,k},k}| \le o_p(1)(n-\ell_{n,k}) + o_p(\sqrt n).
\]
\end{lemma}

\begin{proof}\leanok
\uses{lem:mart_maximal, lem:mart_maximal_dyadic, lem:mart_maximal_pi, lem:increment_control,
lem:expectation_to_O}
Pick $L_n\to\infty$ with $L_n = o(n)$. Lemma~\ref{lem:mart_maximal} and
Lemma~\ref{lem:mart_maximal_dyadic} bound the two maximal quantities
$\E[\max_{m\le L_n}\|Q_n-Q_{n-m}\|]\le 2\sqrt{L_n C_0}$ and
$\E[\max_{L_n\le m<n}\frac{\|Q_n-Q_{n-m}\|}{m}]\le 3\sqrt{C_0/L_n}$; by
Lemma~\ref{lem:expectation_to_O} these are $O_p(\sqrt{L_n})=o_p(\sqrt n)$ and
$O_p(1/\sqrt{L_n})=o_p(1)$ respectively. Plugging into the deterministic
inequality of Lemma~\ref{lem:increment_control} gives the bound for $Q$; the same
argument applied to $M$ gives the second bound.

\emph{Formalization.} The two maximal quantities are formalized as \texttt{wmaxSeq} and
\texttt{vmaxSeq} (with the concrete window $L_n=\lfloor\sqrt n\rfloor$), and the two $o_p$ statements
are \texttt{AlphaRAR.isLittleOpOne\_wmaxSeq\_div\_sqrt} ($o_p(\sqrt n)$) and
\texttt{AlphaRAR.isLittleOpOne\_vmaxSeq} ($o_p(1)$), proved from the \emph{vector} maximal bounds
(Lemma~\ref{lem:mart_maximal_pi}) and Markov's inequality. The deterministic increment control is
\texttt{AlphaRAR.norm\_increment\_le\_vmaxSeq\_wmaxSeq}. For $Q$ (the response martingale, a
square-integrable family) the $o_p$ statements are \texttt{AlphaRAR.qm\_increments\_resp}; for the
scalar assignment martingale $M$ (bounded increments $|\Delta M|\le1$) they are
\texttt{AlphaRAR.qm\_increments\_of\_bdd} (via the singleton family $\{M\}$).
\end{proof}

\begin{lemma}[Upper bound on the $U$-increment]\label{lem:U_increment_bound}
\lean{AlphaRAR.isLittleOpOne_max_of_decomp}
\uses{def:U, def:hitting}
Under Conditions~\textbf{A}--\textbf{B}, for every $k$,
\[
  U_{n,k} - U_{\ell_{n,k},k} \le o_p(\sqrt n),
  \qquad
  U_{n,k} - U_{\ell_{n,k},k} \le O(\sqrt{n\log\log n})\ \text{a.s.}
\]
The $o_p(\sqrt n)$ half is formalized (generically) as \texttt{AlphaRAR.isLittleOpOne\_max\_of\_decomp};
the a.s.\ $O(\sqrt{n\log\log n})$ half is not yet formalized.
\end{lemma}

\begin{proof}
\uses{lem:diff_U_decomp, lem:ell_rho_control, lem:QM_increments, thm:LLN, thm:lil_bounded}
From Lemma~\ref{lem:ell_rho_control} and the $Q$-bound of
Lemma~\ref{lem:QM_increments},
$|\ell_{n,k}(\hat{\rho}_{\ell_{n,k}}-\hat{\rho}_{n,k})| \le o_p(1)(n-\ell_{n,k}) + o_p(\sqrt n)$.
Substituting this and the $M$-bound into the decomposition of
Lemma~\ref{lem:diff_U_decomp},
$U_{n,k}-U_{\ell_{n,k},k} = (n-\ell_{n,k})[-(1-\alpha)v_k+o_p(1)] + o_p(\sqrt n)
 \le o_p(\sqrt n)$, since the drift coefficient is $\le0$ and $n-\ell_{n,k}\ge0$.
For the a.s.\ bound, Theorem~\ref{thm:LLN} gives $n(\hat{\rho}_n-v)=O(\sqrt{n\log\log n})$,
hence $\ell_{n,k}(\hat{\rho}_{\ell_{n,k}}-\hat{\rho}_{n,k})=O(\sqrt{n\log\log n})$,
and the LIL for the martingale $M$ gives $M_{n,k}=O(\sqrt{n\log\log n})$; the same
substitution yields the second inequality.
\end{proof}

\begin{lemma}[Deviation between proportions and plug-in target]\label{lem:prop_dev}
\lean{AlphaRAR.aRTS_prop_dev, AlphaRAR.prop_dev}
\uses{def:generic_cond, def:counts, def:plugin_target}
Under the generic conditions and Conditions~\textbf{A}--\textbf{B}, for every $k$,
\begin{align*}
  |N_{n,k} - n\hat{\rho}_{n,k}| &= o_p(\sqrt n), &
  |N_{n,k} - n\hat{\rho}_{n,k}| &= O(\sqrt{n\log\log n})\ \text{a.s.}, \\
  N_{n,k} - n v_k &= O(\sqrt{n\log\log n})\ \text{a.s.}
\end{align*}
\emph{Formalization.} The generic $o_p(\sqrt n)$ assembly is \texttt{AlphaRAR.prop\_dev}: for a finite
family with $\sum_k \mathrm{Dev}_k = 0$ (\texttt{sum\_count\_sub\_smul\_eq\_zero}), the generic
inequality \eqref{eq:generic_ineq}, and the $U$-increment decomposition (drift $-(1-\alpha)v_k$, the
$M$- and $\rho$-terms increment-controlled by \texttt{qm\_increments\_of\_bdd} and
\texttt{ell\_rho\_control}, and an $o_p(1)$ perturbation from \texttt{diff\_U\_decomp}), it yields
$\mathrm{Dev}_k=o_p(\sqrt n)$ via the drift-sign lemma
\texttt{isLittleOpOne\_max\_div\_sqrt\_of\_drift} and the simplex reverse step
\texttt{isLittleOpOne\_dev\_of\_sum\_zero}. The concrete paper-form statement
$|N_{n,k}-n\hat\rho_{n,k}|=o_p(\sqrt n)$ is \texttt{AlphaRAR.aRTS\_prop\_dev}, the full, self-contained
$\aRTS$ instantiation: from the $\aRTS$ design bundle (an \texttt{IsAlgEnvSeq} sequence, Condition
\textbf{A} $Y\in L^2$, a continuous simplex-valued target $\Target$, the algorithm-level predicate
\texttt{IsARTS} --- whence the throttle via \texttt{throttle\_of\_isARTS} ---, $\alpha\in[0,1)$,
Condition \textbf{B} ($\Target$ Lipschitz near $\Params$ and non-sparsity $0<v_k$), and the
$\mathrm{thm{:}LLN}$ consistencies $N_{n,k}/n\to v_k$ and $\estParams_n\to\Params$), it discharges
\emph{every} hypothesis of \texttt{prop\_dev} internally: the generic inequality via
\texttt{generic\_ineq\_of\_hitting} (a.e., from the throttle), the smallness via
\texttt{preliminary\_small}, the $M$-increment via the seam $M=\texttt{assignMart}$
(\texttt{assignMart\_eq\_assignMG}) and \texttt{qm\_increments\_of\_bdd}, the \texttt{diff\_U\_decomp}
perturbation shown $o_p$ (a.e.\ $\to0$ by $\hat\rho_n\to v$), and the plug-in-target increment control
(Lemma~\ref{lem:ell_rho_control}) via \texttt{ell\_rho\_control}, whose $g$- and $h$-coefficients are
the random-hitting-time bounds \texttt{aRTS\_g\_littleOp} ($O_p\cdot o_p=o_p$, Doob running-max plus
$\sqrt n/(N_n+1)\to0$) and \texttt{aRTS\_h\_bigOp} (a.s.\ bounded). All random-hitting-time terms are
made measurable by \texttt{measurable\_eval\_of\_le}. Only the two a.s.\ $O(\sqrt{n\log\log n})$
statements remain unformalized.
\end{lemma}

\begin{proof}
\uses{lem:U_increment_bound, lem:counts_sum, thm:LLN}
The generic inequality \eqref{eq:generic_ineq} together with
\eqref{eq:generic_small} and Lemma~\ref{lem:U_increment_bound} gives the one-sided
bounds $N_{n,k}-n\hat{\rho}_{n,k} \le o_p(\sqrt n)$ and $\le O(\sqrt{n\log\log n})$
a.s. The reverse inequality follows from $\sum_j N_{n,j}=n$, $\sum_j\hat{\rho}_{n,j}=1$
(Lemma~\ref{lem:counts_sum}): $n\hat{\rho}_{n,k}-N_{n,k}=\sum_{j\ne k}(N_{n,j}-n\hat{\rho}_{n,j})$,
which inherits the same one-sided bounds. Combining gives the absolute-value
statements. The last line adds $n(\hat{\rho}_{n,k}-v_k)=O(\sqrt{n\log\log n})$ from
Theorem~\ref{thm:LLN}.
\end{proof}

\section{Joint central limit theorem for the estimators}
\label{sec:norm_clt}

\begin{lemma}[Componentwise joint CLT]\label{lem:componentwise}
\uses{def:estimator, def:counts, def:condA}
Consider any RAR procedure with $N_{n,k}\to\infty$ a.s.\ for all $k$. Under
Condition~\textbf{A}, with $D_n := \mathrm{diag}(\sqrt{N_{n,1}},\dots,\sqrt{N_{n,K}})$,
\[
  D_n(\estParams_n - \Params) \eqdist \mathcal{N}\!\big(0, \mathrm{diag}(V_1,\dots,V_K)\big).
\]
\end{lemma}

\begin{proof}
\uses{lem:Q_martingale, lem:Q_quad_var, thm:mart_clt, cor:mart_clt_componentwise}
The vector martingale $M_{n,k}=\sum_{i\le n} X_{i,k}Y_{i,k}$ (with $Y_{i,k}=\xi_{i,k}-\theta_k$)
has quadratic variation $\langle M_{\cdot,k}\rangle_n = V_k N_{n,k}$ and zero
cross-variation (Lemma~\ref{lem:Q_quad_var}). The Lindeberg condition holds: for
$\varepsilon>0$, $L_{n,k}(\varepsilon)\le \frac1{V_k}\E[Y_{1,k}^2\indic_{\{|Y_{1,k}|>\varepsilon\sqrt{V_k}\sqrt m\}}]$
on $\{N_{n,k}\ge m\}$, which is small for $m$ large by dominated convergence
(using $\E|Y_{1,k}|^2<\infty$), and $\Prob(N_{n,k}<m)\to0$. The martingale CLT
gives $M_n/\sqrt{\langle M\rangle_n}\eqdist \mathcal{N}(0,I_K)$. Finally
$D_n(\estParams_n-\Params) = \mathrm{diag}(\sqrt{V_1},\dots,\sqrt{V_K})\,M_n/\sqrt{\langle M\rangle_n}$
(up to a Bahadur $o_p(1)$ term), giving the stated limit.
\end{proof}

\textbf{Lean remark.} Lemma~\ref{lem:componentwise} uses only $N_{n,k}\to\infty$,
not any property of the limiting proportions. It is the crux of both the
non-sparse and the sparse (Chapter~\ref{chap:forced}) analyses and should be a
standalone theorem depending on a martingale-CLT lemma from Mathlib.

\section{The normality theorem}
\label{sec:norm_main}

\begin{lemma}[CLT for the estimator]\label{lem:clt_theta}
\lean{AlphaRAR.estimator_sqrtN_joint_tendsto_multivariateGaussian}\leanok
\uses{def:estimator, lem:mart_clt_joint_selfnorm, lem:estimator_bahadur}
Suppose $N_{n,k}/n\to v_k$ a.s.\ with $v_k>0$ for every arm $k$. Then, under
Condition~\textbf{A}, $\sqrt n(\estParams_n - \Params)\eqdist\mathcal{N}(0,V)$ with
$V=\mathrm{diag}(V_1/v_1,\dots,V_K/v_K)$.
\end{lemma}

\begin{proof}\leanok
\uses{lem:mart_clt_joint_selfnorm, lem:estimator_bahadur}
Exactly as in Corollary~\ref{cor:mart_clt_componentwise}, but with the deterministic
normalization $\sqrt n$ in place of $D_n=\mathrm{diag}(\sqrt{N_{n,k}})$. By the Bahadur
identity (Lemma~\ref{lem:estimator_bahadur}), coordinate $k$ of
$\sqrt n(\estParams_n-\Params)$ is
$(Q_{n,k}/\sqrt{N_{n,k}})\cdot\sqrt{nN_{n,k}}/(N_{n,k}+1)
 + (\theta_{0,k}-\theta_k)\sqrt n/(N_{n,k}+1)$. The self-normalized martingale vector
$(Q_{n,k}/\sqrt{N_{n,k}})_k$ converges to $\mathcal N(0,\mathrm{diag}(V_k))$
(Lemma~\ref{lem:mart_clt_joint_selfnorm}); the multiplicative factor tends to $1/\sqrt{v_k}$ and
the additive one to $0$ in probability. Multivariate Slutsky then gives the diagonal rescaling
$\mathrm{diag}(1/\sqrt{v_k})$ of the Gaussian, whose covariance is $\mathrm{diag}(V_k/v_k)=V$.
\end{proof}

\begin{lemma}[CLT for the plug-in target]\label{lem:clt_rho}
\lean{AlphaRAR.clt_rho}\leanok
\uses{def:plugin_target}
Under the standing hypotheses, $\sqrt n(\hat{\rho}_n - v)\eqdist\mathcal{N}(0,GVG^\top)$.
\end{lemma}

\begin{proof}\leanok
\uses{lem:clt_theta, lem:rho_rate, lem:taylor_rho}
By the Taylor expansion (as in Lemma~\ref{lem:ell_rho_control}),
$\sqrt n(\hat{\rho}_n-v)=G\sqrt n(\estParams_n-\Params)+O_p(1/\sqrt n)$. Apply
Lemma~\ref{lem:clt_theta} and Slutsky.

\emph{Formalization.} The delta method is assembled in
\texttt{AlphaRAR.clt\_rho}. First, the product-space Slutsky lemma
(\texttt{tendsto\_map\_comp\_of\_tendstoInMeasure\_const}) with $g(x,r)=Gx+r$ combines the estimator
CLT (Lemma~\ref{lem:clt_theta}) with the fact that the first-order remainder
$\sqrt n(\Target(\estParams_n)-\Target(\Params))-G\sqrt n(\estParams_n-\Params)=\sqrt n\,\varphi(\estParams_n-\Params)$
tends to $0$ in probability, and the limiting law is the linear pushforward of the estimator's Gaussian
(\texttt{AlphaRAR.multivariateGaussian\_map\_matrix}, giving the covariance $GVG^\top$). The
remainder-in-probability is discharged by \texttt{MeasureTheory.tendstoInMeasure\_smul\_littleO\_of\_tight}
(an abstract $\varepsilon$--$\delta$ lemma: if $X_n=a_n\cdot S_n$ is tight, $S_n\to0$ in probability and
$\varphi=o(\|\cdot\|)$ at the origin, then $a_n\cdot\varphi(S_n)\to0$ in probability), whose tightness
input is supplied from the estimator CLT via
\texttt{MeasureTheory.tight\_of\_tendsto\_probabilityMeasure} (portmanteau plus continuity from above of
the limit Gaussian over $\{\|x\|\ge j\}\downarrow\emptyset$). The differentiability of $\Target$ at
$\Params$ (Jacobian $G$; Condition~\textbf{B}) and the a.s.\ consistency $\estParams_n\to\Params$
(Theorem~\ref{thm:LLN}) enter as the hypotheses \texttt{hTderiv} and \texttt{hcons}.
\end{proof}

\begin{theorem}[Asymptotic normality]\label{thm:normality}
\uses{def:generic_cond, def:condA, def:condB, def:counts, def:plugin_target, def:estimator}
Under the generic conditions and Conditions~\textbf{A}--\textbf{B}, as
$n\to\infty$:
\begin{enumerate}
  \item[(i)] for every $k$, $|N_{n,k}-n\hat{\rho}_{n,k}|=o_p(\sqrt n)$,
        $|N_{n,k}-n\hat{\rho}_{n,k}|=O(\sqrt{n\log\log n})$ a.s., and
        $N_{n,k}-nv_k=O(\sqrt{n\log\log n})$ a.s.;
  \item[(ii)] $\sqrt n(\estParams_n-\Params)\eqdist\mathcal{N}(0,V)$ and
        \[
          \begin{pmatrix}\sqrt n(N_n/n - v)\\ \sqrt n(\hat{\rho}_n - v)\end{pmatrix}
          \eqdist \mathcal{N}(0,\Omega).
        \]
\end{enumerate}
\end{theorem}

\begin{proof}
\uses{lem:prop_dev, lem:clt_theta, lem:clt_rho}
Part (i) is Lemma~\ref{lem:prop_dev}. For (ii), the marginal CLT for
$\estParams_n$ is Lemma~\ref{lem:clt_theta}. By part (i),
$\sqrt n(N_n/n-\hat{\rho}_n)=o_p(1)$, so $\sqrt n(N_n/n-v)$ and
$\sqrt n(\hat{\rho}_n-v)$ have the same limit $\mathcal{N}(0,GVG^\top)$
(Lemma~\ref{lem:clt_rho}) and are asymptotically equal; the joint vector
$(\sqrt n(N_n/n-v),\sqrt n(\hat{\rho}_n-v))$ therefore equals
$(\sqrt n(\hat{\rho}_n-v),\sqrt n(\hat{\rho}_n-v))+o_p(1)$, which converges to
$\mathcal{N}(0,\Omega)$ by Slutsky.

\emph{Formalization.} The joint statement of part (ii) is
\texttt{AlphaRAR.clt\_joint}: it takes the proportion-deviation fact
$\sqrt n(N_n/n-\hat{\rho}_n)\to_p 0$ of part (i) as a hypothesis and assembles the joint CLT from
\texttt{AlphaRAR.clt\_rho}. The joint vector is the image of $\sqrt n(\hat{\rho}_n-v)$ under the
duplication map $x\mapsto(x,x)$ (the stacked matrix $[I;I]$) up to an $o_p(1)$ remainder, so the
same product-space Slutsky lemma as in Lemma~\ref{lem:clt_rho} applies, and the limiting covariance
$[I;I]\,GVG^\top\,[I;I]^\top$ is the block matrix $\Omega$ via the rectangular Gaussian pushforward
\texttt{AlphaRAR.multivariateGaussian\_map\_matrix}. Part (i)
(Lemma~\ref{lem:prop_dev}, the $o_p(\sqrt n)$ deviation machinery) is not yet formalized, so the
theorem is not marked complete.
\end{proof}

\begin{remark}[Asymptotic efficiency]\label{rem:efficiency}
\uses{thm:normality, prop:exp_fam}
When the arm responses belong to one-parameter exponential families with mean
parameter, Proposition~\ref{prop:exp_fam} gives $V_k = I_k(\theta_k)^{-1}$, so the
asymptotic variance $GVG^\top$ of $N_n/n$ equals the lower bound
$G\,I(\Params)^{-1}G^\top$ of \cite{hu2006asymptotically}, with
$I(\Params)=\mathrm{diag}(v_1 I_1(\theta_1),\dots,v_K I_K(\theta_K))$. Hence every
$\aRTS$ design is asymptotically efficient.
\end{remark}

\chapter{Forced exploration and sparse targets}
\label{chap:forced}

This chapter augments the $\aRTS$ family with a forced-exploration mechanism,
yielding the $\aRTSFE$ family. Forced exploration guarantees that every arm is
sampled infinitely often, which (a) lets the previous asymptotic theory go
through unchanged for non-sparse targets, and (b) extends consistency and a
componentwise CLT to \emph{sparse} targets. It corresponds to Section~4 and
Appendix~B of the paper.

\section{The \texorpdfstring{$\aRTSFE$}{aRTS-FE} family}
\label{sec:forced_def}

\begin{definition}[Exploration schedule]\label{def:exploration_schedule}
An \emph{exploration schedule} is a function $h : \N \to \R$ with
\begin{enumerate}
  \item[(i)] $\lim_{n\to\infty} h(n) = +\infty$;
  \item[(ii)] $h(n) = o(\sqrt n)$ as $n\to\infty$.
\end{enumerate}
\end{definition}

\begin{example}\label{ex:h_choice}
The D-Tracking choice $h(n) = (\sqrt n - K/2)^+$ of \cite{garivier2016optimal}
violates condition (ii). The choice $h(n) = (n^{1/3} - K/2)^+$ satisfies both
conditions and is used in the experiments.
\end{example}

\begin{definition}[Under-sampled and least-sampled sets]\label{def:underexplored_sets}
\uses{def:counts, def:exploration_schedule}
For a schedule $h$ and $m\ge1$,
\[
  \mathbf{U}_m := \{ j\in[K] : N_{m,j} \le h(m)\}, \qquad
  \mathbf{S}_m := \argmin_{j\in\mathbf{U}_m} N_{m,j}.
\]
\end{definition}

\begin{definition}[$\aRTSFE$ family]\label{def:aRTSFE}
\uses{def:aRTS, def:underexplored_sets}
Fix $\alpha\in[0,1)$ and a schedule $h$. A RAR procedure is in the $\aRTSFE$
family if its allocation probabilities obey:
\begin{align}
  \textbf{(Forced exploration)}\quad
    &\mathbf{U}_m\ne\emptyset \Rightarrow \forall a\in\mathbf{S}_m:\ p_{m+1,a}=\tfrac{1}{|\mathbf{S}_m|};
    \label{eq:fe_phase}\\
  \textbf{(Otherwise)}\quad
    &\mathbf{U}_m=\emptyset \Rightarrow p_{m+1}\ \text{satisfies the throttling condition \eqref{eq:throttle}.}
    \label{eq:fe_otherwise}
\end{align}
\end{definition}

\begin{remark}\label{rem:d_tracking_fe}
The full D-Tracking algorithm of \cite{garivier2016optimal} is the $\aRTSFE$
member with $\alpha=0$ (up to the choice of $h$).
\end{remark}

\section{The forced-exploration hitting time}
\label{sec:forced_hitting}

\begin{definition}[$\aRTSFE$ last hitting time]\label{def:hitting_fe}
\uses{def:hitting, def:underexplored_sets}
For $k\in[K]$ and $n\ge1$,
\[
  \ell_{n,k}^{\mathrm{FE}} := \max\Big\{ m\le n :
    \tfrac{N_{m,k}}{m}\le\hat{\rho}_{m,k}
    \ \text{or}\ \big(\mathbf{U}_m\ne\emptyset\ \text{and}\ k\in\argmin_{j\in\mathbf{U}_m}N_{m,j}\big)\Big\}.
\]
It is the last time before $n$ at which arm $k$ is under-sampled or is a
least-sampled arm among the under-explored set.
\end{definition}

\begin{lemma}[Sign after the $\aRTSFE$ hitting time]\label{lem:hitting_fe_sign}
\uses{def:hitting_fe, def:aRTSFE}
For an $\aRTSFE$ design and $m$ with $\ell_{n,k}^{\mathrm{FE}} < m \le n$,
we have $\frac{N_{m,k}}{m} > \hat{\rho}_{m,k}$ and
$\big(\mathbf{U}_m=\emptyset\ \text{or}\ k\notin\argmin_{j\in\mathbf{U}_m}N_{m,j}\big)$;
consequently $p_{m+1,k}\le\alpha\hat{\rho}_{m,k}$ (equal to $0$ during a forced
exploration step in which $k$ is not least-sampled).
\end{lemma}

\begin{proof}
By maximality of $\ell_{n,k}^{\mathrm{FE}}$, neither disjunct in
Definition~\ref{def:hitting_fe} holds at such $m$. If $\mathbf{U}_m\ne\emptyset$
then $k\notin\argmin_{j\in\mathbf{U}_m}N_{m,j}$, so \eqref{eq:fe_phase} gives
$p_{m+1,k}=0\le\alpha\hat{\rho}_{m,k}$. If $\mathbf{U}_m=\emptyset$ then
\eqref{eq:fe_otherwise} and the over-sampling $N_{m,k}/m>\hat{\rho}_{m,k}$ give
$p_{m+1,k}\le\alpha\hat{\rho}_{m,k}$.
\end{proof}

\begin{lemma}[$\aRTSFE$ satisfies the generic conditions]\label{lem:preliminary_fe}
\uses{def:generic_cond, def:aRTSFE, def:hitting_fe}
For any $\aRTSFE$ design, the generic conditions
(Definition~\ref{def:generic_cond}) hold with $\ell_{n,k}^{\mathrm{FE}}$.
\end{lemma}

\begin{proof}
\uses{lem:hitting_fe_sign, lem:U_telescope, def:exploration_schedule}
(i) Monotonicity and $\ell_{n,k}^{\mathrm{FE}}\le n$ are immediate. (ii) The
telescoping identity (Lemma~\ref{lem:U_telescope}, whose proof only uses the
increment of $U$) combined with Lemma~\ref{lem:hitting_fe_sign} — every
intermediate $p_{m,k}-\alpha\hat{\rho}_{m-1,k}\le0$ — gives inequality
\eqref{eq:generic_ineq}. (iii) At $m=\ell_{n,k}^{\mathrm{FE}}$, either
$N_{m,k}/m\le\hat{\rho}_{m,k}$, giving $N_{m,k}-m\hat{\rho}_{m,k}\le0$, or $k$ is
least-sampled in $\mathbf{U}_m$, giving
$N_{\ell_{n,k}^{\mathrm{FE}},k}-\ell_{n,k}^{\mathrm{FE}}\hat{\rho}_{\ell_{n,k}^{\mathrm{FE}},k}
 \le N_{\ell_{n,k}^{\mathrm{FE}},k}\le h(\ell_{n,k}^{\mathrm{FE}})=o(\sqrt n)$
by Definition~\ref{def:exploration_schedule}.
\end{proof}

\section{Validity under forced exploration (non-sparse targets)}
\label{sec:forced_nonsparse}

\begin{theorem}[Validity under forced exploration]\label{thm:forced_valid}
\uses{def:aRTSFE, def:condA, def:condB}
Under Conditions~\textbf{A}--\textbf{B}, Theorem~\ref{thm:LLN} and
Theorem~\ref{thm:normality} hold for every design in the $\aRTSFE$ family.
\end{theorem}

\begin{proof}
\uses{lem:preliminary_fe, thm:generic_main}
By Lemma~\ref{lem:preliminary_fe} the $\aRTSFE$ design satisfies the generic
conditions with $\ell^{\mathrm{FE}}$; apply Theorem~\ref{thm:generic_main}.
\end{proof}

\section{Sparse target allocations}
\label{sec:forced_sparse}

We now drop Condition~\textbf{B} and allow $v_k=0$ for some $k$. The forced
exploration guarantees infinite sampling of all arms, which is the only
ingredient the componentwise CLT needs.

\begin{lemma}[Forced exploration prevents starvation]\label{lem:fe_no_starvation}
\uses{def:aRTSFE, def:exploration_schedule}
For any $\aRTSFE$ design, $N_{n,k}\to\infty$ a.s.\ for every $k\in[K]$ (no
assumption on the target).
\end{lemma}

\begin{proof}
Suppose for contradiction that some arm $i$ stalls: $N_{n,i}=M$ for all $n\ge N$.
Since $h(n)\to\infty$, enlarge $N$ so that $h(n)\ge M$ for $n\ge N$. Then for
$n\ge N$, $N_{n,i}\le h(n)$ gives $i\in\mathbf{U}_n\ne\emptyset$, so the arm
$I_{n+1}$ selected by forced exploration is a least-sampled arm of $\mathbf{U}_n$,
whence $N_{n,I_{n+1}}\le N_{n,i}=M$ and $N_{n+1,I_{n+1}}\le M+1$. The potential
$D(n):=\sum_j (M+1-N_{n,j})^+$ then satisfies $D(n+1)-D(n)=-1$ for all $n\ge N$,
so $D$ becomes negative, contradicting $D\ge0$. Hence every arm is sampled
infinitely often.
\end{proof}

\begin{theorem}[Consistency and rate for sparse targets]\label{thm:sparse_rate}
\uses{def:aRTSFE, def:condA, def:counts, def:estimator}
Under Condition~\textbf{A} only, any $\aRTSFE$ design satisfies, for every $k$,
\[
  N_{n,k}\xrightarrow{a.s.}+\infty,\qquad
  \frac{N_{n,k}}{n}\xrightarrow{a.s.}v_k,\qquad
  \estParams_n-\Params = O\!\Big(\sqrt{\tfrac{\log\log N_{n,k}}{N_{n,k}}}\Big)\ \text{a.s.}
\]
\end{theorem}

\begin{proof}
\uses{lem:fe_no_starvation, lem:theta_limit_dichotomy, lem:match, lem:convergence, lem:preliminary_fe, lem:lil_truncation}
The first statement is Lemma~\ref{lem:fe_no_starvation}. Given it, every arm's
estimator converges: $\hat{\theta}_{n,k}\to\theta_k$ by the SLLN, so
$\hat{\rho}_n\to v$ by continuity of $\Target$; the matching argument (as in
Lemma~\ref{lem:match}, using Lemma~\ref{lem:convergence} and the generic
inequality valid by Lemma~\ref{lem:preliminary_fe}) gives $N_{n,k}/n\to v_k$,
allowing $v_k=0$. The rate is the LIL for the martingale $Q_{\cdot,k}$ scaled by
its own sample size $N_{n,k}$, as in Lemma~\ref{lem:theta_LIL} but without
dividing by $n$.
\end{proof}

\begin{corollary}[Componentwise CLT for sparse targets]\label{cor:sparse_clt}
\uses{def:aRTSFE, def:condA, def:estimator}
Under Condition~\textbf{A} only, any $\aRTSFE$ design satisfies
\[
  D_n(\estParams_n-\Params)\eqdist\mathcal{N}\!\big(0,\mathrm{diag}(V_1,\dots,V_K)\big),
  \qquad D_n=\mathrm{diag}(\sqrt{N_{n,1}},\dots,\sqrt{N_{n,K}}).
\]
\end{corollary}

\begin{proof}
\uses{lem:fe_no_starvation, lem:componentwise}
By Lemma~\ref{lem:fe_no_starvation}, $N_{n,k}\to\infty$ for all $k$, which is the
sole hypothesis of Lemma~\ref{lem:componentwise}. Apply it.
\end{proof}

\begin{remark}\label{rem:sparse_normalization}
The normalization in Corollary~\ref{cor:sparse_clt} is componentwise (by
$\sqrt{N_{n,k}}$) rather than by $\sqrt n$, which is what makes it meaningful when
$v_k=0$. When $v\in(0,1)^K$ it recovers the classical result of
\cite{hu2006asymptotically}.
\end{remark}

\chapter{Application: hypothesis testing}
\label{chap:testing}

As an application of the asymptotic theory, we calibrate a Pearson chi-square
homogeneity test on data collected by an $\aRTS$ or $\aRTSFE$ design. This
corresponds to Section~5.3 and Appendix~B.3 of the paper.

Consider $K$ Bernoulli arms and the homogeneity test
\[
  \mathcal{H}_0 : \forall k,\ \theta_k = \theta_1
  \qquad\text{against}\qquad
  \mathcal{H}_1 : \exists i,j,\ \theta_i \ne \theta_j.
\]

\begin{definition}[Pearson statistic]\label{def:pearson_stat}
\uses{def:estimator, def:counts}
Let $S_{n,k} := N_{n,k}\hat{\theta}_{n,k}$ and
$\tilde{\theta}_n := \sum_{k=1}^K \frac{N_{n,k}}{n}\hat{\theta}_{n,k}$ be the
pooled estimator. The \emph{Pearson chi-square statistic} is
\[
  \mathcal{X}_n := \sum_{k=1}^K \frac{(S_{n,k} - N_{n,k}\tilde{\theta}_n)^2}
                                        {N_{n,k}\tilde{\theta}_n(1-\tilde{\theta}_n)}.
\]
\end{definition}

\begin{lemma}[Quadratic-form rewriting]\label{lem:pearson_quadform}
\uses{def:pearson_stat}
With $Y_n := (\sqrt{N_{n,k}}(\hat{\theta}_{n,k}-\theta))_{k\in[K]}$ and
$s_n := (\sqrt{\omega_{n,k}})_{k\in[K]}$ where $\omega_{n,k}:=N_{n,k}/n$, under
$\mathcal{H}_0$ (common value $\theta$),
\[
  \mathcal{X}_n = \frac{1}{\tilde{\theta}_n(1-\tilde{\theta}_n)}\,
                  Y_n^\top\,(I - s_n s_n^\top)\,Y_n.
\]
\end{lemma}

\begin{proof}
Under $\mathcal{H}_0$, $\sqrt{N_{n,k}}(\tilde{\theta}_n-\theta)
 = \sqrt{\omega_{n,k}}\sum_j\sqrt{\omega_{n,j}}[Y_n]_j$, so
$\sqrt{N_{n,k}}(\hat{\theta}_{n,k}-\tilde{\theta}_n) = [Y_n - s_n s_n^\top Y_n]_k$.
Substituting $S_{n,k}-N_{n,k}\tilde{\theta}_n = N_{n,k}(\hat{\theta}_{n,k}-\tilde{\theta}_n)$
into Definition~\ref{def:pearson_stat} gives the quadratic form.
\end{proof}

\begin{lemma}[Idempotent projection]\label{lem:projection}
For a unit vector $s\in\R^K$ (here $s=(\sqrt{v_1},\dots,\sqrt{v_K})$ with
$\sum_k v_k=1$), the matrix $I - s s^\top$ is an orthogonal projection of rank
$K-1$.
\end{lemma}

\begin{proof}
$s^\top s = \sum_k v_k = 1$, so $(I-ss^\top)^2 = I - 2ss^\top + s(s^\top s)s^\top
 = I - ss^\top$, and it is symmetric; its trace $K-1$ equals its rank.
\end{proof}

\begin{lemma}[Asymptotic null distribution of the Pearson statistic]\label{lem:pearson}
\uses{def:pearson_stat}
Consider a RAR procedure with $N_{n,k}\to\infty$ a.s.\ for all $k$ (in
particular any $\aRTS$ design under Condition~\textbf{B}, or any $\aRTSFE$
design). Under $\mathcal{H}_0$,
\[
  \mathcal{X}_n \eqdist \chi^2_{K-1}.
\]
\end{lemma}

\begin{proof}
\uses{lem:pearson_quadform, lem:projection, lem:componentwise, thm:LLN, thm:cochran}
By Lemma~\ref{lem:componentwise}, $\frac{1}{\sqrt{\theta(1-\theta)}}Y_n
\eqdist\mathcal{N}(0,I_K)$ (Bernoulli variance $V_k=\theta(1-\theta)$ under
$\mathcal{H}_0$). By Theorem~\ref{thm:LLN}, $\tilde{\theta}_n\to\theta$ and
$s_n\to s=(\sqrt{v_1},\dots,\sqrt{v_K})$, so $I-s_ns_n^\top\to I-ss^\top$, an
orthogonal projection of rank $K-1$ by Lemma~\ref{lem:projection}. Slutsky's
lemma and Lemma~\ref{lem:pearson_quadform} give
$\mathcal{X}_n\eqdist Z^\top(I-ss^\top)Z$ with $Z\sim\mathcal{N}(0,I_K)$, which is
$\chi^2_{K-1}$.
\end{proof}

\begin{remark}\label{rem:test_level}
Lemma~\ref{lem:pearson} shows that rejecting $\mathcal{H}_0$ when $\mathcal{X}_n$
exceeds the $0.95$-quantile of $\chi^2_{K-1}$ is asymptotically a level-$0.05$
test, valid for all $\aRTS$ designs (under Condition~\textbf{B}) and all
$\aRTSFE$ designs.
\end{remark}


\part{Probabilistic prerequisites to be added to Mathlib}
\label{part:prereq}

\textbf{Overview}

The proofs in Part~I rely on several standard results of probability theory that
are not yet available in Mathlib. This part develops them at the level of a
blueprint: statements decomposed into small lemmas, with proof sketches and, for
each, a note on the Mathlib building blocks already available. Nothing here is
formalized yet, and the results are stated for their own sake (independently of
the $\aRTS$ application), so that each can become a self-contained Mathlib
contribution.

Ordered by how much of Part~I depends on them:
\begin{itemize}
  \item Chapter~\ref{chap:pre_orders}: supporting API (stochastic orders
        $o_p/O_p$, predictable quadratic variation, a multivariate Taylor
        remainder). Not deep, but pervasive.
  \item Chapter~\ref{chap:pre_slln}: the strong law of large numbers for
        martingales, used for the consistency of the allocation.
  \item Chapter~\ref{chap:pre_lil}: a one-sided law of the iterated logarithm for
        martingales, used for every almost-sure rate.
  \item Chapter~\ref{chap:pre_llil}: the general martingale law of the iterated
        logarithm at the sharp $\log\log$ rate, sharpening
        Chapter~\ref{chap:pre_lil}, with the i.i.d.\ Hartman--Wintner LIL and
        optional skipping as applications.
  \item Chapter~\ref{chap:pre_clt}: the martingale central limit theorem, the
        analytic core of the asymptotic normality results.
  \item Chapter~\ref{chap:pre_cochran}: the $\chi^2$ distribution and Cochran's
        theorem, used to calibrate the Pearson test.
\end{itemize}
The variance/inverse-Fisher-information identity for exponential families
(Proposition~\ref{prop:exp_fam}) is deliberately omitted here: it is used only in
the efficiency remark and can be discharged concretely per distribution family.

\chapter{Stochastic orders, quadratic variation, and Taylor}
\label{chap:pre_orders}

This chapter collects the supporting API that the asymptotic proofs use but that
Mathlib does not yet package: a calculus for the stochastic Landau symbols
$o_p, O_p$; the predictable quadratic variation of a discrete-time martingale;
and a multivariate first-order Taylor expansion with a quadratic remainder. None
of these is a deep theorem, but each is used pervasively.

\section{The \texorpdfstring{$o_p$}{op} / \texorpdfstring{$O_p$}{Op} calculus}
\label{sec:pre_op}

Recall Definition~\ref{def:op_Op}. We reduce everything to statements about
$o_p(1)$ (convergence to $0$ in probability) and $O_p(1)$ (boundedness in
probability, i.e.\ tightness).

\textbf{Formalization note.} Mathlib provides the primitives: convergence in
probability is \texttt{MeasureTheory.TendstoInMeasure} to a constant
(\texttt{MeasureTheory/Function/ConvergenceInMeasure.lean}), and boundedness in
probability is tightness of the laws. The lemmas below are an API layer on top of
these, plus the bridge to Slutsky's theorem
(\texttt{ConvergenceInDistribution.lean}). This is engineering rather than new
mathematics, but it must exist before any of the rate computations can be
written.

\begin{lemma}[Reduction to order one]\label{lem:op_factor}
\uses{def:op_Op}
$X_n = o_p(a_n)$ if and only if $X_n/a_n = o_p(1)$, and $X_n = O_p(a_n)$ if and
only if $X_n/a_n = O_p(1)$.
\end{lemma}

\begin{proof}
Unfold Definition~\ref{def:op_Op}: both sides are the same statement about the
sequence $|X_n|/a_n$.
\end{proof}

\begin{lemma}[Convergence in probability gives $o_p$]\label{lem:op_of_tendsto}
\lean{AlphaRAR.isLittleOpOne_of_tendsto_ae}
\uses{def:op_Op}
If $X_n \to 0$ in probability (in particular, if $X_n \to 0$ a.s.), then
$X_n = o_p(1)$. If $X_n \to c$ in distribution for some random variable, then
$X_n = O_p(1)$.
\end{lemma}

\begin{proof}
The first statement is the definition of convergence in probability. For the
second, convergence in distribution implies tightness of the laws
($\{X_n\}$ is uniformly tight), which is exactly $O_p(1)$.
\end{proof}

\begin{lemma}[Arithmetic of stochastic orders]\label{lem:op_arith}
\lean{AlphaRAR.IsLittleOpOne.add, AlphaRAR.IsBigOpOne.add, AlphaRAR.IsLittleOpOne.add_bigOpOne, AlphaRAR.IsBigOpOne.mul, AlphaRAR.IsBigOpOne.mul_littleOp, AlphaRAR.IsBigOpOne.bdd_mul}\leanok
\uses{def:op_Op}
For a fixed rate $a_n$:
\begin{enumerate}
  \item[(i)] $o_p(a_n) + o_p(a_n) = o_p(a_n)$ and $O_p(a_n) + O_p(a_n) = O_p(a_n)$;
  \item[(ii)] $o_p(a_n) + O_p(a_n) = O_p(a_n)$;
  \item[(iii)] if $X_n = O_p(1)$ and $Y_n = o_p(1)$ then $X_n Y_n = o_p(1)$;
  \item[(iv)] if $X_n = O_p(a_n)$ and $c_n \to c \in \R$ then $c_n X_n = O_p(a_n)$.
\end{enumerate}
\end{lemma}

\begin{proof}\leanok
\uses{lem:op_factor}
By Lemma~\ref{lem:op_factor} reduce to $a_n = 1$. Then (i)--(ii) are union bounds
on the events $\{|X_n| > \varepsilon\}$, $\{|Y_n| > \varepsilon\}$; (iii) follows
by conditioning on the tightness of $X_n$: for $\eta$ pick $M$ with
$\sup_n \Prob(|X_n|>M)\le\eta$, then $\Prob(|X_nY_n|>\varepsilon)\le\eta
+\Prob(|Y_n|>\varepsilon/M)$; (iv) is immediate from boundedness of $c_n$.
\end{proof}

\begin{lemma}[Slutsky bridge]\label{lem:op_slutsky}
\uses{def:op_Op}
If $X_n \eqdist X$ (convergence in distribution) and $R_n = o_p(1)$, then
$X_n + R_n \eqdist X$. More generally, if $(X_n) \eqdist X$ and $Y_n \to c$ in
probability for a constant $c$, then $(X_n, Y_n) \eqdist (X, c)$.
\end{lemma}

\begin{proof}
\uses{lem:op_of_tendsto}
This is Slutsky's theorem, available in Mathlib as
\texttt{TendstoInDistribution.prodMk\_of\_tendstoInMeasure\_const}, combined with
the continuous mapping theorem \texttt{TendstoInDistribution.continuous\_comp}
applied to addition.
\end{proof}

\section{Predictable quadratic variation}
\label{sec:pre_qv}

\begin{definition}[Predictable quadratic variation]\label{def:pred_qv}
\lean{AlphaRAR.predQuadVar}\leanok
Let $(M_n)_{n\ge0}$ be a square-integrable $(\filt_n)$-martingale with $M_0 = 0$.
Its \emph{predictable quadratic variation} is
$\langle M\rangle := \mathrm{predictablePart}(M^2)$, the predictable part of the
submartingale $M^2$ in the Doob decomposition. Equivalently,
$\langle M\rangle_0 = 0$ and
$\langle M\rangle_n = \sum_{i=1}^n \E[(\Delta M_i)^2 \mid \filt_{i-1}]$.
\end{definition}

\textbf{Formalization note.} Mathlib already has the discrete Doob decomposition
(\texttt{predictablePart} / \texttt{martingalePart} in
\texttt{Probability/Martingale/Centering.lean}); $\langle M\rangle$ is a thin
wrapper. The lemmas below are the basic API and are straightforward from that
decomposition.

\begin{lemma}[Increment of the quadratic variation]\label{lem:qv_incr}
\lean{AlphaRAR.predQuadVar_succ_sub_eq, AlphaRAR.predQuadVar_ae_eq_sum, AlphaRAR.predQuadVar_mono,
AlphaRAR.isStronglyPredictable_predQuadVar}\leanok
\uses{def:pred_qv}
For a square-integrable martingale $M$ with $M_0=0$,
$\langle M\rangle_n - \langle M\rangle_{n-1} = \E[(\Delta M_n)^2\mid\filt_{n-1}]$,
and $\langle M\rangle$ is nondecreasing and predictable.
\end{lemma}

\begin{proof}\leanok
The Doob decomposition of the submartingale $M^2$ has predictable part with
increments $\E[M_n^2 - M_{n-1}^2\mid\filt_{n-1}]
= \E[(\Delta M_n)^2\mid\filt_{n-1}]$, using that $M$ is a martingale so the cross
term $\E[2M_{n-1}\Delta M_n\mid\filt_{n-1}]=0$. Nonnegativity of the increments
gives monotonicity.
\end{proof}

\begin{lemma}[Linear bound on $\langle M\rangle$ for bounded increments]\label{lem:qv_linear_bound}
\lean{AlphaRAR.predQuadVar_le_of_bound}\leanok
\uses{def:pred_qv, lem:qv_incr}
If $|\Delta M_i|\le c$ a.s., then $\langle M\rangle_n \le c^2 n$ a.s.
\end{lemma}

\begin{proof}\leanok
Each increment satisfies $\langle M\rangle_{k+1}-\langle M\rangle_k
= \E[(\Delta M_k)^2\mid\filt_k] \le \E[c^2\mid\filt_k] = c^2$; telescoping from
$\langle M\rangle_0 = 0$ gives $\langle M\rangle_n \le c^2 n$.
\end{proof}

\begin{lemma}[$M^2 - \langle M\rangle$ is a martingale]\label{lem:qv_mart}
\lean{AlphaRAR.martingale_sq_sub_predQuadVar}\leanok
\uses{def:pred_qv}
$M^2 - \langle M\rangle$ is an $(\filt_n)$-martingale.
\end{lemma}

\begin{proof}\leanok
This is the martingale part in the Doob decomposition of $M^2$
(\texttt{martingale\_martingalePart}).
\end{proof}

\begin{lemma}[$\langle M\rangle$ is the unique compensator]\label{lem:qv_unique}
\lean{AlphaRAR.IsPredQuadVar, AlphaRAR.isPredQuadVar_predQuadVar,
AlphaRAR.IsPredQuadVar.indistinguishable_predQuadVar}\leanok
\uses{def:pred_qv, lem:qv_mart}
Call a process $A$ a \emph{compensator} of $M^2$ if $M^2 - A$ is an $(\filt_n)$-martingale, $A$ is
predictable, $A_0 = 0$, and each $A_n$ is integrable. Then $\langle M\rangle$ is a compensator of
$M^2$, and any compensator $A$ of $M^2$ is indistinguishable from it:
$\P[\forall n,\ A_n = \langle M\rangle_n] = 1$. So Definition~\ref{def:pred_qv}, which is a
formula, defines the object the name describes, and does so up to the a.s.\ equality of paths that
is the best any statement about conditional expectations can offer.
\end{lemma}

\begin{proof}\leanok
\uses{lem:qv_mart, lem:qv_incr}
Existence is Lemma~\ref{lem:qv_mart} together with the predictability of the Doob predictable part
(Lemma~\ref{lem:qv_incr}). For uniqueness, $M^2 = (M^2 - A) + A$ exhibits $M^2$ as a martingale
plus a predictable process null at $0$, and the Doob decomposition is unique up to null sets
(\texttt{predictablePart\_add\_ae\_eq}), so $A_n = \langle M\rangle_n$ a.s.\ for each fixed $n$;
the index set is countable, so the countably many exceptional sets union to a null set.
\end{proof}

\begin{lemma}[Expected quadratic variation equals second moment]\label{lem:qv_second_moment}
\lean{AlphaRAR.integral_sq_eq_integral_predQuadVar,
AlphaRAR.integral_sq_eq_sum_integral_increment_sq}\leanok
\uses{def:pred_qv, lem:qv_mart}
For a square-integrable martingale $M$ with $M_0 = 0$ and every $n$,
$\E[M_n^2] = \E[\langle M\rangle_n]$.
\end{lemma}

\begin{proof}\leanok
$M^2 - \langle M\rangle$ is a martingale (Lemma~\ref{lem:qv_mart}) starting at
$M_0^2 - \langle M\rangle_0 = 0$, so its expectation is constant equal to $0$;
that is $\E[M_n^2] - \E[\langle M\rangle_n] = 0$.
\end{proof}

\begin{lemma}[Martingale $L^2$ growth]\label{lem:mart_sq_growth}
\lean{AlphaRAR.integral_sq_le_of_increment_bound}\leanok
\uses{def:pred_qv, lem:qv_second_moment, lem:qv_incr}
Let $M$ be a square-integrable martingale with $M_0 = 0$. If every increment has
second moment $\E[(\Delta M_{n+1})^2] \le \sigma^2$, then $\E[M_n^2] \le \sigma^2 n$.
\end{lemma}

\begin{proof}\leanok
By Lemma~\ref{lem:qv_second_moment}, $\E[M_n^2] = \E[\langle M\rangle_n]$.
Telescoping and Lemma~\ref{lem:qv_incr},
$\E[\langle M\rangle_n] = \sum_{i=0}^{n-1}\E[(\Delta M_{i+1})^2] \le n\sigma^2$.
\end{proof}

\begin{lemma}[Orthogonal martingales]\label{lem:qv_orthogonal}
\lean{AlphaRAR.martingale_mul, AlphaRAR.predQuadVar_add_of_martingale_mul}\leanok
\uses{def:pred_qv}
Let $M, N$ be square-integrable martingales such that
$\E[\Delta M_i\, \Delta N_i \mid \filt_{i-1}] = 0$ for all $i$ (e.g.\ if
$\Delta M_i\, \Delta N_i = 0$ a.s.). Then $MN$ is a martingale, so the cross
variation vanishes, and $\langle M + N\rangle = \langle M\rangle + \langle N\rangle$.
\end{lemma}

\begin{proof}\leanok
$\E[\Delta(MN)_i\mid\filt_{i-1}] = M_{i-1}\E[\Delta N_i\mid\cdot]
+ N_{i-1}\E[\Delta M_i\mid\cdot] + \E[\Delta M_i\Delta N_i\mid\cdot] = 0$, so $MN$
is a martingale (\texttt{martingale\_mul}). For additivity, expand
$(M+N)^2 = M^2 + N^2 + 2MN$; the predictable part is linear and the predictable
part of the martingale $MN$ vanishes, so
$\langle M+N\rangle = \langle M\rangle + \langle N\rangle$.
\end{proof}

\section{A multivariate Taylor remainder}
\label{sec:pre_taylor}

\begin{lemma}[First-order expansion with quadratic remainder]\label{lem:taylor_first_order}
\lean{AlphaRAR.norm_sub_fderiv_le_mul_sq}\leanok
Let $f : \R^K \to \R^d$ be $C^2$ on a convex open set $\hilb$, and suppose
$\|\mathrm{D}^2 f(z)\| \le C$ for all $z$ in a convex neighborhood of $y \in \hilb$.
Then for $x$ in that neighborhood,
\[
  \big\| f(x) - f(y) - \mathrm{D}f(y)(x - y) \big\| \le \tfrac{C}{2}\,\|x - y\|^2.
\]
\end{lemma}

\begin{proof}\leanok
Apply the mean-value / Taylor formula to $t \mapsto f(y + t(x-y))$ on $[0,1]$:
$f(x) - f(y) - \mathrm{D}f(y)(x-y)
= \int_0^1 (1-t)\, \mathrm{D}^2 f(y + t(x-y))[x-y, x-y]\, dt$, and bound the
integrand by $C\|x-y\|^2$.

\emph{Formalization.} \texttt{AlphaRAR.norm\_sub\_fderiv\_le\_mul\_sq} states the equivalent bound
$\|f(x)-f(y)-\mathrm{D}f(y)(x-y)\| \le K\|x-y\|^2$ under the (weaker, and here more convenient)
hypothesis that the derivative is $K$-Lipschitz relative to $y$ on a convex set,
$\|\mathrm{D}f(z)-\mathrm{D}f(y)\| \le K\|z-y\|$ --- which the second-derivative bound
$\|\mathrm{D}^2 f\|\le C$ supplies with $K = C$ by the mean value inequality. The proof applies the
mean value inequality \texttt{Convex.norm\_image\_sub\_le\_of\_norm\_hasFDerivWithin\_le'} on the
segment $[y,x]$, where $\|\mathrm{D}f(z)-\mathrm{D}f(y)\|\le K\|z-y\|\le K\|x-y\|$ gives the constant.
\end{proof}

\textbf{Formalization note.} Mathlib has the pieces —
\texttt{fderiv}, \texttt{ContDiff}, iterated derivatives, and one-dimensional
Taylor with remainder (\texttt{Analysis/Calculus/Taylor.lean}). What is missing
is this packaged multivariate first-order statement with an explicit quadratic
bound on a neighborhood; it can be assembled from the existing bounds on the
second derivative and the integral form of the remainder.

\begin{corollary}[Expansion of the target function]\label{lem:taylor_rho}
\uses{def:condB, lem:taylor_first_order}
Under Condition~\textbf{B}, there is a constant $C$ and a neighborhood of
$\Params$ on which
$\big\|\Target(x) - \Target(\Params) - G(x - \Params)\big\| \le C\|x - \Params\|^2$,
where $G = \nabla\Target|_\Params$. In particular, if $\hat X \to \Params$ with
$\hat X - \Params = O_p(r_n)$, then
$\Target(\hat X) - \Target(\Params) = G(\hat X - \Params) + O_p(r_n^2)$.
\end{corollary}

\begin{proof}
\uses{lem:op_arith}
Condition~\textbf{B} makes $\Target$ twice differentiable near $\Params$, so
Lemma~\ref{lem:taylor_first_order} applies with $y = \Params$; the second
statement follows from the order arithmetic of Lemma~\ref{lem:op_arith}.
\end{proof}

\chapter{Strong law of large numbers for martingales}
\label{chap:pre_slln}

We formalize the $L^2$ strong law of large numbers for martingales: if the
increment variances decay fast enough (in particular for uniformly bounded
increments), then $M_n / n \to 0$ a.s. This is the result used in
Lemma~\ref{lem:M_lln} to control the assignment martingale.

\textbf{Formalization note.} Mathlib has the i.i.d.\ strong law
(\texttt{Probability/StrongLaw.lean}) but no martingale version. The classical
route below is Kolmogorov's: turn the claim into the a.s.\ convergence of a
series via the $L^2$ martingale convergence theorem (which Mathlib \emph{does}
have, \texttt{Probability/Martingale/Convergence.lean}), then apply Kronecker's
lemma. The only genuinely new pieces are Kronecker's lemma (elementary real
analysis) and the assembly.

\section{Ingredients}
\label{sec:pre_slln_ingredients}

\begin{lemma}[Orthogonality of martingale increments]\label{lem:slln_orthogonal}
Let $(M_n)$ be a square-integrable martingale. For $i \ne j$,
$\E[\Delta M_i\, \Delta M_j] = 0$; consequently
$\E[M_n^2] = \sum_{i=1}^n \E[(\Delta M_i)^2]$ when $M_0 = 0$.
\end{lemma}

\begin{proof}
For $i < j$, condition on $\filt_{j-1}$: $\E[\Delta M_i \Delta M_j]
= \E[\Delta M_i\, \E[\Delta M_j\mid\filt_{j-1}]] = 0$. Expanding
$M_n^2 = (\sum_i \Delta M_i)^2$ and using orthogonality kills the cross terms.
\end{proof}

\begin{lemma}[Weighted increment martingale is $L^2$-bounded]\label{lem:slln_weighted}
\uses{lem:slln_orthogonal}
Let $(M_n)$ be a square-integrable martingale with $M_0=0$ and
$\sum_{k\ge1}\frac{\E[(\Delta M_k)^2]}{k^2} < \infty$. Then
$S_n := \sum_{k=1}^n \frac{\Delta M_k}{k}$ is a square-integrable martingale with
$\sup_n \E[S_n^2] = \sum_{k\ge1}\frac{\E[(\Delta M_k)^2]}{k^2} < \infty$.
\end{lemma}

\begin{proof}
Each increment $\Delta S_k = \Delta M_k / k$ is $\filt_k$-measurable with
$\E[\Delta S_k\mid\filt_{k-1}] = 0$, so $S_n$ is a martingale. By
Lemma~\ref{lem:slln_orthogonal}, $\E[S_n^2] = \sum_{k\le n}\E[(\Delta M_k)^2]/k^2$,
which is bounded by the convergent series.
\end{proof}

\begin{lemma}[$L^2$-bounded martingales converge a.s.]\label{lem:slln_l2_conv}
A square-integrable martingale $(S_n)$ with $\sup_n \E[S_n^2] < \infty$ converges
a.s.\ to an integrable limit.
\end{lemma}

\begin{proof}
This is the $L^2$ (hence $L^1$-bounded) martingale convergence theorem, available
in Mathlib as \texttt{MeasureTheory.Submartingale.ae\_tendsto\_limitProcess}
applied to the $L^1$-bounded martingale $S_n$.
\end{proof}

\begin{lemma}[Kronecker's lemma]\label{lem:kronecker}
Let $(b_n)$ be a positive nondecreasing sequence with $b_n \to \infty$, and let
$(x_n)$ be reals such that $\sum_n x_n$ converges. Then
$\frac{1}{b_n}\sum_{k=1}^n b_k x_k \to 0$.
\end{lemma}

\begin{proof}
Let $s_n = \sum_{k\le n} x_k \to s$. Abel summation gives
$\frac{1}{b_n}\sum_{k\le n} b_k x_k = s_n - \frac{1}{b_n}\sum_{k=1}^{n-1}(b_{k+1}-b_k)s_k$.
The weights $(b_{k+1}-b_k)/b_n$ are nonnegative, sum to $(b_n - b_1)/b_n \to 1$,
and concentrate on large $k$; since $s_k \to s$, the weighted average tends to
$s$, so the whole expression tends to $s - s = 0$.
\end{proof}

\textbf{Formalization note.} Kronecker's lemma is not in Mathlib (only unrelated
``Kronecker'' matrix results exist). It is a short, self-contained real-analysis
lemma and a natural standalone contribution.

\section{The theorem}
\label{sec:pre_slln_thm}

\begin{theorem}[Martingale SLLN]\label{thm:mart_slln}
\uses{lem:slln_weighted, lem:slln_l2_conv, lem:kronecker}
Let $(M_n)$ be a square-integrable martingale with $M_0 = 0$ and
$\sum_{k\ge1}\frac{\E[(\Delta M_k)^2]}{k^2} < \infty$. Then $\frac{M_n}{n}\to 0$
a.s.
\end{theorem}

\begin{proof}
By Lemmas~\ref{lem:slln_weighted} and \ref{lem:slln_l2_conv}, the series
$S_n = \sum_{k\le n}\Delta M_k/k$ converges a.s. Apply Kronecker's lemma
(Lemma~\ref{lem:kronecker}) pathwise with $b_k = k$ and $x_k = \Delta M_k/k$:
$\frac{1}{n}\sum_{k\le n} k\cdot\frac{\Delta M_k}{k} = \frac{M_n}{n}\to 0$ a.s.
\end{proof}

\begin{corollary}[Bounded increments]\label{cor:mart_slln_bounded}
\uses{thm:mart_slln}
If $(M_n)$ is a martingale with $M_0=0$ and uniformly bounded increments
$|\Delta M_k| \le c$, then $\frac{M_n}{n}\to 0$ a.s.
\end{corollary}

\begin{proof}
$\E[(\Delta M_k)^2]\le c^2$, so $\sum_k \E[(\Delta M_k)^2]/k^2 \le c^2\sum_k k^{-2}
< \infty$; apply Theorem~\ref{thm:mart_slln}.
\end{proof}

\chapter{A one-sided law of the iterated logarithm}
\label{chap:pre_lil}

The blueprint's almost-sure rates all take the form
$M_n = O(\sqrt{\langle M\rangle_n \log\log\langle M\rangle_n})$. This is the
\emph{upper} half of the law of the iterated logarithm; the sharp constant is
never used. We therefore only formalize the one-sided statement, which is
considerably easier than the full Hartman--Wintner/Stout LIL.

\textbf{Formalization note.} Mathlib has \emph{no} law of the iterated logarithm.
However, \texttt{Probability/Moments/SubGaussian.lean} already provides
conditionally sub-Gaussian moment-generating-function bounds and an
Azuma--Hoeffding tail inequality for martingale differences
(\texttt{measure\_sum\_ge\_le\_of\_HasCondSubgaussianMGF}), together with the
Borel--Cantelli lemmas (\texttt{Probability/BorelCantelli.lean}). These are
exactly the ingredients of the argument below. The bounded-increment case is
within reach of the existing API; the general finite-variance case needs an
additional truncation step.

\section{Exponential tail bound}
\label{sec:pre_lil_tail}

\begin{lemma}[Exponential supermartingale]\label{lem:lil_exp_supermart}
\lean{AlphaRAR.supermartingale_expProcess}\leanok
Let $(M_n)$ be a martingale with $M_0=0$ and increments bounded by
$|\Delta M_i| \le c$. For $|\theta| \le 1/c$, the process
\[
  Z_n(\theta) := \exp\!\Big(\theta M_n - \theta^2 \langle M\rangle_n\Big)
\]
is a supermartingale, where $\langle M\rangle$ is the predictable quadratic
variation (Definition~\ref{def:pred_qv}).
\end{lemma}

\begin{proof}\leanok
\uses{def:pred_qv}
Using $e^x \le 1 + x + x^2$ for $|x| \le 1$ and $\E[\Delta M_i\mid\filt_{i-1}]=0$,
\[
  \E[e^{\theta\Delta M_i}\mid\filt_{i-1}]
  \le 1 + \theta^2\E[(\Delta M_i)^2\mid\filt_{i-1}]
  \le \exp\!\big(\theta^2(\langle M\rangle_i - \langle M\rangle_{i-1})\big).
\]
Hence $\E[Z_i(\theta)\mid\filt_{i-1}] \le Z_{i-1}(\theta)$.
\end{proof}

\begin{lemma}[Freedman-type inequality]\label{lem:lil_freedman}
\lean{AlphaRAR.measure_exists_ge_le_exp}
\uses{lem:lil_exp_supermart}
Under the hypotheses of Lemma~\ref{lem:lil_exp_supermart}, for all $\lambda > 0$
and $v > 0$,
\[
  \Prob\!\Big(\exists n:\ M_n \ge \lambda \ \text{and}\ \langle M\rangle_n \le v\Big)
  \le \exp\!\Big(-\frac{\lambda^2}{2(v + c\lambda)}\Big).
\]
The formalized statement \texttt{AlphaRAR.measure\_exists\_ge\_le\_exp} is the pre-optimization,
finite-horizon form: for each $n$ and each admissible $0<\theta\le 1/c$,
$\Prob(\exists k\le n:\ M_k\ge\lambda,\ \langle M\rangle_k\le v)\le \exp(-\theta\lambda+\theta^2 v)$,
obtained from Ville's maximal inequality
(\texttt{AlphaRAR.smul\_measure\_sup\_le\_integral\_zero}) applied to the exponential
supermartingale. Optimizing over $\theta$ and letting $n\to\infty$ recovers the displayed bound.
\end{lemma}

\begin{proof}
Optional stopping applied to the supermartingale $Z_n(\theta)$ at the first time
$M_n\ge\lambda$ with $\langle M\rangle_n\le v$ gives
$\Prob(\cdots)\le \exp(-\theta\lambda + \theta^2 v)$; optimize over
$0<\theta\le 1/c$. Equivalently, this is the Azuma--Hoeffding/Freedman bound that
follows from the conditionally sub-Gaussian MGF machinery of
\texttt{Probability/Moments/SubGaussian.lean}.
\end{proof}

\begin{corollary}[Optimized Freedman bound]\label{cor:lil_freedman_opt}
\lean{AlphaRAR.measure_exists_ge_le_exp_optimized, AlphaRAR.measure_exists_ge_le_exp_all}\leanok
\uses{lem:lil_freedman}
Under the hypotheses of Lemma~\ref{lem:lil_freedman}, for $\lambda, v > 0$ with the
admissibility condition $\lambda c \le 2v$ (so that the optimizer $\theta=\lambda/(2v)$
satisfies $\theta c\le 1$),
\[
  \Prob\!\Big(\exists n:\ M_n \ge \lambda \ \text{and}\ \langle M\rangle_n \le v\Big)
  \le \exp\!\Big(-\frac{\lambda^2}{4v}\Big).
\]
The variance proxy $2v$ (rather than the sharp $v$) comes from the one-step bound
$e^x\le 1+x+x^2$, costing a factor $\sqrt2$ in the eventual constant.
\end{corollary}

\begin{proof}\leanok
\uses{lem:lil_freedman}
Evaluate Lemma~\ref{lem:lil_freedman} at $\theta=\lambda/(2v)$, giving exponent
$-\theta\lambda+\theta^2 v = -\lambda^2/(4v)$. The finite-horizon events increase to the
event over all $n$, so their measures pass to the limit
(\texttt{measure\_exists\_ge\_le\_exp\_all}).
\end{proof}

\begin{lemma}[Borel--Cantelli over blocks]\label{lem:lil_bc}
\lean{AlphaRAR.ae_eventually_forall_lt_of_summable}\leanok
\uses{cor:lil_freedman_opt}
Fix threshold schedules $(v_k)_{k}$, $(\lambda_k)_k$ with $v_k,\lambda_k>0$ and
$\lambda_k c\le 2v_k$, such that $\sum_k \exp(-\lambda_k^2/(4v_k))<\infty$. Then almost
surely only finitely many blocks are ``bad'': for a.e.\ $\omega$, for all large $k$ and
every $n$, $\langle M\rangle_n\le v_k \Rightarrow M_n < \lambda_k$.
\end{lemma}

\begin{proof}\leanok
\uses{cor:lil_freedman_opt}
By Corollary~\ref{cor:lil_freedman_opt}, $\Prob(s_k)\le\exp(-\lambda_k^2/(4v_k))$ where
$s_k=\{\exists n:\ M_n\ge\lambda_k,\ \langle M\rangle_n\le v_k\}$. The right-hand sides are
summable, so $\sum_k\Prob(s_k)<\infty$, and the first Borel--Cantelli lemma gives that a.e.\
$\omega$ lies in only finitely many $s_k$.
\end{proof}

\section{The upper bound}
\label{sec:pre_lil_upper}

\begin{theorem}[One-sided LIL, bounded increments]\label{thm:lil_bounded}
\uses{lem:lil_bc}
Let $(M_n)$ be a martingale with $M_0=0$ and $|\Delta M_i|\le c$, and suppose
$\langle M\rangle_n \to \infty$ a.s. Then
\[
  \limsup_{n\to\infty}
  \frac{M_n}{\sqrt{2\langle M\rangle_n \log\log\langle M\rangle_n}} \le 1
  \quad\text{a.s.},
\]
and in particular $M_n = O\!\big(\sqrt{\langle M\rangle_n \log\log\langle M\rangle_n}\big)$
a.s.
\end{theorem}

\begin{proof}
\uses{lem:lil_bc}
Fix $\eta>1$ and set $v_k = \eta^k$. Apply Lemma~\ref{lem:lil_bc} with
$\lambda_k=(1+\varepsilon)\sqrt{2v_k\log\log v_k}$: the tail bounds
$\exp(-\lambda_k^2/(4v_k))$ are summable in $k$, so a.s.\ only finitely many blocks are
bad. On the block $\{\langle M\rangle_n\in(v_{k-1},v_k]\}$ this yields
$M_n < \lambda_k \le C\sqrt{\langle M\rangle_n\log\log\langle M\rangle_n}$; letting
$\eta\downarrow1$ and $\varepsilon\downarrow0$ yields the $\limsup$ bound.

\textbf{Formalization status.} The formalized statement
\texttt{AlphaRAR.ae\_eventually\_forall\_lt\_dyadic} is the dyadic-block form: with the schedule
$v_k=2^k$ and $\lambda_k = K\sqrt{v_k(k+1)}$ (where $0<K$ and $Kc\le2$, so admissibility
$\lambda_k c\le 2v_k$ holds for \emph{every} $k$ via $k+1\le 2^k$, and the tail bounds
$\exp(-\lambda_k^2/(4v_k))=\exp(-K^2(k+1)/4)$ form a geometric series), almost surely for all
large $k$ and every $n$, $\langle M\rangle_n\le 2^k \Rightarrow M_n < K\sqrt{2^k(k+1)}$. This is
the complete probabilistic content, built from Ville's inequality, the optimized Freedman bound
(Corollary~\ref{cor:lil_freedman_opt}), and Borel--Cantelli (Lemma~\ref{lem:lil_bc}). Since on the
block $\langle M\rangle_n\in(2^{k-1},2^k]$ one has $\sqrt{2^k(k+1)}\asymp
\sqrt{\langle M\rangle_n\log\langle M\rangle_n}$, it gives the displayed
$O(\sqrt{\langle M\rangle_n\log\langle M\rangle_n})$ upper bound, formalized in
Corollary~\ref{cor:lil_upper} below. One deliberate deviation from the displayed statement remains
unformalized: the $\log\log$ (rather than $\log$) rate with sharp constant $1$, which would require
refining the one-step MGF bound $e^x\le1+x+x^2$ to variance proxy $1$ (valid only as the increment
$\to0$). The $\log$ rate already gives $M_n=o(\langle M\rangle_n)$, which suffices for all
downstream uses.
\end{proof}

\begin{corollary}[Normalized one-sided upper bound]\label{cor:lil_upper}
\uses{thm:lil_bounded}
Under the hypotheses of Theorem~\ref{thm:lil_bounded} (in particular $\langle M\rangle_n\to\infty$
a.s.), there is a deterministic constant $C$ (depending only on $c$) such that almost surely
\[
  M_n \le C\,\sqrt{\langle M\rangle_n\,\log\langle M\rangle_n}\qquad\text{for all large }n.
\]
This is the consumable $O(\sqrt{\langle M\rangle_n\log\langle M\rangle_n})$ a.s.\ upper bound used
downstream (Lemma~\ref{lem:U_increment_bound}, via the bounded assignment martingale).
\end{corollary}

\begin{proof}
\uses{thm:lil_bounded}
Repackage Theorem~\ref{thm:lil_bounded}: for large $n$ let $k$ be the least index with
$\langle M\rangle_n\le 2^k$. Minimality gives $2^k\le 2\langle M\rangle_n$ and
$k\le\log_2\langle M\rangle_n+1$, so
$K\sqrt{2^k(k+1)}\le C\sqrt{\langle M\rangle_n\log\langle M\rangle_n}$ with
$C=K\sqrt{2(1/\log 2+2)}$, using $\log\langle M\rangle_n\ge1$ eventually.
\end{proof}

\begin{corollary}[Upper bound at scale $\sqrt{n\log n}$]\label{cor:lil_upper_nat}
\uses{cor:lil_upper, lem:qv_linear_bound}
Under the hypotheses of Corollary~\ref{cor:lil_upper}, almost surely
$M_n \le C\sqrt{n\log n}$ for all large $n$, for a deterministic constant $C$. This is the form
$M_n = O(\sqrt{n\log n})$ a.s.\ used for the bounded assignment martingale
$M_{\cdot,k}$ in Lemma~\ref{lem:U_increment_bound}.
\end{corollary}

\begin{proof}
\uses{cor:lil_upper, lem:qv_linear_bound}
By Lemma~\ref{lem:qv_linear_bound}, $\langle M\rangle_n\le c^2 n$, hence
$\log\langle M\rangle_n\le\log(c^2 n)\le 2\log n$ once $n\ge c^2$. Substituting into
Corollary~\ref{cor:lil_upper} bounds $\sqrt{\langle M\rangle_n\log\langle M\rangle_n}
\le\sqrt{2c^2}\,\sqrt{n\log n}$.
\end{proof}

\section{The finite-variance case via truncation}
\label{sec:pre_lil_trunc}

The remaining case is a martingale $M_n = \sum_{i<n}\xi_i D_i$ whose increments are
\emph{not} bounded: $\xi_i$ are i.i.d.\ centered with only $\E[\xi_0^2]=\sigma^2<\infty$,
and $D_i\in\{0,1\}$ is predictable (each $\xi_i$ is independent of the past $\filt_{i-1}$,
while $D_i$ is $\filt_{i-1}$-measurable). Then $\langle M\rangle_n=\sigma^2 N_n$ with
$N_n:=\sum_{i<n}D_i$. Theorem~\ref{thm:lil_bounded} does not apply directly, so we
truncate at level $b_i:=\sqrt{i}$ --- the largest level for which the tail is still
summable under only a second moment --- and control three pieces separately. This section
records the decomposition into small steps; each is a candidate Lean lemma to be formalized
in turn. As throughout, we prove the $\log$ (not $\log\log$) rate.

\begin{definition}[Truncated decomposition]\label{def:lil_trunc}
With $b_i:=\sqrt i$, set
$\xi_i^{\mathrm t}:=\xi_i\indic_{\{|\xi_i|\le b_i\}}$,
$\xi_i^{\mathrm r}:=\xi_i\indic_{\{|\xi_i|> b_i\}}$,
$m_i:=\E[\xi_i^{\mathrm t}]$, $\tilde\xi_i:=\xi_i^{\mathrm t}-m_i$, and
\[
  \tilde M_n:=\sum_{i<n}\tilde\xi_i D_i,\quad
  R_n:=\sum_{i<n}\xi_i^{\mathrm r} D_i,\quad
  \mathrm{Dr}_n:=\sum_{i<n} m_i D_i,\quad
  N_n:=\sum_{i<n} D_i,
\]
so that $M_n=\tilde M_n+R_n+\mathrm{Dr}_n$.
\end{definition}

\subsection{The tail term $R_n$}

\begin{lemma}[Summable tail probabilities]\label{lem:trunc_tail_summable}
\lean{AlphaRAR.tsum_measure_abs_sub_gt_sqrt_ne_top}\leanok
\uses{def:lil_trunc}
$\sum_{i}\Prob(|\xi_i|>\sqrt i)<\infty$. (The formalized statement
\texttt{tsum\_measure\_abs\_sub\_gt\_sqrt\_ne\_top} is the law-level version: for any finite measure
$\rho$ on $\R$ with $\int(x-\theta)^2\,d\rho<\infty$, $\sum_i\rho(|x-\theta|>\sqrt i)<\infty$; apply
it to the reward law $\rho=\nu_k$, $\theta=\theta_k$. This rests on a general-measure layer-cake
bound \texttt{tsum\_measure\_Ioi\_ne\_top}, the $\mu$-general form of Mathlib's
\texttt{tsum\_prob\_mem\_Ioi\_lt\_top}.)
\end{lemma}
\begin{proof}\leanok
Since $|\xi_i|>\sqrt i \iff (x-\theta_k)^2>i$, the sum equals $\sum_i\nu_k((\cdot-\theta_k)^2>i)$,
which is finite by the layer-cake bound $\sum_i\mu\{X>i\}\le\int X\,d\mu+\mu(\text{univ})$ applied to
$X=(\cdot-\theta_k)^2$ (integrable, being the arm-$k$ second central moment $\sigma^2$).
\end{proof}

\begin{lemma}[Tail is eventually zero]\label{lem:trunc_tail_const}
\lean{AlphaRAR.ae_eventually_abs_le_of_tsum_ne_top,
AlphaRAR.measure_action_and_tail_le, AlphaRAR.tsum_measure_action_and_tail_ne_top,
AlphaRAR.ae_eventually_tailRespPart_const}\leanok
\uses{lem:trunc_tail_summable}
Almost surely $|\xi_i|\le b_i$ for all large $i$; consequently $\xi_i^{\mathrm r}=0$ eventually and
$R_n$ is a.s.\ eventually constant, so $R_n=O(1)$ a.s.

The formalization has two layers. The abstract eventually-bounded clause
\texttt{ae\_eventually\_abs\_le\_of\_tsum\_ne\_top} is stated for a general level sequence $b$ and
takes $\sum_i\Prob(|\xi_i|>b_i)<\infty$ as hypothesis. The sampled version
\texttt{ae\_eventually\_tailRespPart\_const} concludes that the tail remainder
$R_n=Q_n-\tilde M_n-\mathrm{Dr}_n$ (\texttt{tailRespPart}) is a.s.\ eventually constant, via a
\emph{conditional bookkeeping} step: the sampled tail event has
$\Prob(A_i=k,\ \sqrt i\le|Y_i-\theta_k|)\le\nu_k(\sqrt i\le|\cdot-\theta_k|)$
(\texttt{measure\_action\_and\_tail\_le}, by conditioning on $\filtAction_i$ and
\texttt{condExp\_feedback\_comp}), so these probabilities are summable
(\texttt{tsum\_measure\_action\_and\_tail\_ne\_top}, using the closed-window law-level bound
\texttt{tsum\_measure\_abs\_sub\_ge\_sqrt\_ne\_top}).
\end{lemma}
\begin{proof}\leanok
\uses{lem:trunc_tail_summable}
By Lemma~\ref{lem:trunc_tail_summable} (in the sampled form above) the summability hypothesis holds,
and the first Borel--Cantelli lemma (\texttt{ae\_eventually\_notMem}) gives a.s.\ that for all large
$i$ arm $k$ is not sampled with an off-window response, whence $\xi_i^{\mathrm r}=0$ and the
increments of $R_n$ (\texttt{tailRespPart\_succ\_sub}, where the truncated mean cancels) vanish for
large $i$; a sequence with eventually-vanishing increments is eventually constant.
\end{proof}

\subsection{The drift term $\mathrm{Dr}_n$}

\begin{lemma}[Truncated-mean bound]\label{lem:trunc_mean_bound}
\lean{AlphaRAR.abs_integral_truncation_le}\leanok
\uses{def:lil_trunc}
$|m_i|\le \sigma^2/\sqrt i$ for $i\ge1$. (The formalized statement
\texttt{abs\_integral\_truncation\_le} is the abstract single-variable bound: for a centred $X$
with $X^2$ integrable, $|\E[\operatorname{trunc}(X,A)]|\le \E[X^2]/A$; apply it to $X=\xi_0$ and
$A=\sqrt i$.)
\end{lemma}
\begin{proof}\leanok
Since $\E[\xi_0]=0$, $m_i=\E[\operatorname{trunc}(\xi_0,\sqrt i)]
=\E[\operatorname{trunc}(\xi_0,\sqrt i)-\xi_0]$, so by the pointwise bound
$|\operatorname{trunc}(\xi_0,A)-\xi_0|\le\xi_0^2/A$ (the truncation is either $\xi_0$ or, off the
window where $|\xi_0|\ge A$, is $0$ with $|\xi_0|\le\xi_0^2/A$),
$|m_i|\le\E[\xi_0^2]/\sqrt i=\sigma^2/\sqrt i$.
\end{proof}

\begin{lemma}[Drift is $O(\sqrt{n})$]\label{lem:trunc_drift}
\lean{AlphaRAR.abs_truncDrift_le}\leanok
\uses{lem:trunc_mean_bound}
$|\mathrm{Dr}_n|\le 2\sigma^2\sqrt{n}$ (everywhere, with $\sigma^2=V_k$). In particular
$\mathrm{Dr}_n=O(\sqrt n)=o(\sqrt{N_n\log N_n})$ a.s.\ whenever $N_n$ grows linearly (the case of
interest). The sharper bound $|\mathrm{Dr}_n|\le 2\sigma^2\sqrt{N_n}$ holds by reindexing to the
sampled subsequence $i_j\ge j$, but is not needed and not formalized.
\end{lemma}
\begin{proof}\leanok
\uses{lem:trunc_mean_bound}
By Lemma~\ref{lem:trunc_mean_bound}, $|m_i|\le\sigma^2/\sqrt i$ (and $m_0=0$), so
$|\mathrm{Dr}_n|\le\sum_{i<n}|m_i|D_i\le\sum_{i<n}|m_i|\le\sigma^2\sum_{i<n}\tfrac1{\sqrt i}
\le 2\sigma^2\sqrt n$, using $\sum_{i<n}1/\sqrt i\le 2\sqrt n$
(\texttt{sum\_one\_div\_sqrt\_le}).
\end{proof}

\subsection{The centred truncated martingale $\tilde M_n$}

\begin{lemma}[$\tilde M$ is a martingale]\label{lem:trunc_mart}
\lean{AlphaRAR.martingale_truncRespMart}\leanok
\uses{def:lil_trunc, lem:Q_martingale}
$\tilde M$ is a martingale with $\tilde M_0=0$: $\E[\tilde\xi_i D_i\mid\filt_{i-1}]
=D_i\,\E[\tilde\xi_i\mid\filt_{i-1}]=D_i(\E[\xi_i^{\mathrm t}]-m_i)=0$, using that $\xi_i$ is
independent of $\filt_{i-1}$ and $D_i$ is $\filt_{i-1}$-measurable.
(Formalized in the \texttt{IsAlgEnvSeq}/\texttt{stationaryEnv} framework: $\tilde M$ is the
\emph{truncated response martingale} $\tilde Q_{n,k}=\sum_{i<n}\indic_{\{A_i=k\}}
(\operatorname{trunc}(Y_i-\theta_k,\sqrt i)-m_i)$, a martingale for the action-augmented
filtration $\filtAction$. It is an instance of a general functional response martingale
\texttt{genRespMart} --- of which the response martingale $Q$ (Lemma~\ref{lem:Q_martingale}) is the
$g=\mathrm{id}$ case --- whose martingale property follows from \texttt{condExp\_feedback\_comp}.)
\end{lemma}
\begin{proof}\leanok
\uses{def:lil_trunc}
Pull out the $\filtAction_i$-measurable indicator $\indic_{\{A_i=k\}}$; the conditional
expectation of $g(Y_i)$ (with $g=\operatorname{trunc}(\cdot-\theta_k,\sqrt i)$) given
$\filtAction_i$ is $(\nu(A_i))[g]=m_i$ on $\{A_i=k\}$ (\texttt{condExp\_feedback\_comp}), so the
centred increment has conditional mean $m_i-m_i=0$.
\end{proof}

\begin{lemma}[Quadratic variation of $\tilde M$]\label{lem:trunc_qv}
\lean{AlphaRAR.predQuadVar_truncRespMart_eq}\leanok
\uses{lem:trunc_mart, def:pred_qv, lem:qv_incr}
$\langle\tilde M\rangle_n=\sum_{i<n}v_i D_i$, where $D_i=\indic_{\{A_i=k\}}$ and
$v_i:=\int(\operatorname{trunc}(x-\theta_k,\sqrt i)-m_i)^2\,d\nu_k$ is the truncated-response
variance of arm $k$.
\end{lemma}
\begin{proof}\leanok
\uses{lem:trunc_mart, def:pred_qv, lem:qv_incr}
Instance of the general functional quadratic variation
\texttt{predQuadVar\_genRespMart\_eq}: the compensator increment is the $\filtAction_i$-conditional
second moment $\E[(\tilde\xi_i D_i)^2\mid\filtAction_i]=v_i D_i$
(\texttt{condExp\_genRespMart\_increment\_sq}), and these sum to $\sum_{i<n}v_i D_i$. (Unlike the
response martingale $Q$, whose constant variance gives $V_k N$, the truncated variances $v_i$ vary
with $i$, so the weighted sum is retained.)
\end{proof}

\begin{lemma}[Quadratic variation upper bound]\label{lem:trunc_qv_le}
\lean{AlphaRAR.predQuadVar_truncRespMart_le}\leanok
\uses{lem:trunc_qv}
$\langle\tilde M\rangle_n\le V_k N_{n,k}$ a.s., since each $v_i\le V_k$ (truncation reduces the
absolute value, and variance is at most the second moment). This is the only bound on
$\langle\tilde M\rangle$ used by the block LIL argument (the matching lower bound and the
convergence $v_i\to\sigma^2$ of Lemma~\ref{lem:trunc_var_conv} are not needed for the $O$-rate).
\end{lemma}
\begin{proof}\leanok
\uses{lem:trunc_qv}
By Lemma~\ref{lem:trunc_qv}, $\langle\tilde M\rangle_n=\sum_{i<n}v_i D_i\le V_k\sum_{i<n}D_i
=V_k N_{n,k}$, using $v_i=\operatorname{Var}_{\nu_k}(\operatorname{trunc})\le
\E_{\nu_k}[\operatorname{trunc}^2]\le\E_{\nu_k}[(\cdot-\theta_k)^2]=V_k$.
\end{proof}

\begin{lemma}[Truncated variances converge]\label{lem:trunc_var_conv}
\uses{lem:trunc_qv}
(Not needed for the $O$-rate.) $v_i\to\sigma^2$; in particular $\tfrac12\sigma^2\le v_i\le\sigma^2$
for all large $i$, so
$\tfrac12\sigma^2 N_n\le\langle\tilde M\rangle_n\le\sigma^2 N_n$ eventually and
$\langle\tilde M\rangle_n\asymp\sigma^2 N_n$.
\end{lemma}
\begin{proof}
$v_i=\E[(\xi_i^{\mathrm t})^2]-m_i^2$. As $\sqrt i\to\infty$, $\E[(\xi_i^{\mathrm t})^2]
=\int\operatorname{trunc}(x-\theta_k,\sqrt i)^2\,d\nu_k\to\int(x-\theta_k)^2\,d\nu_k=\sigma^2$
by monotone/dominated convergence, and $m_i\to0$ (Lemma~\ref{lem:trunc_mean_bound}), so
$v_i\to\sigma^2$. The upper bound $v_i\le\sigma^2$ is immediate from
$|\operatorname{trunc}(\xi_i,\sqrt i)|\le|\xi_i|$.
\end{proof}

\subsection{Core: the growing-increment LIL for $\tilde M$}
\label{sec:pre_lil_grow}

The increments $|\tilde\xi_i D_i|\le 2b_i=2\sqrt i$ \emph{grow}, so Theorem~\ref{thm:lil_bounded}
(fixed $c$) is unavailable. We rerun the dyadic argument with a block-dependent increment bound.
On block $k$ we use horizon $2^k$, increment bound $c_k:=2\sqrt{2^k}$, and the \emph{constrained}
exponential parameter $\theta_k:=1/c_k$ (the optimizer is inadmissible). The point is that
$\theta_k^2\langle\tilde M\rangle_{2^k}\le\theta_k^2\sigma^2 2^k=\sigma^2/4$ stays bounded while
$\theta_k\lambda_k$ grows like $\sqrt k$.

\begin{lemma}[Per-block Freedman with constrained $\theta$]\label{lem:trunc_block}
\uses{lem:trunc_mart, lem:lil_freedman}
Fix $C>0$ and set $\lambda_k:=C\sqrt{2^k\,k}$. For each $k$,
\[
  \Prob\!\Big(\exists\,m\le 2^k:\ \tilde M_m\ge\lambda_k,\ \langle\tilde M\rangle_m\le\sigma^2 2^k\Big)
  \le \exp\!\big(-\tfrac{C}{2}\sqrt k+\tfrac{\sigma^2}{4}\big).
\]
Here $\sigma^2=V_k$ is the arm variance and $\langle\tilde M\rangle_m\le V_k 2^k$ uses only the
upper bound Lemma~\ref{lem:trunc_qv_le}.
\end{lemma}
\begin{proof}
\uses{lem:lil_freedman, lem:trunc_qv_le}
Apply the Freedman inequality (Lemma~\ref{lem:lil_freedman}) to the martingale stopped at time
$2^k$ --- whose increments are bounded by $c_k=2\sqrt{2^k}$, so $\theta_k c_k=1$ is admissible ---
with $\theta=\theta_k$, $\lambda=\lambda_k$, $v=\sigma^2 2^k$. The exponent is
$-\theta_k\lambda_k+\theta_k^2\sigma^2 2^k=-\tfrac C2\sqrt k+\tfrac{\sigma^2}4$, since
$\theta_k\lambda_k=\tfrac{C\sqrt{2^k k}}{2\sqrt{2^k}}=\tfrac C2\sqrt k$ and
$\theta_k^2\sigma^2 2^k=\tfrac{\sigma^2 2^k}{4\cdot2^k}=\tfrac{\sigma^2}4$. Stopping at $2^k$ is
needed so the increment bound holds up to the horizon; on the event only times $m\le 2^k$ occur,
where the stopped process agrees with $\tilde M$.

The formalization packages the stopping into a reusable device
\texttt{AlphaRAR.measure\_exists\_ge\_le\_exp\_horizon}: a version of the pre-optimization Freedman
bound (Lemma~\ref{lem:lil_freedman}) whose increment bound $|\Delta M_i|\le c$ is required only for
$i<N$, and whose event is restricted to $m\le N$. It is proved by applying
Lemma~\ref{lem:lil_freedman} to the deterministically stopped process
$M^N_m:=M_{\min(m,N)}$ (Mathlib's \texttt{stoppedProcess} at the constant time $N$), which is a martingale
(\texttt{martingale\_stoppedProcess\_const}), whose increments vanish past $N$ hence are globally bounded by
$c$, and whose predictable quadratic variation agrees with $\langle M\rangle$ below the horizon
(\texttt{predQuadVar\_stoppedProcess\_const\_of\_le}). Here it is used with $N=2^k$, $c=c_k=2\sqrt{2^k}$, the
constrained $\theta_k=1/c_k$, and the increment bound $|\Delta\tilde M_i|\le2\sqrt i\le c_k$ for
$i<2^k$ (\texttt{abs\_truncRespMart\_increment\_le}).
\end{proof}

\begin{lemma}[Block probabilities are summable]\label{lem:trunc_block_summable}
\lean{AlphaRAR.summable_exp_neg_mul_sqrt_add}\leanok
For every $C>0$, $\sum_k \exp(-\tfrac C2\sqrt k+\tfrac{\sigma^2}4)<\infty$.
\end{lemma}
\begin{proof}\leanok
$\exp(-\tfrac C2\sqrt k)$ is summable in $k$: since $\log=o(\sqrt\cdot)$, eventually
$2\log k\le \tfrac C2\sqrt k$, so the terms are $\le k^{-2}$.
\end{proof}

\begin{lemma}[LIL for the truncated martingale]\label{lem:trunc_mart_lil}
\lean{AlphaRAR.ae_eventually_truncRespMart_lt_block,
AlphaRAR.ae_eventually_truncRespMart_le_sqrt_nat_mul_log,
AlphaRAR.ae_eventually_abs_truncRespMart_le_sqrt_nat_mul_log}\leanok
\uses{lem:trunc_block, lem:trunc_block_summable, lem:trunc_qv_le}
$\tilde M_n=O(\sqrt{n\log n})$ a.s.

\textbf{Reusable tool.} The block argument is formalized \emph{abstractly}, for any $L^2$-martingale
$M$ with a square-root-growing increment bound $|\Delta M_i|\le a\sqrt i$ and a linear quadratic
variation bound $\langle M\rangle_n\le v\,n$: the general lemmas
\texttt{measure\_exists\_ge\_le\_exp\_block} (per-block Freedman),
\texttt{ae\_eventually\_lt\_block\_of\_growing} (block Borel--Cantelli),
\texttt{ae\_eventually\_le\_sqrt\_nat\_mul\_log\_of\_growing} (one-sided $O(\sqrt{n\log n})$) and
\texttt{ae\_eventually\_abs\_le\_sqrt\_nat\_mul\_log\_of\_growing} (\emph{two-sided}, via $\pm M$ and
$\langle -M\rangle=\langle M\rangle$) live in \texttt{LILTruncation.lean} and belong upstream. The
truncated martingale $\tilde M$ is the instance $a=2$, $v=V_k$; in particular the two-sided form
\texttt{ae\_eventually\_abs\_truncRespMart\_le\_sqrt\_nat\_mul\_log} ($|\tilde M_n|\le
C\sqrt{n\log n}$) is free.
\end{lemma}
\begin{proof}\leanok
\uses{lem:trunc_block, lem:trunc_block_summable, lem:trunc_qv_le}
Borel--Cantelli over the blocks (Lemmas~\ref{lem:trunc_block}--\ref{lem:trunc_block_summable})
gives that a.s.\ for all large $k$ and every $m\le 2^k$, $\langle\tilde M\rangle_m\le\sigma^2 2^k
\Rightarrow \tilde M_m<\lambda_k$ (formalized as
\texttt{ae\_eventually\_truncRespMart\_lt\_block}, the summable first Borel--Cantelli step). To
bound a fixed time $n$, take the \emph{least} $k$ with $n\le 2^k$; then $2^k\le 2n$ and
$k\le\log_2 n+1$, and the quadratic-variation bound $\langle\tilde M\rangle_n\le V_k N_{n,k}\le V_k
n\le V_k 2^k$ (Lemma~\ref{lem:trunc_qv_le}) supplies the block hypothesis at $m=n\le 2^k$, giving
$\tilde M_n<\lambda_k=C\sqrt{2^k k}\le C'\sqrt{n\log n}$.

\textbf{Scale.} Because the increment bound $|\Delta\tilde M_i|\le2\sqrt i$ \emph{grows}, the block
horizon $m\le 2^k$ forces the least admissible $k$ to satisfy $2^k\ge n$, so time-blocking delivers
the \emph{time} scale $O(\sqrt{n\log n})$ rather than the variance scale
$O(\sqrt{\langle\tilde M\rangle_n\log\langle\tilde M\rangle_n})=O(\sqrt{N_n\log N_n})$ of the
displayed statement. The two coincide when arm $k$ is sampled a positive fraction of the time
($N_{n,k}\asymp n$, the regime of interest); the sharper variance-scale bound would require blocking
by the quadratic variation rather than by the natural time, and is not formalized. The time-scale
form is exactly parallel to the bounded-increment Corollary~\ref{cor:lil_upper_nat}
(\texttt{ae\_eventually\_le\_sqrt\_nat\_mul\_log}).
\end{proof}

\subsection{Assembling}

\begin{lemma}[Truncation to finite variance]\label{lem:lil_truncation}
\lean{AlphaRAR.ae_eventually_abs_respMart_le_sqrt_nat_mul_log,
AlphaRAR.ae_eventually_respMart_le_sqrt_nat_mul_log}\leanok
\uses{lem:trunc_tail_const, lem:trunc_drift, lem:trunc_mart_lil}
$M_n = O(\sqrt{n\log n})$ a.s. The formalized statement is \emph{two-sided},
\texttt{ae\_eventually\_abs\_respMart\_le\_sqrt\_nat\_mul\_log} ($|M_n|\le C\sqrt{n\log n}$
eventually, a.s.), with the one-sided form its immediate corollary. The two-sidedness comes for free
from the reusable growing-increment LIL tool (see Lemma~\ref{lem:trunc_mart_lil}).
\end{lemma}
\begin{proof}\leanok
\uses{lem:trunc_tail_const, lem:trunc_drift, lem:trunc_mart_lil}
By Definition~\ref{def:lil_trunc}, $M_n=\tilde M_n+R_n+\mathrm{Dr}_n$. The main term is
$O(\sqrt{n\log n})$ (Lemma~\ref{lem:trunc_mart_lil}); the tail is $O(1)$
(Lemma~\ref{lem:trunc_tail_const}); the drift is $O(\sqrt{n})$ (Lemma~\ref{lem:trunc_drift}).
Both remainders are $o(\sqrt{n\log n})$, so $M_n=O(\sqrt{n\log n})$. In the regime of interest
$N_{n,k}\asymp n$, this is the $O(\sqrt{N_n\log N_n})=O(\sqrt{\langle M\rangle_n\log\langle
M\rangle_n})$ rate (using $\langle M\rangle_n=\sigma^2 N_n$) required downstream.
\end{proof}

\textbf{Remark.} Lemma~\ref{lem:lil_truncation} is exactly the shape used in
Lemma~\ref{lem:theta_LIL}: there $M = Q_{\cdot,k}$, $D_i = X_{i,k}$,
$\xi_i = \xi_{i,k}-\theta_k$, and $\langle Q_{\cdot,k}\rangle_n = V_k N_{n,k}$, so
the conclusion reads $Q_{n,k} = O(\sqrt{N_{n,k}\log N_{n,k}})$ a.s. The same
statement applied to the bounded martingale $M_{\cdot,k}$
(Definition~\ref{def:M}) covers the use in Lemma~\ref{lem:U_increment_bound}; that direction is
already formalized as Corollary~\ref{cor:lil_upper_nat}.

\textbf{Formalization status.} All of \S\ref{sec:pre_lil_trunc} is now formalized: the tail
(Lemmas~\ref{lem:trunc_tail_summable}--\ref{lem:trunc_tail_const}), the drift
(Lemmas~\ref{lem:trunc_mean_bound}--\ref{lem:trunc_drift}), the martingale and its quadratic
variation (Lemmas~\ref{lem:trunc_mart}--\ref{lem:trunc_qv_le}), the core growing-increment LIL
(Lemmas~\ref{lem:trunc_block}--\ref{lem:trunc_mart_lil}), and the assembling
Lemma~\ref{lem:lil_truncation} (\texttt{ae\_eventually\_respMart\_le\_sqrt\_nat\_mul\_log}). The core
reuses \texttt{measure\_exists\_ge\_le\_exp} on a \emph{deterministically stopped} martingale --- the
one genuinely new device, packaged as \texttt{measure\_exists\_ge\_le\_exp\_horizon} --- followed by
the Borel--Cantelli/repackaging pattern of \S\ref{sec:pre_lil_upper}. The whole chapter delivers the
time-scale $O(\sqrt{n\log n})$ form (see the scale remark in Lemma~\ref{lem:trunc_mart_lil});
the variance-scale $O(\sqrt{\langle M\rangle_n\log\langle M\rangle_n})$, which agrees in the regime
$N_{n,k}\asymp n$, would require blocking by the quadratic variation and is not formalized.

\chapter{A general martingale law of the iterated logarithm}
\label{chap:pre_llil}

Chapter~\ref{chap:pre_lil} proves a one-sided LIL at the \emph{$\log$} rate,
$M_n=O(\sqrt{\langle M\rangle_n\log\langle M\rangle_n})$, using a dyadic schedule
tuned so that the Borel--Cantelli tail bounds form a trivially summable
\emph{geometric} series. That $\log$ (rather than $\log\log$) factor is the price of
that convenience, and it suffices for every $o(\langle M\rangle_n)$ use in Part~I.

This chapter sharpens the rate to the true $\log\log$ and, more importantly,
repackages the result as a \emph{general} martingale law of the iterated logarithm:
a theorem about an arbitrary $L^2$-martingale whose increments are small on the LIL
scale, with no bandit or i.i.d.\ structure baked in. Everything is stated for its own
sake, as a candidate Mathlib contribution. The chapter reuses, unchanged, the
formalized engine of Chapter~\ref{chap:pre_lil}: the exponential supermartingale
(Lemma~\ref{lem:lil_exp_supermart}), Ville's maximal inequality, the Freedman tail
bound (Lemma~\ref{lem:lil_freedman}, Corollary~\ref{cor:lil_freedman_opt}), the
horizon-local variant \texttt{measure\_exists\_ge\_le\_exp\_horizon}, and the
Borel--Cantelli step (Lemma~\ref{lem:lil_bc}).

\medskip
\noindent\textbf{Where the $\log\log$ comes from.} The dyadic argument controls, for a
threshold schedule $(v_k,\lambda_k)$, the events $s_k=\{\exists n:\ M_n\ge\lambda_k,\
\langle M\rangle_n\le v_k\}$, whose measures the optimized Freedman bound estimates by
$\exp(-\lambda_k^2/(4v_k))$. With geometric blocks $v_k=\eta^k$ the two natural
thresholds differ only in what summability they demand:
\[
  \lambda_k=\kappa\sqrt{v_k\,\log v_k}\ \Longrightarrow\ \exp(-\tfrac{\kappa^2}{4}\log v_k)
  \ \text{(geometric)},\qquad
  \lambda_k=\kappa\sqrt{v_k\,\log\log v_k}\ \Longrightarrow\
  \exp(-\tfrac{\kappa^2}{4}\log\log v_k)\ \text{(a $p$-series)}.
\]
Since $\log\log\eta^k=\log(k\log\eta)$, the second right-hand side is
$(k\log\eta)^{-\kappa^2/4}$, summable exactly when $\kappa>2$. So the only real change
from Chapter~\ref{chap:pre_lil} is the block \emph{threshold} and the corresponding
switch from geometric to $p$-series summability. The one structural subtlety is that
the near-optimal exponential parameter $\theta_k=\lambda_k/(2v_k)\to0$, so its
admissibility $\theta_k c\le1$ holds only \emph{eventually} in $k$; the Borel--Cantelli
step must be relaxed accordingly.

\medskip
\noindent\textbf{Layout.} Section~\ref{sec:llil_bounded} does the bounded-increment
case (the workhorse, reusing Chapter~\ref{chap:pre_lil} verbatim);
Section~\ref{sec:llil_general} isolates the general hypothesis \textbf{(H)} that makes
the argument go through for growing increments;
Section~\ref{sec:llil_sharp} refines the one-step estimate to obtain the sharp
constant $\limsup\le1$; Section~\ref{sec:llil_forms} records the consumable normalized
forms; Section~\ref{sec:llil_hw} applies the general theorem, via truncation, to the
classical i.i.d.\ Hartman--Wintner LIL under a bare second moment; and
Section~\ref{sec:llil_skip} records Doob's optional-skipping lemma, which connects a
predictably-sampled i.i.d.\ sequence to a genuine i.i.d.\ partial sum.

Throughout, $(M_n)_{n\in\N}$ is an $L^2$-martingale with $M_0=0$ on a filtered
probability space $(\Omega,\filt,\Prob)$, and $\langle M\rangle$ is its predictable
quadratic variation (Definition~\ref{def:pred_qv}). We abbreviate
$s_n^2:=\langle M\rangle_n$.

\section{The bounded-increment case}
\label{sec:llil_bounded}

\begin{lemma}[$p$-series tail summability]\label{lem:llil_tail_summable}
\lean{AlphaRAR.summable_exp_neg_mul_log_add}\leanok
For every $\kappa>2$ and $\eta>1$,
$\sum_{k}\exp\!\big(-\tfrac{\kappa^2}{4}\log\log\eta^k\big)<\infty$. (Formalized with base $\eta=2$
and the surrogate $\log(k+2)\asymp\log\log 2^k$: for $1<p$, $\sum_k\exp(-p\log(k+2))<\infty$.)
\end{lemma}
\begin{proof}\leanok
For $k\ge1$, $\log\log\eta^k=\log(k\log\eta)$, so the $k$-th term equals
$(k\log\eta)^{-\kappa^2/4}=(\log\eta)^{-\kappa^2/4}\,k^{-\kappa^2/4}$. Since
$\kappa^2/4>1$ this is a convergent $p$-series (Mathlib
\texttt{summable\_one\_div\_nat\_rpow}); dropping the finitely many indices with
$\log\log\eta^k\le0$ does not affect summability.
\end{proof}

\begin{lemma}[Borel--Cantelli with eventual admissibility]\label{lem:llil_evt_adm}
\lean{AlphaRAR.ae_eventually_forall_lt_of_summable_eventually}\leanok
\uses{cor:lil_freedman_opt, lem:lil_bc}
Let $M$ be a square-integrable martingale with $M_0=0$ and $|\Delta M_i|\le c$ a.e.
Let $(v_k),(\lambda_k)$ be positive schedules such that
\begin{enumerate}
  \item $\lambda_k c\le 2v_k$ for all large $k$ \emph{(eventual admissibility)}, and
  \item $\sum_k\exp(-\lambda_k^2/(4v_k))<\infty$.
\end{enumerate}
Then almost surely, for all large $k$ and every $n$,
$\langle M\rangle_n\le v_k\Rightarrow M_n<\lambda_k$.
\end{lemma}
\begin{proof}\leanok
As in Lemma~\ref{lem:lil_bc}, for indices $k$ where admissibility holds the optimized
Freedman bound (Corollary~\ref{cor:lil_freedman_opt}, infinite-horizon form
\texttt{measure\_exists\_ge\_le\_exp\_all}) gives
$\Prob(s_k)\le\exp(-\lambda_k^2/(4v_k))$, where
$s_k=\{\exists n:\ M_n\ge\lambda_k,\ \langle M\rangle_n\le v_k\}$; for the finitely many
remaining $k$ use the trivial bound $\Prob(s_k)\le1$. The total $\sum_k\Prob(s_k)$ is
still finite, so the first Borel--Cantelli lemma (\texttt{ae\_eventually\_notMem})
places a.e.\ $\omega$ in only finitely many $s_k$; the complement is the claim. This is
Lemma~\ref{lem:lil_bc} with its ``$\forall k$'' admissibility weakened to
``$\forall^\infty k$''.
\end{proof}

\begin{lemma}[Bounded-increment $\log\log$ exceedance]\label{lem:llil_bounded_block}
\lean{AlphaRAR.ae_eventually_forall_lt_dyadic_loglog}\leanok
\uses{lem:llil_tail_summable, lem:llil_evt_adm}
Let $|\Delta M_i|\le c$ a.e.\ with $c>0$. Fix $\eta>1$ and $\kappa>2$, set $v_k=\eta^k$
and $\lambda_k=\kappa\sqrt{\eta^k\,\log\log\eta^k}$ (for $k$ large enough that
$\log\log\eta^k>0$). Then almost surely, for all large $k$ and every $n$,
$\langle M\rangle_n\le\eta^k\Rightarrow M_n<\kappa\sqrt{\eta^k\log\log\eta^k}$.
(Formalized with $\eta=2$ and the surrogate $\lambda_k=\kappa\sqrt{2^k\log(k+2)}$.)
\end{lemma}
\begin{proof}\leanok
Eventual admissibility: $\lambda_k c\le 2v_k$ reads
$\kappa c\sqrt{\log\log\eta^k/\eta^k}\le2$, which holds for large $k$ because
$\log\log\eta^k/\eta^k\to0$ (in the formalization, $\lambda_k c\le2\cdot2^k$ reduces to
$(\kappa c)^2\log(k+2)\le4\cdot2^k$, eventual since $\log(k+2)=o(2^k)$). Summability:
$\exp(-\lambda_k^2/(4v_k))=\exp(-\tfrac{\kappa^2}{4}\log\log\eta^k)$ is summable by
Lemma~\ref{lem:llil_tail_summable}. Apply Lemma~\ref{lem:llil_evt_adm}.
\end{proof}

\begin{theorem}[General bounded-increment LIL, $O$-rate]\label{thm:llil_bounded}
\lean{AlphaRAR.ae_eventually_le_sqrt_predQuadVar_mul_loglog}\leanok
\uses{lem:llil_bounded_block}
Let $M$ be a square-integrable martingale with $M_0=0$, $|\Delta M_i|\le c$ a.e., and
$\langle M\rangle_n\to\infty$ a.s. Then for every constant $C>2$, almost surely
\[
  M_n\le C\,\sqrt{\langle M\rangle_n\,\log\log\langle M\rangle_n}\qquad
  \text{for all large }n.
\]
\end{theorem}
\begin{proof}\leanok
Pick $\eta>1$ and $\kappa\in(2,C/\sqrt\eta)$ (possible once $\eta<(C/2)^2$), and use
Lemma~\ref{lem:llil_bounded_block}. For large $n$ let $k$ be the least index with
$\langle M\rangle_n\le\eta^k$; minimality gives $\eta^k\le\eta\langle M\rangle_n$, and
$\langle M\rangle_n\to\infty$ makes $\log\log\eta^k\le\log\log(\eta\langle M\rangle_n)
\le(1+o(1))\log\log\langle M\rangle_n$. Hence
$M_n<\kappa\sqrt{\eta^k\log\log\eta^k}\le\kappa\sqrt\eta\,(1+o(1))
\sqrt{\langle M\rangle_n\log\log\langle M\rangle_n}\le C\sqrt{\langle M\rangle_n
\log\log\langle M\rangle_n}$ for large $n$.
\end{proof}

\section{The general hypothesis: small increments on the LIL scale}
\label{sec:llil_general}

Bounded increments are the special case of a condition that also covers martingales
with \emph{growing} increments, provided they are negligible relative to the LIL scale
$\sqrt{s_n^2/\log\log s_n^2}$. This is the martingale analogue of the Lindeberg
condition and the natural hypothesis of the Hall--Heyde martingale LIL.

\begin{definition}[Condition \textbf{(H)}]\label{def:llil_H}
$M$ satisfies \textbf{(H)} with constant $A>0$ if there is an $\filt$-predictable
sequence $(c_i)$ with $|\Delta M_i|\le c_i$ a.e.\ and, for all large $i$,
\[
  c_i\ \le\ A\,\sqrt{\langle M\rangle_i\,/\,\log\log\langle M\rangle_i}\qquad\text{a.e.}
\]
\end{definition}

A fixed bound $|\Delta M_i|\le c$ satisfies \textbf{(H)} with \emph{every} $A>0$ once
$\langle M\rangle_n\to\infty$, since then $\sqrt{\langle M\rangle_i/\log\log\langle
M\rangle_i}\to\infty$. The point of \textbf{(H)} is that under it the near-optimal
$\theta$ is admissible \emph{and} small, so no truncation is needed inside the theorem.

\begin{lemma}[Quadratic-variation stopping]\label{lem:llil_qv_stop}
\uses{def:pred_qv}
For $v>0$ let $\tau_v:=\inf\{n:\ \langle M\rangle_{n+1}>v\}$. Because $\langle M\rangle$
is predictable ($\langle M\rangle_{n+1}$ is $\filt_n$-measurable), $\tau_v$ is a
stopping time, and on $\{n\le\tau_v\}$ one has $\langle M\rangle_n\le v$. The stopped
process $M^{\tau_v}_n:=M_{n\wedge\tau_v}$ is a martingale with
$\langle M^{\tau_v}\rangle_n=\langle M\rangle_{n\wedge\tau_v}$, and under \textbf{(H)}
with constant $A$ its increments obey $|\Delta M^{\tau_v}_i|\le C(v):=A\sqrt{v/\log\log v}$
a.e.\ for all $i$ (for $v$ large), since for $i<\tau_v$ we have $\langle M\rangle_i\le v$
and past $\tau_v$ the increments vanish. Moreover
\[
  \{\exists n:\ M_n\ge\lambda,\ \langle M\rangle_n\le v\}\ \subseteq\
  \{\exists n:\ M^{\tau_v}_n\ge\lambda,\ \langle M^{\tau_v}\rangle_n\le v\}.
\]
\end{lemma}
\begin{proof}
Predictability of $\langle M\rangle$ gives $\{\tau_v\ge n+1\}=\{\langle M\rangle_{n+1}
\le v\}\in\filt_n$, so $\tau_v$ is a stopping time (constructible with Mathlib's
\texttt{hittingBtwn}, the stopping-time API already used in the proof of Ville's
inequality). That $M^{\tau_v}$ is a martingale is Mathlib's
\texttt{Submartingale.stoppedProcess} applied to $\pm M$; the increment bound uses
monotonicity of $\langle M\rangle$ and \textbf{(H)}. The only genuinely new ingredient is
the quadratic-variation identity $\langle M^{\tau_v}\rangle=\langle M\rangle_{\cdot\wedge
\tau_v}$, the stopping-time analogue of the \emph{formalized} deterministic-horizon
\texttt{predQuadVar\_stopMart\_of\_le} of Chapter~\ref{chap:pre_lil}. For the inclusion: if
$\langle M\rangle_n\le v$ then $n\le\tau_v$, so $M^{\tau_v}_n=M_n$ and
$\langle M^{\tau_v}\rangle_n=\langle M\rangle_n$.
\end{proof}

\begin{lemma}[Per-block bound under \textbf{(H)}]\label{lem:llil_H_measure}
\uses{cor:lil_freedman_opt, lem:llil_qv_stop}
Suppose $M$ satisfies \textbf{(H)} with constant $A$. Fix $\kappa>0$ with
$\kappa A\le2$, set $v_k=\eta^k$ ($\eta>1$) and
$\lambda_k=\kappa\sqrt{v_k\log\log v_k}$. Then for all large $k$,
\[
  \Prob\big(\exists n:\ M_n\ge\lambda_k,\ \langle M\rangle_n\le v_k\big)\ \le\
  \exp\!\big(-\tfrac{\kappa^2}{4}\log\log v_k\big).
\]
\end{lemma}
\begin{proof}
By Lemma~\ref{lem:llil_qv_stop} the event is contained in the corresponding event for
the stopped martingale $M^{\tau_{v_k}}$, whose increments are globally bounded by
$C(v_k)=A\sqrt{v_k/\log\log v_k}$. Apply the optimized infinite-horizon Freedman bound
(Corollary~\ref{cor:lil_freedman_opt}) to $M^{\tau_{v_k}}$ with $c=C(v_k)$,
$\lambda=\lambda_k$, $v=v_k$: admissibility $\lambda_k C(v_k)\le2v_k$ reads
$\kappa\sqrt{v_k\log\log v_k}\cdot A\sqrt{v_k/\log\log v_k}=\kappa A\,v_k\le2v_k$, i.e.\
$\kappa A\le2$; and the exponent is
$-\lambda_k^2/(4v_k)=-\tfrac{\kappa^2}{4}\log\log v_k$.
\end{proof}

\begin{lemma}[\textbf{(H)} exceedance]\label{lem:llil_H_block}
\uses{lem:llil_H_measure, lem:llil_tail_summable}
Under the hypotheses of Lemma~\ref{lem:llil_H_measure} with additionally $\kappa>2$
(so $A<1$ and $A\le2/\kappa$), almost surely for all large $k$ and every $n$,
$\langle M\rangle_n\le\eta^k\Rightarrow M_n<\kappa\sqrt{\eta^k\log\log\eta^k}$.
\end{lemma}
\begin{proof}
The per-block bounds of Lemma~\ref{lem:llil_H_measure} are summable by
Lemma~\ref{lem:llil_tail_summable} ($\kappa>2$), so $\sum_k\Prob(s_k)<\infty$ with
$s_k=\{\exists n:\ M_n\ge\lambda_k,\ \langle M\rangle_n\le v_k\}$; the first
Borel--Cantelli lemma finishes as in Lemma~\ref{lem:llil_evt_adm}.
\end{proof}

\begin{theorem}[General martingale LIL under \textbf{(H)}, $O$-rate]\label{thm:llil_general}
\uses{lem:llil_H_block}
Let $M$ be a square-integrable martingale with $M_0=0$ satisfying \textbf{(H)} with some
constant $A<1$, and $\langle M\rangle_n\to\infty$ a.s. Then there is a constant $C$
(any $C>2$ suffices, once $A$ is small enough) such that almost surely
\[
  M_n\le C\,\sqrt{\langle M\rangle_n\,\log\log\langle M\rangle_n}\qquad
  \text{for all large }n.
\]
Bounded increments (Theorem~\ref{thm:llil_bounded}) are the case $A\to0$.
\end{theorem}
\begin{proof}
Choose $\kappa>2$ with $\kappa A\le2$ and $\eta>1$, apply
Lemma~\ref{lem:llil_H_block}, and repackage exactly as in
Theorem~\ref{thm:llil_bounded} (least $k$ with $\langle M\rangle_n\le\eta^k$).
\end{proof}

\section{The sharp constant}
\label{sec:llil_sharp}

The $O$-rate above uses the crude one-step estimate $e^x\le1+x+x^2$
(Chapter~\ref{chap:pre_lil}), whose variance proxy $2$ inflates the constant. Replacing
it by the second-order estimate, valid as the increment $\to0$, yields the sharp
$\limsup\le1$. This section is the ``definitive theorem'' upgrade; the $O$-rate of
Sections~\ref{sec:llil_bounded}--\ref{sec:llil_general} already suffices for every
almost-sure rate in Part~I.

\begin{lemma}[Refined elementary inequality]\label{lem:llil_refined_ineq}
For every $\delta>0$ there is $\eta_\delta>0$ such that $e^x\le1+x+\tfrac12(1+\delta)x^2$
for all $|x|\le\eta_\delta$.
\end{lemma}
\begin{proof}
Mathlib's \texttt{Real.exp\_bound} at $n=3$ gives, for $|x|\le1$,
$\big|e^x-(1+x+\tfrac12 x^2)\big|\le|x|^3\cdot\tfrac{4}{3!\cdot3}=\tfrac29|x|^3$, hence
$e^x\le1+x+\tfrac12 x^2+\tfrac29|x|^3\le1+x+\tfrac12\big(1+\tfrac49|x|\big)x^2$; take
$\eta_\delta:=\min(1,\tfrac94\delta)$. (So this lemma is a Mathlib one-liner, not new work.)
\end{proof}

\begin{lemma}[Refined conditional MGF bound]\label{lem:llil_refined_mgf}
\uses{lem:llil_refined_ineq, def:pred_qv}
If $|\Delta M_i|\le c$ a.e.\ and $|\theta|c\le\eta_\delta$, then
$\E[e^{\theta\Delta M_i}\mid\filt_i]\le\exp\!\big(\tfrac12(1+\delta)\theta^2
(\langle M\rangle_{i+1}-\langle M\rangle_i)\big)$ a.e.
\end{lemma}
\begin{proof}
As in Lemma~\ref{lem:lil_exp_supermart}'s one-step bound, but with
Lemma~\ref{lem:llil_refined_ineq} in place of $e^x\le1+x+x^2$: using
$\E[\Delta M_i\mid\filt_i]=0$ and $\E[(\Delta M_i)^2\mid\filt_i]=\langle M\rangle_{i+1}
-\langle M\rangle_i$, one gets
$\E[e^{\theta\Delta M_i}\mid\filt_i]\le1+\tfrac12(1+\delta)\theta^2
(\langle M\rangle_{i+1}-\langle M\rangle_i)\le\exp(\cdots)$.
\end{proof}

\begin{lemma}[Refined exponential supermartingale and Freedman bound]\label{lem:llil_refined_freedman}
\uses{lem:llil_refined_mgf}
Under $|\Delta M_i|\le c$ and $|\theta|c\le\eta_\delta$, the process
$Z_n=\exp\!\big(\theta M_n-\tfrac12(1+\delta)\theta^2\langle M\rangle_n\big)$ is a
supermartingale, and hence for $\lambda,v>0$,
\[
  \Prob\big(\exists n:\ M_n\ge\lambda,\ \langle M\rangle_n\le v\big)\ \le\
  \exp\!\big(-\theta\lambda+\tfrac12(1+\delta)\theta^2 v\big).
\]
Optimizing at $\theta=\lambda/((1+\delta)v)$ (admissible when
$\lambda c\le(1+\delta)v\,\eta_\delta$) gives the bound
$\exp\!\big(-\lambda^2/(2(1+\delta)v)\big)$.
\end{lemma}
\begin{proof}
Verbatim the proofs of Lemma~\ref{lem:lil_exp_supermart} (supermartingale, via
Lemma~\ref{lem:llil_refined_mgf}), Lemma~\ref{lem:lil_freedman} (Ville's inequality),
and Corollary~\ref{cor:lil_freedman_opt} (optimization), with the proxy $2v$ replaced
by $(1+\delta)v$.
\end{proof}

\begin{lemma}[Sharp exceedance]\label{lem:llil_sharp_block}
\uses{lem:llil_refined_freedman, lem:llil_qv_stop, lem:llil_tail_summable}
Let $M$ satisfy \textbf{(H)} with constant $A$ (or have bounded increments), and
$\langle M\rangle_n\to\infty$. Fix $\eta>1$ and $\varepsilon>\delta>0$, set $v_k=\eta^k$
and $\lambda_k=\sqrt{2(1+\varepsilon)\,v_k\log\log v_k}$. Then almost surely for all
large $k$ and every $n$,
$\langle M\rangle_n\le\eta^k\Rightarrow M_n<\sqrt{2(1+\varepsilon)\eta^k\log\log\eta^k}$.
\end{lemma}
\begin{proof}
Stop at $\tau_{v_k}$ (Lemma~\ref{lem:llil_qv_stop}); the stopped increments are bounded
by $C(v_k)=A\sqrt{v_k/\log\log v_k}$ (or by the fixed $c$). The optimizer
$\theta_k=\lambda_k/((1+\delta)v_k)=\sqrt{2(1+\varepsilon)\log\log v_k/v_k}/(1+\delta)$
satisfies $\theta_k C(v_k)=A\sqrt{2(1+\varepsilon)}/(1+\delta)$, which is $\le\eta_\delta$
for $A$ small (bounded case: $\theta_k c\to0$). Lemma~\ref{lem:llil_refined_freedman}
gives $\Prob(s_k)\le\exp(-\lambda_k^2/(2(1+\delta)v_k))
=\exp(-\tfrac{1+\varepsilon}{1+\delta}\log\log v_k)$, summable by
Lemma~\ref{lem:llil_tail_summable}'s argument since $(1+\varepsilon)/(1+\delta)>1$.
Borel--Cantelli finishes.
\end{proof}

\begin{theorem}[Sharp general martingale LIL]\label{thm:llil_sharp}
\uses{lem:llil_sharp_block}
Let $M$ be a square-integrable martingale with $M_0=0$, $\langle M\rangle_n\to\infty$
a.s., satisfying \textbf{(H)} (in particular any bounded-increment martingale). Then
\[
  \limsup_{n\to\infty}\frac{M_n}{\sqrt{2\,\langle M\rangle_n\,\log\log\langle M\rangle_n}}
  \ \le\ 1\qquad\text{a.s.},
\]
and, applied to $-M$ (same quadratic variation, Definition~\ref{def:pred_qv}), the two-sided
bound $\limsup_n|M_n|/\sqrt{2\langle M\rangle_n\log\log\langle M\rangle_n}\le1$ a.s.
\end{theorem}
\begin{proof}
Fix a countable family $\eta_j\downarrow1$, $\varepsilon_j\downarrow0$ with
$\varepsilon_j>\delta_j>0$. For each $j$, Lemma~\ref{lem:llil_sharp_block} and the
repackaging of Theorem~\ref{thm:llil_bounded} give, a.s.,
$\limsup_n M_n/\sqrt{2\langle M\rangle_n\log\log\langle M\rangle_n}\le
\sqrt{(1+\varepsilon_j)\eta_j}$. Intersect the countably many a.s.\ events and let
$j\to\infty$.
\end{proof}

\section{Consumable forms}
\label{sec:llil_forms}

\begin{corollary}[Two-sided normalized bound]\label{cor:llil_norm}
\uses{thm:llil_general}
Under the hypotheses of Theorem~\ref{thm:llil_general} (or \ref{thm:llil_bounded}),
almost surely $|M_n|\le C\sqrt{\langle M\rangle_n\log\log\langle M\rangle_n}$ for all
large $n$, some deterministic $C$. Apply Theorem~\ref{thm:llil_general} to $M$ and $-M$.
\end{corollary}

\begin{corollary}[Bound at scale $\sqrt{n\log\log n}$]\label{cor:llil_nat}
\uses{cor:llil_norm}
If moreover $\langle M\rangle_n\le K n$ for a constant $K$ (e.g.\ $|\Delta M_i|\le c$
gives $\langle M\rangle_n\le c^2 n$, Lemma~\ref{lem:qv_linear_bound}), then almost
surely $|M_n|\le C'\sqrt{n\log\log n}$ for all large $n$.
\end{corollary}
\begin{proof}
$\langle M\rangle_n\le Kn$ gives $\log\log\langle M\rangle_n\le\log\log(Kn)\le
\log\log n+o(1)$ for large $n$; substitute into Corollary~\ref{cor:llil_norm}.
\end{proof}

These are the $\log\log$ replacements for Corollaries~\ref{cor:lil_upper} and
\ref{cor:lil_upper_nat}: the bounded assignment martingale $M_{\cdot,k}$ of
Lemma~\ref{lem:U_increment_bound} is covered by Corollary~\ref{cor:llil_nat} via
Theorem~\ref{thm:llil_bounded}.

\section{Application: the i.i.d.\ Hartman--Wintner LIL}
\label{sec:llil_hw}

The general theorem needs \textbf{(H)}; a sum of i.i.d.\ increments with a fixed heavy
law does not satisfy it and must be truncated. This is the classical Hartman--Wintner
LIL \cite{HartmanWintner1941}: it holds under a bare second moment, and that sharpness
is exactly what makes it delicate. A single truncation at the LIL scale
$\sqrt{j/\log\log j}$ makes the main part satisfy \textbf{(H)} but leaves a tail whose
one-term exceedance probabilities are \emph{not} summable ($\sum_j\Prob(|Y_j|>
\sqrt{j/\log\log j})=\infty$ when only $\E[Y^2]<\infty$), so the pathwise
Borel--Cantelli route of Chapter~\ref{chap:pre_lil} is unavailable.

The resolution is a \emph{three-level} truncation. Cut at two scales: the LIL scale
$b_j\asymp\sqrt{j/\log\log j}$ and the higher scale $\sqrt j$.
\begin{itemize}
  \item The \emph{high} part $|Y_j|>\sqrt j$ has \emph{summable} exceedance
        probabilities ($\sum_j\Prob(Y^2>j)\le\E[Y^2]+1<\infty$), so it vanishes
        eventually and contributes $O(1)$ --- exactly the tail step already formalized
        in Chapter~\ref{chap:pre_lil} (Lemma~\ref{lem:trunc_tail_summable}).
  \item The \emph{low} part $|Y_j|\le b_j$ satisfies \textbf{(H)} and gets the sharp LIL
        from Theorem~\ref{thm:llil_sharp}.
  \item The \emph{medium} band $b_j<|Y_j|\le\sqrt j$ is the crux, and is handled by
        Kolmogorov's convergence theorem and Kronecker's lemma: the key point
        (Lemma~\ref{lem:hw_medium_var}) is that this band is only a
        \emph{$\log\log$-factor wide} multiplicatively, so its variance series against
        the LIL scale \emph{converges under the bare second moment}, with no extra
        moment. This is the classical Kolmogorov method; see
        \cite{Stout1974,heyde1980}.
\end{itemize}
A one-shot alternative for the whole upper bound is de~Acosta's proof
\cite{deAcosta1983}; see the remark after Theorem~\ref{thm:hw}.

Let $(Y_j)_{j\ge1}$ be i.i.d., $\E[Y_1]=0$, $\E[Y_1^2]=\sigma^2<\infty$, and
$S_m=\sum_{j\le m}Y_j$. With its natural filtration $S$ is a martingale with
$\langle S\rangle_m=\sigma^2 m$. Write $L(t):=\log\log t$.

\begin{definition}[Three-level truncated decomposition]\label{def:hw_trunc}
Fix $A_0\in(0,\tfrac12)$, with low cutoff $b_j:=A_0\sigma\sqrt{j/L(j)}$ and high cutoff
$\sqrt j$. Split each increment into low, medium and high parts,
\[
  Y_j^{\mathrm L}:=Y_j\indic_{\{|Y_j|\le b_j\}},\quad
  Y_j^{\mathrm M}:=Y_j\indic_{\{b_j<|Y_j|\le\sqrt j\}},\quad
  Y_j^{\mathrm H}:=Y_j\indic_{\{|Y_j|>\sqrt j\}},
\]
so $Y_j=Y_j^{\mathrm L}+Y_j^{\mathrm M}+Y_j^{\mathrm H}$, and set
\[
  \tilde S_m:=\sum_{j\le m}\!\big(Y_j^{\mathrm L}-\E Y_j^{\mathrm L}\big),\quad
  W_m:=\sum_{j\le m}\!\big(Y_j^{\mathrm M}-\E Y_j^{\mathrm M}\big),\quad
  H_m:=\sum_{j\le m}Y_j^{\mathrm H},\quad
  D_m:=\sum_{j\le m}\E\big[Y_j\indic_{\{|Y_j|\le\sqrt j\}}\big].
\]
Since $\E Y_j^{\mathrm L}+\E Y_j^{\mathrm M}=\E[Y_j\indic_{\{|Y_j|\le\sqrt j\}}]$, one has
$S_m=\tilde S_m+W_m+H_m+D_m$.
\end{definition}

\begin{lemma}[High part is eventually zero]\label{lem:hw_high}
\uses{def:hw_trunc, lem:trunc_tail_summable}
Almost surely $Y_j^{\mathrm H}=0$ for all large $j$, so $H_m$ is eventually constant and
$H_m=O(1)$ a.s.
\end{lemma}
\begin{proof}
$\sum_j\Prob(|Y_j|>\sqrt j)=\sum_j\Prob(Y_1^2>j)\le\E[Y_1^2]+1<\infty$ by the layer-cake
bound (Lemma~\ref{lem:trunc_tail_summable}). The first Borel--Cantelli lemma places a.e.\
$\omega$ in only finitely many $\{|Y_j|>\sqrt j\}$; past that index $Y_j^{\mathrm H}=0$.
\end{proof}

\begin{lemma}[Low part satisfies \textbf{(H)}]\label{lem:hw_low_H}
\uses{def:hw_trunc, thm:llil_general, thm:llil_sharp}
$\tilde S$ is a martingale with $|\Delta\tilde S_j|\le2b_j$ and
$\langle\tilde S\rangle_m=\sum_{j\le m}\Var(Y_j^{\mathrm L})$. Since
$\Var(Y_j^{\mathrm L})\to\sigma^2$ (as $b_j\to\infty$), $\langle\tilde S\rangle_m\sim
\sigma^2 m$; in particular $\tfrac12\sigma^2 m\le\langle\tilde S\rangle_m\le\sigma^2 m$
eventually. Hence $\tilde S$ satisfies \textbf{(H)} with $A=2A_0<1$, so
Theorem~\ref{thm:llil_general} gives $\tilde S_m=O(\sqrt{m\,L(m)})$ a.s., and the sharp
Theorem~\ref{thm:llil_sharp} gives $\limsup_m\tilde S_m/\sqrt{2\sigma^2 m\,L(m)}\le1$.
\end{lemma}
\begin{proof}
$|\Delta\tilde S_j|=|Y_j^{\mathrm L}-\E Y_j^{\mathrm L}|\le2b_j=2A_0\sigma\sqrt{j/L(j)}
\le2A_0\sqrt{\langle\tilde S\rangle_j/L(\langle\tilde S\rangle_j)}$ for large $j$, using
$\langle\tilde S\rangle_j\ge\tfrac12\sigma^2 j$ and monotonicity of $t\mapsto t/L(t)$;
this is \textbf{(H)} with $A=2A_0$. The variance facts $\Var(Y_j^{\mathrm L})\le\sigma^2$
and $\Var(Y_j^{\mathrm L})=\E[Y^2\indic_{\{|Y|\le b_j\}}]-(\E Y_j^{\mathrm L})^2\to\sigma^2$
(dominated convergence, $\E Y_j^{\mathrm L}\to0$) are the same statements formalized at
level $\sqrt i$ as \texttt{AlphaRAR.truncVar\_le\_armVar} and the variance convergence
$v_i\to\sigma^2$ (Lemma~\ref{lem:trunc_var_conv}). Then $\langle\tilde S\rangle_m/(\sigma^2
m)\to1$, so $L(\langle\tilde S\rangle_m)\sim L(m)$, fixing the constant in the sharp bound.
\end{proof}

\begin{lemma}[Medium variance series converges under a bare second moment]\label{lem:hw_medium_var}
\uses{def:hw_trunc}
$\displaystyle\sum_{j\ge3}\frac{\Var(Y_j^{\mathrm M})}{j\,L(j)}<\infty$.
\end{lemma}
\begin{proof}
$\Var(Y_j^{\mathrm M})\le\E[(Y_j^{\mathrm M})^2]=\E[Y^2\indic_{\{b_j<|Y|\le\sqrt j\}}]$,
so by Tonelli
\[
  \sum_{j\ge3}\frac{\Var(Y_j^{\mathrm M})}{j\,L(j)}\ \le\
  \E\Big[Y^2\!\!\sum_{j:\,b_j<|Y|\le\sqrt j}\frac{1}{j\,L(j)}\Big].
\]
Fix $y=|Y|$. The condition $b_j<y\le\sqrt j$ means $y^2\le j$ and
$j/L(j)<(y/(A_0\sigma))^2$; as $t\mapsto t/L(t)$ is increasing the latter is $j<j_0(y)$,
where $j_0(y)/L(j_0(y))=(y/(A_0\sigma))^2$, so $j_0(y)\asymp y^2 L(y^2)$. Thus $j$ runs
over $[y^2,j_0(y))$, an interval of multiplicative width $\asymp L(y^2)$ on which
$L(j)\asymp L(y^2)$; hence
\[
  \sum_{y^2\le j<j_0(y)}\frac{1}{j\,L(j)}\ \asymp\
  \frac{1}{L(y^2)}\log\frac{j_0(y)}{y^2}\ \asymp\ \frac{\log L(y^2)}{L(y^2)}\ =\ O(1),
\]
bounded uniformly in $y$ (indeed $\to0$). The inner factor being bounded, the double sum
is $\le C\,\E[Y^2]<\infty$. The mechanism: the medium band is only a $\log\log$-factor
wide, so the harmonic inner sum is $O(\log\log\log/\log\log)=O(1)$; this is exactly where
the bare second moment suffices --- a $\sqrt j$ high cutoff makes the band any wider and
the sum would need $\E[Y^2\log\log|Y|]$.
\end{proof}

\begin{lemma}[Medium part is negligible]\label{lem:hw_medium}
\uses{lem:hw_medium_var, lem:slln_l2_conv, lem:kronecker_general}
$W_m=o(\sqrt{m\,L(m)})$ a.s.
\end{lemma}
\begin{proof}
Put $a_j:=\sqrt{2\,j\,L(j)}$, positive nondecreasing with $a_j\to\infty$. The partial
sums of the independent, mean-zero terms $(Y_j^{\mathrm M}-\E Y_j^{\mathrm M})/a_{j+1}$
form an $L^2$-bounded martingale, since by independence and Lemma~\ref{lem:hw_medium_var}
\[
  \int\Big(\sum_{j<n}\tfrac{Y_j^{\mathrm M}-\E Y_j^{\mathrm M}}{a_{j+1}}\Big)^2
  =\sum_{j<n}\frac{\Var(Y_j^{\mathrm M})}{a_{j+1}^2}
  \le\tfrac12\sum_{j}\frac{\Var(Y_j^{\mathrm M})}{j\,L(j)}<\infty .
\]
Hence $\sum_j(Y_j^{\mathrm M}-\E Y_j^{\mathrm M})/a_{j+1}$ converges a.s.\ by the $L^2$
martingale convergence theorem (Lemma~\ref{lem:slln_l2_conv}), and Kronecker's lemma with
weights $b_j=a_j$ (Lemma~\ref{lem:kronecker_general}, whose conclusion
$a_m^{-1}\sum_{j<m}a_{j+1}y_j\to0$ with $y_j=(Y_j^{\mathrm M}-\E Y_j^{\mathrm M})/a_{j+1}$
is exactly $a_m^{-1}\sum_{j<m}(Y_j^{\mathrm M}-\E Y_j^{\mathrm M})\to0$) gives
$W_m=o(a_m)=o(\sqrt{m\,L(m)})$ a.s. This is the Kolmogorov-convergence-plus-Kronecker
pattern of the martingale SLLN (Chapter~\ref{chap:pre_slln}); both ingredients are
already formalized.
\end{proof}

\begin{lemma}[Drift is negligible]\label{lem:hw_drift}
\uses{def:hw_trunc}
$|D_m|\le2\sigma^2\sqrt m$, so $D_m=O(\sqrt m)=o(\sqrt{m\,L(m)})$.
\end{lemma}
\begin{proof}
Each summand $\E[Y_j\indic_{\{|Y_j|\le\sqrt j\}}]$ is the mean of the truncation of the
centred $Y_j$ at level $\sqrt j$, so the already-formalized centred truncated-mean bound
$|\!\int\operatorname{trunc}(X,A)|\le\E[X^2]/A$
(\texttt{AlphaRAR.abs\_integral\_truncation\_le}, about Mathlib's
\texttt{MeasureTheory.truncation}) gives $|\E[Y_j\indic_{\{|Y_j|\le\sqrt j\}}]|\le
\sigma^2/\sqrt j$. Summing with $\sum_{j\le m}1/\sqrt j\le2\sqrt m$
(\texttt{AlphaRAR.sum\_one\_div\_sqrt\_le}, also formalized) gives $|D_m|\le2\sigma^2
\sqrt m$; and $\sqrt m/\sqrt{m\,L(m)}\to0$, so $D_m=o(\sqrt{m\,L(m)})$. Both cited bounds
are exactly those used for the drift in the existing (log-rate) truncation route
(Lemmas~\ref{lem:trunc_mean_bound}--\ref{lem:trunc_drift}), applied here at level
$\sqrt j$ instead of $\sqrt i$.
\end{proof}

\begin{theorem}[i.i.d.\ Hartman--Wintner, upper half]\label{thm:hw}
\uses{lem:hw_high, lem:hw_low_H, lem:hw_medium, lem:hw_drift}
For i.i.d.\ $(Y_j)$ with $\E[Y_1]=0$ and $\E[Y_1^2]=\sigma^2<\infty$, almost surely
\[
  \limsup_{m\to\infty}\frac{S_m}{\sqrt{2\sigma^2 m\,\log\log m}}\ \le\ 1,
\]
in particular $S_m=O(\sqrt{m\log\log m})$; the two-sided form follows from $-Y$.
\end{theorem}
\begin{proof}
By Definition~\ref{def:hw_trunc}, $S_m=\tilde S_m+W_m+H_m+D_m$, and the three remainders
are $o(\sqrt{m\,L(m)})$: $H_m=O(1)$ (Lemma~\ref{lem:hw_high}), $W_m=o(\sqrt{m\,L(m)})$
(Lemma~\ref{lem:hw_medium}) and $D_m=o(\sqrt m)$ (Lemma~\ref{lem:hw_drift}). Hence
$\limsup_m S_m/\sqrt{2\sigma^2 m\,L(m)}=\limsup_m\tilde S_m/\sqrt{2\sigma^2 m\,L(m)}\le1$
by the sharp low-part bound (Lemma~\ref{lem:hw_low_H}). No $A_0\to0$ limit is needed: the
low truncation $b_j\to\infty$ already gives $\Var(Y_j^{\mathrm L})\to\sigma^2$, so the
constant is sharp. The two-sided statement applies this to $-Y$.
\end{proof}

\begin{remark}[de~Acosta's one-shot alternative]\label{rem:hw_deacosta}
The bound $\limsup_m S_m/\sqrt{2\sigma^2 m\log\log m}\le1$ can instead be obtained in one
stroke by de~Acosta's method \cite{deAcosta1983}, which avoids the three-level split and
may be cleaner to formalize as a self-contained i.i.d.\ result. In that route
Theorem~\ref{thm:hw} is proved directly, while the general Theorem~\ref{thm:llil_general}
is still used for the genuinely martingale (bounded or \textbf{(H)}) applications ---
notably the assignment martingale $M_{\cdot,k}$.
\end{remark}

\section{Optional skipping and the response martingale}
\label{sec:llil_skip}

The finite-variance martingales of Part~I are of the form
$\sum_{j<n}D_j(Y_j-\theta)$ with $(Y_j)$ i.i.d.\ independent of the past and $(D_j)$
predictable $\{0,1\}$-valued. Doob's optional-skipping theorem turns such a martingale
into a genuine i.i.d.\ partial sum, so Theorem~\ref{thm:hw} applies to it directly ---
no separate martingale finite-variance LIL is needed.

\begin{lemma}[Optional skipping]\label{lem:opt_skip}
Let $(Y_j)_{j\ge1}$ be i.i.d.\ with law $\rho$, adapted to $\filt$ with $Y_j$
independent of $\filt_{j-1}$, and $(D_j)$ predictable ($D_j$ is $\filt_{j-1}$-measurable)
$\{0,1\}$-valued with $\sum_j D_j=\infty$ a.s. Let $\tau_1<\tau_2<\cdots$ enumerate
$\{j:D_j=1\}$. Then $(Y_{\tau_m})_{m\ge1}$ is i.i.d.\ with law $\rho$; consequently, for
every $\theta$,
\[
  \sum_{j<n}D_j\,(Y_j-\theta)\ =\ S_{N_n},\qquad
  S_m:=\sum_{i\le m}(Y_{\tau_i}-\theta),\quad N_n:=\sum_{j<n}D_j.
\]
\end{lemma}
\begin{proof}
The classical optional-skipping / optional-sampling argument for i.i.d.\ sequences: for
any bounded measurable $g_1,\dots,g_p$, conditioning successively on
$\filt_{\tau_i-1}$ and using $Y_{\tau_i}\perp\filt_{\tau_i-1}$ with law $\rho$ (as
$\{\tau_i=j\}\in\filt_{j-1}$) gives $\E\prod_i g_i(Y_{\tau_i})=\prod_i\int g_i\,d\rho$,
which is the i.i.d.\ law. The displayed identity is just a reindexing of the nonzero
summands. Mathlib building blocks: the independence/`Kernel` API and `Filtration`
machinery; this general statement is not yet in Mathlib and is a worthwhile
contribution.
\end{proof}

\begin{corollary}[LIL for a predictably-sampled i.i.d.\ sum]\label{cor:subsampled_lil}
\uses{lem:opt_skip, thm:hw}
With the notation of Lemma~\ref{lem:opt_skip}, if $\E[Y_1]=\theta$ and
$\Var(Y_1)=V<\infty$, then almost surely
\[
  \Big|\sum_{j<n}D_j(Y_j-\theta)\Big|\ =\ O\!\big(\sqrt{N_n\,\log\log N_n}\big).
\]
\end{corollary}
\begin{proof}
By Lemma~\ref{lem:opt_skip} the left side is $|S_{N_n}|$ for the i.i.d.\ centred sum
$S$; since $N_n\to\infty$ a.s., Theorem~\ref{thm:hw} (two-sided) gives
$|S_{N_n}|=O(\sqrt{N_n\log\log N_n})$.
\end{proof}

\textbf{Application.} This is the shape used in Lemma~\ref{lem:theta_LIL}: with
$Y_j=\xi_{j,k}$, $\theta=\theta_k$, $D_j=X_{j,k}$, and $N_n=N_{n,k}$, it yields
$Q_{n,k}=O(\sqrt{N_{n,k}\log\log N_{n,k}})$ a.s., the $\log\log$ upgrade of the
$Q$-bound. The bounded assignment martingale $M_{\cdot,k}$ instead goes through
Corollary~\ref{cor:llil_nat} (Theorem~\ref{thm:llil_bounded}), as it is genuinely
adaptive rather than a subsampled i.i.d.\ sum.

\medskip
\noindent\textbf{Relationship to Chapter~\ref{chap:pre_lil}.} The bounded-increment
$O$-corollary (Corollary~\ref{cor:llil_nat}) supersedes Corollaries~\ref{cor:lil_upper}
and \ref{cor:lil_upper_nat} with the $\log\log$ rate; Theorem~\ref{thm:hw} +
Corollary~\ref{cor:subsampled_lil} supersede the growing-increment truncation route
(the $\log$-rate \texttt{measure\_exists\_ge\_le\_exp\_block} family), which remains a
valid $\log$-rate fallback. Nothing in this chapter is formalized yet; it reuses the
formalized supermartingale, Ville, Freedman, horizon and Borel--Cantelli lemmas of
Chapter~\ref{chap:pre_lil}, and, for the i.i.d.\ layer, the already-formalized tools of
the SLLN chapter: the layer-cake summability
Lemma~\ref{lem:trunc_tail_summable} (high part), the $L^2$ martingale convergence
theorem Lemma~\ref{lem:slln_l2_conv} and Kronecker's lemma
Lemma~\ref{lem:kronecker_general} (medium part). The genuinely new devices are just the
stopping-time quadratic-variation stop (Lemma~\ref{lem:llil_qv_stop}), the refined
one-step estimate (Lemma~\ref{lem:llil_refined_ineq}), the medium-band variance estimate
that carries the bare second moment (Lemma~\ref{lem:hw_medium_var}, the analytic crux),
and optional skipping (Lemma~\ref{lem:opt_skip}).

\chapter{Central limit theorem for martingales}
\label{chap:pre_clt}

We formalize the martingale central limit theorem in Lindeberg form, in dimension
one. This is the analytic core of Lemma~\ref{lem:componentwise}, and through it of
the asymptotic normality Theorem~\ref{thm:normality} and the test calibration
Lemma~\ref{lem:pearson}. A one-dimensional statement suffices: the coordinates of
the estimator error are driven by martingales with vanishing cross-variation
(Lemma~\ref{lem:qv_orthogonal}), so the vector CLT follows coordinatewise via the
Cramér--Wold device or directly from the orthogonality.

\textbf{Formalization note.} Mathlib has the i.i.d.\ one-dimensional CLT
(\texttt{Probability/CentralLimitTheorem.lean}), proved by convergence of
characteristic functions, and the supporting API: \texttt{charFun}
(\texttt{MeasureTheory/Measure/CharacteristicFunction/Basic.lean}), the fact that
\texttt{charFun} determines a measure, and the Gaussian characteristic function
(\texttt{Distributions/Gaussian/CharFun.lean}). What is missing is the martingale
version. The proof reuses the same final step (charFun convergence $\Rightarrow$
convergence in distribution) but replaces independence by a
conditional-characteristic-function argument.

\section{Analytic ingredient}
\label{sec:pre_clt_analytic}

\begin{lemma}[Complex exponential expansion]\label{lem:clt_exp_bound}
For all $x \in \R$,
\[
  \Big| e^{ix} - 1 - ix + \tfrac{x^2}{2} \Big| \le \min\!\Big(x^2,\ \tfrac{|x|^3}{6}\Big).
\]
\end{lemma}

\begin{proof}
Standard: integrate the remainder of the Taylor expansion of $e^{ix}$; the two
bounds come from estimating the remainder integral to first and to second order.
\end{proof}

\begin{lemma}[Conditional characteristic increment]\label{lem:clt_cond_char}
\uses{lem:clt_exp_bound, def:pred_qv}
Let $(M_{n,i})_{i\le k_n}$ be a martingale array with increments $\Delta_{n,i}$.
For fixed $t$,
\[
  \E\big[e^{it\Delta_{n,i}}\mid\filt_{i-1}\big]
  = 1 - \tfrac{t^2}{2}\,\E[\Delta_{n,i}^2\mid\filt_{i-1}] + r_{n,i},
\]
where $\sum_i |r_{n,i}|$ is controlled by the conditional Lindeberg quantity
$\sum_i \E[\Delta_{n,i}^2\indic_{\{|\Delta_{n,i}|>\varepsilon\}}\mid\filt_{i-1}]$
plus $\varepsilon\, t\, \langle M_n\rangle_{k_n}$.
\end{lemma}

\begin{proof}
Take conditional expectation in Lemma~\ref{lem:clt_exp_bound}: the linear term
vanishes since $\E[\Delta_{n,i}\mid\filt_{i-1}]=0$, and the remainder is bounded
by $\E[\min(\Delta_{n,i}^2,|\Delta_{n,i}|^3/6)\mid\filt_{i-1}]$, split at the
threshold $\varepsilon$ into a Lindeberg part and a part $\le \varepsilon t
\E[\Delta_{n,i}^2\mid\filt_{i-1}]$.
\end{proof}

\section{Convergence of characteristic functions}
\label{sec:pre_clt_char}

\begin{lemma}[Telescoping product]\label{lem:clt_product}
\uses{lem:clt_cond_char}
With the notation above, if $\langle M_n\rangle_{k_n} \xrightarrow{\Prob}\sigma^2$
and the conditional Lindeberg condition holds, then
\[
  \E\big[e^{it M_{n,k_n}}\big] \longrightarrow \exp\!\Big(-\tfrac{\sigma^2 t^2}{2}\Big)
  \qquad\text{for every } t\in\R.
\]
\end{lemma}

\begin{proof}
\uses{lem:clt_cond_char, lem:qv_incr}
Write $e^{itM_{n,k_n}} = \prod_i e^{it\Delta_{n,i}}$ and compare with the
predictable product $\prod_i(1 - \tfrac{t^2}{2}\E[\Delta_{n,i}^2\mid\filt_{i-1}])
\approx \exp(-\tfrac{t^2}{2}\langle M_n\rangle_{k_n})$. The telescoping bound uses
$|\prod a_i - \prod b_i| \le \sum_i |a_i - b_i|$ for terms of modulus $\le1$,
with $a_i - b_i = r_{n,i}$ from Lemma~\ref{lem:clt_cond_char}, whose sum vanishes
by Lindeberg. Then $\langle M_n\rangle_{k_n}\xrightarrow{\Prob}\sigma^2$ gives the
limit $\exp(-\sigma^2 t^2/2)$.
\end{proof}

\begin{lemma}[Characteristic functions determine the limit]\label{lem:levy}
The map sending a probability measure to its characteristic function is injective,
and if $\E[e^{itX_n}]\to\exp(-\sigma^2 t^2/2)$ for all $t$, then $X_n$ converges in
distribution to $\mathcal{N}(0,\sigma^2)$.
\end{lemma}

\begin{proof}
Injectivity is \texttt{MeasureTheory.Measure.ext\_of\_charFun} in Mathlib. The
target $\exp(-\sigma^2 t^2/2)$ is the characteristic function of
$\mathcal{N}(0,\sigma^2)$ (\texttt{Distributions/Gaussian/CharFun.lean}). The
implication ``pointwise charFun convergence to a continuous limit charFun
$\Rightarrow$ convergence in distribution'' is Lévy's continuity theorem; Mathlib
has the special case used in its i.i.d.\ CLT
(\texttt{Probability/CentralLimitTheorem.lean}), which is exactly what is needed
here since the limit is a fixed Gaussian.
\end{proof}

\textbf{Formalization note.} A fully general Lévy continuity theorem (arbitrary
tight limits) is not yet in Mathlib, but is not required: we only need
convergence to a fixed, known Gaussian limit, which the CLT file already handles.

\section{The theorem}
\label{sec:pre_clt_thm}

\begin{theorem}[Martingale CLT, Lindeberg form]\label{thm:mart_clt}
\uses{lem:clt_product, lem:levy}
Let $(M_{n,i})_{i \le k_n}$ be a martingale array (with respect to filtrations
$(\filt_{n,i})_i$) with $M_{n,0}=0$, predictable quadratic variation
$\langle M_n\rangle_{k_n}\xrightarrow{\Prob}\sigma^2$, and satisfying the
conditional Lindeberg condition: for every $\varepsilon>0$,
\[
  \sum_{i}\E\big[\Delta_{n,i}^2\indic_{\{|\Delta_{n,i}|>\varepsilon\}}\mid\filt_{n,i-1}\big]
  \xrightarrow{\Prob} 0.
\]
Then $M_{n,k_n}\eqdist\mathcal{N}(0,\sigma^2)$.
\end{theorem}

\begin{proof}
Combine Lemma~\ref{lem:clt_product} (convergence of characteristic functions) with
Lemma~\ref{lem:levy} (the limit is $\mathcal{N}(0,\sigma^2)$).
\end{proof}

\begin{corollary}[Self-normalized componentwise form]\label{cor:mart_clt_componentwise}
\uses{thm:mart_clt, lem:qv_orthogonal}
Let $M_{n,k}=\sum_{i\le n}X_{i,k}(\xi_{i,k}-\theta_k)$ as in
Definition~\ref{def:M}--\ref{def:Q}, with $N_{n,k}\to\infty$ a.s.\ for each $k$.
Then, with $D_n=\mathrm{diag}(\sqrt{N_{n,1}},\dots,\sqrt{N_{n,K}})$,
$D_n(\estParams_n-\Params)\eqdist\mathcal{N}(0,\mathrm{diag}(V_1,\dots,V_K))$.
\end{corollary}

\begin{proof}
\uses{thm:mart_clt, lem:qv_orthogonal}
Apply Theorem~\ref{thm:mart_clt} to the time-changed martingale
$i\mapsto Q_{\tau_k(i),k}$ stopped when $N_{\cdot,k}=i$; its quadratic variation
normalized by $N_{n,k}$ tends to $V_k$, and the conditional Lindeberg condition
holds because the $\xi_{i,k}$ are i.i.d.\ with finite variance (dominated
convergence, cf.\ the proof sketch of Lemma~\ref{lem:componentwise}). Distinct
coordinates have vanishing cross-variation by Lemma~\ref{lem:qv_orthogonal}, so
the joint limit is the stated diagonal Gaussian.
\end{proof}

\chapter{Chi-square distribution and Cochran's theorem}
\label{chap:pre_cochran}

To calibrate the Pearson test (Lemma~\ref{lem:pearson}) we need the $\chi^2$
distribution and the special case of Cochran's theorem stating that a Gaussian
quadratic form through an orthogonal projection of rank $r$ is $\chi^2_r$.

\textbf{Formalization note.} Mathlib has the Gamma and exponential distributions
(\texttt{Distributions/Gamma.lean}, \texttt{Distributions/Exponential.lean}) and
a rich Gaussian theory including rotation invariance of the standard Gaussian
(\texttt{Distributions/Gaussian/}, \texttt{IsGaussian}, \texttt{stdGaussian}). It
lacks a named $\chi^2$ distribution, the additivity of Gamma laws, and any
Cochran-type result. All three are self-contained additions given the existing
Gaussian and Gamma APIs.

\section{The chi-square distribution}
\label{sec:pre_chisq_def}

\begin{definition}[Chi-square distribution]\label{def:chiSquared}
For an integer $r\ge1$, the \emph{chi-square distribution with $r$ degrees of
freedom} $\chi^2_r$ is the Gamma distribution with shape $r/2$ and rate $1/2$.
\end{definition}

\begin{lemma}[Square of a standard Gaussian]\label{lem:sq_gaussian_chisq1}
\uses{def:chiSquared}
If $Z\sim\mathcal{N}(0,1)$ then $Z^2\sim\chi^2_1$.
\end{lemma}

\begin{proof}
Compute the law of $Z^2$ by the change of variables $x = z^2$ on the standard
Gaussian density; the resulting density $\frac{1}{\sqrt{2\pi x}}e^{-x/2}$ on
$(0,\infty)$ is the $\mathrm{Gamma}(1/2,1/2)$ density, i.e.\ $\chi^2_1$.
\end{proof}

\begin{lemma}[Additivity of Gamma laws]\label{lem:gamma_add}
If $U\sim\mathrm{Gamma}(a,\lambda)$ and $V\sim\mathrm{Gamma}(b,\lambda)$ are
independent, then $U+V\sim\mathrm{Gamma}(a+b,\lambda)$.
\end{lemma}

\begin{proof}
The characteristic function of $\mathrm{Gamma}(a,\lambda)$ is
$(1 - it/\lambda)^{-a}$; by independence the characteristic function of $U+V$ is
the product $(1-it/\lambda)^{-(a+b)}$, which identifies the law by injectivity of
characteristic functions.
\end{proof}

\textbf{Formalization note.} Mathlib's \texttt{Gamma.lean} provides the density
but, as of writing, not this convolution/additivity identity; it can be proved via
the characteristic-function API used elsewhere in this part, or directly by a Beta
integral convolution of densities.

\begin{lemma}[Chi-square as a sum of squares]\label{lem:sum_sq_chisq}
\uses{def:chiSquared, lem:sq_gaussian_chisq1, lem:gamma_add}
If $Z_1,\dots,Z_r\sim\mathcal{N}(0,1)$ are independent, then
$\sum_{i=1}^r Z_i^2\sim\chi^2_r$.
\end{lemma}

\begin{proof}
Each $Z_i^2\sim\chi^2_1=\mathrm{Gamma}(1/2,1/2)$ by
Lemma~\ref{lem:sq_gaussian_chisq1}, and by independence and
Lemma~\ref{lem:gamma_add} the sum is $\mathrm{Gamma}(r/2,1/2)=\chi^2_r$.
\end{proof}

\section{Cochran's theorem for a projection}
\label{sec:pre_cochran_thm}

\begin{lemma}[Rotation invariance of the standard Gaussian]\label{lem:gaussian_rotation}
If $Z\sim\mathcal{N}(0,I_K)$ and $O$ is an orthogonal $K\times K$ matrix, then
$OZ\sim\mathcal{N}(0,I_K)$.
\end{lemma}

\begin{proof}
$OZ$ is Gaussian with mean $0$ and covariance $O I_K O^\top = I_K$. In Mathlib
this is an instance of the invariance of \texttt{stdGaussian}/\texttt{IsGaussian}
under orthogonal (more generally, measure-preserving linear) maps.
\end{proof}

\begin{lemma}[Diagonalization of an orthogonal projection]\label{lem:proj_diag}
An orthogonal projection matrix $P$ of rank $r$ can be written
$P = O^\top \mathrm{diag}(I_r, 0)\, O$ for some orthogonal matrix $O$.
\end{lemma}

\begin{proof}
$P$ is symmetric and idempotent, hence diagonalizable in an orthonormal
eigenbasis with eigenvalues in $\{0,1\}$, exactly $r$ of which are $1$. This is
the spectral theorem for symmetric matrices, available in Mathlib.
\end{proof}

\begin{theorem}[Gaussian quadratic form under a projection]\label{thm:cochran}
\uses{lem:sum_sq_chisq, lem:gaussian_rotation, lem:proj_diag}
Let $Z\sim\mathcal{N}(0,I_K)$ and let $P$ be an orthogonal projection of rank $r$.
Then $Z^\top P Z\sim\chi^2_r$.
\end{theorem}

\begin{proof}
Write $P = O^\top\mathrm{diag}(I_r,0)O$ (Lemma~\ref{lem:proj_diag}) and set
$W = OZ$, which is $\mathcal{N}(0,I_K)$ by Lemma~\ref{lem:gaussian_rotation}.
Then $Z^\top P Z = W^\top\mathrm{diag}(I_r,0)W = \sum_{i=1}^r W_i^2$, which is
$\chi^2_r$ by Lemma~\ref{lem:sum_sq_chisq}.
\end{proof}

\textbf{Application.} In Lemma~\ref{lem:pearson}, $P = I - ss^\top$ is an
orthogonal projection of rank $K-1$ (Lemma~\ref{lem:projection}), and the limiting
statistic is $Z^\top P Z$ with $Z\sim\mathcal{N}(0,I_K)$; Theorem~\ref{thm:cochran}
then gives the $\chi^2_{K-1}$ limit.

